% Leçon 8 — Éléments socioculturels : processus d'hérédité au Cameroun et en Chine (suite : comparaison et exposé) — Chinois Tle A — Cours signé OnBuch+
% Programme officiel MINESEC. Compiler avec : tectonic ZH08-socio-processus-dheredite-au-cameroun-et-en-chine.tex
\newcommand{\DOCMATIERE}{Chinois}
\newcommand{\DOCNIVEAU}{Tle A}
\newcommand{\DOCMODULE}{Module 1 : Vie familiale et intégration sociale}
\newcommand{\DOCLECON}{ZH08}
\newcommand{\DOCTITRE}{Leçon 8 — Éléments socioculturels : processus d'hérédité au Cameroun et en Chine (suite : comparaison et exposé)}
% OnBuch+ — préambule « cours signés » ENRICHI (compilé avec tectonic / XeLaTeX)
% Système visuel « Pop » d'OnBuch+ : fond crème (jamais blanc), bordures encre
% épaisses, ombres « dures » décalées, titres Archivo Black en capitales.
% Version MATHÉMATIQUES : graphes (pgfplots/tikz), boîtes pédagogiques
% (définition, propriété, méthode, attention, le savais-tu, exercice résolu,
% exercice, TP, bilan), page de garde et signature OnBuch+ ancrée en bas.
%
% Macros attendues dans main.tex AVANT \input{preamble} :
%   \DOCMATIERE  \DOCNIVEAU  \DOCMODULE  \DOCLECON (numéro)  \DOCTITRE
\documentclass[11pt]{article}

\usepackage{fontspec}
\usepackage[french]{babel} % français = langue principale ; le chinois via \zh{..} / textebox (xeCJK)
\usepackage{xeCJK}
\usepackage{pifont} % \ding{51} \ding{55} utilisés par les modèles
\usepackage[a4paper,margin=2cm,top=3.4cm,bottom=2.8cm,headheight=1.4cm,footskip=1.3cm]{geometry}
\usepackage{amsmath,amssymb,amsfonts}
\usepackage{mathtools} % \xmapsto, \coloneqq... souvent utilisés par les modèles
\usepackage{cancel} % simplifications barrées (bilans de forces)
% \abs{z} (module, valeur absolue) et \norm{u} (norme), utilisés par les modèles
\providecommand{\abs}{}\let\abs\relax\DeclarePairedDelimiter{\abs}{\lvert}{\rvert}
\providecommand{\norm}{}\let\norm\relax\DeclarePairedDelimiter{\norm}{\lVert}{\rVert}
\usepackage[table]{xcolor} % [table] : fournit aussi \rowcolors (alternance de lignes)
\usepackage{pagecolor}
\usepackage{tikz}
\usetikzlibrary{arrows.meta,positioning,calc,shapes.geometric,shapes.misc,shapes.arrows,decorations.pathmorphing,decorations.pathreplacing,decorations.markings,patterns,shadows,mindmap,trees,fit,backgrounds,matrix,intersections,angles,quotes}
\usepackage{pgfplots}
\usepackage[version=4]{mhchem} % équations nucléaires \ce{^{235}_{92}U}
\usepackage[european,siunitx]{circuitikz} % schémas électriques
\pgfplotsset{compat=1.18}
\usepgfplotslibrary{fillbetween,groupplots} % aires entre courbes (calcul intégral)
\usepackage{enumitem}
\usepackage{fancyhdr}
% [most] charge toutes les bibliothèques tcolorbox (listings, minted, raster,
% documentation...) et gonfle inutilement la mémoire TeX sur un document long
% (~300 boîtes) : on ne charge que ce qui est réellement utilisé.
\usepackage[skins,breakable]{tcolorbox}
\usepackage{graphicx}
\usepackage{colortbl}
\usepackage{booktabs}
\usepackage{tabularx}
\usepackage{diagbox} % case d'en-tête diagonale des tableaux à double entrée
% Colonnes extensibles centrée / gauche / droite (C, L, R), souvent utilisées par les modèles
\newcolumntype{C}{>{\centering\arraybackslash}X}
\newcolumntype{L}{>{\raggedright\arraybackslash}X}
\newcolumntype{R}{>{\raggedleft\arraybackslash}X}
\usepackage{array}
\usepackage{multicol}
\usepackage{multirow}
\usepackage{siunitx}
% parse-numbers=false : accepte aussi \SI{2,5 \times 10^{-3}}{...}
\sisetup{locale=FR,output-decimal-marker={,},group-minimum-digits=4,parse-numbers=false}
\DeclareSIUnit\ohm{\ensuremath{\symup{\Omega}}} % U+2126 absent de la police
\DeclareSIUnit\division{div} % réglages d'oscilloscope (ms/div, V/div)
\DeclareSIUnit\tr{tr} % tours (vitesse de rotation, ex. \SI{1500}{\tr\per\minute})
\DeclareSIUnit\molar{M} % concentration molaire (ex. \SI{3}{\molar}, \SI{1}{\micro\molar})
\DeclareSIUnit\an{an} % année (le modèle confond parfois avec \year, un compteur TeX) — ex. \SI{5}{\kilo\meter\per\an}
\DeclareSIUnit\FCFA{FCFA} % franc CFA (ex. \SI{500}{\FCFA\per\kilo\gram})
\DeclareSIUnit\degreeBrix{\textdegree Brix} % teneur en sucre d'un sirop/jus (réfractométrie)
\DeclareSIUnit\month{mois} % le modèle utilise parfois \month, un compteur TeX, comme unité de durée
\usepackage{qrcode}
% \degree : commande fréquemment utilisée par les modèles pour un angle ou une
% température en dehors de \SI{}{} (non fournie par les paquets chargés).
\providecommand{\degree}{^{\circ}}
% \qed (fin de démonstration), utilisé par les modèles sans amsthm
\providecommand{\qed}{\hfill$\square$}
% \barre : double barre des valeurs interdites dans un tableau de variations (tkz-tab)
\providecommand{\barre}{\ensuremath{\|}}
% environnement proof (amsthm non chargé) utilisé par les modèles
\ifdefined\proof\else\newenvironment{proof}[1][Démonstration]{\par\noindent\textit{#1.}\ }{\qed\par}\fi
\usepackage{lastpage}
\usepackage{hyperref}
\hypersetup{hidelinks}

% ============================================================
% Palette « Pop » OnBuch+ — mêmes hex que la classe P (plus_ui.dart)
% ============================================================
\definecolor{poporange}{HTML}{F59321}
\definecolor{popdark}{HTML}{E07A0C}
\definecolor{popcream}{HTML}{FFF6E8}
\definecolor{popink}{HTML}{1C1714}
\definecolor{poppurple}{HTML}{7A5AE0}
\definecolor{popgreen}{HTML}{2E9E6A}
\definecolor{poppink}{HTML}{E0457B}
\definecolor{popblue}{HTML}{2F7DE1}
\definecolor{popgold}{HTML}{C98A12}
\definecolor{popmuted}{HTML}{8A7F76}
\definecolor{popline}{HTML}{EADFD2}
\definecolor{popgray}{HTML}{8A7F76}
\colorlet{popgrey}{popgray}
% Alias tolérants (le modèle nomme parfois les couleurs différemment)
\colorlet{popwhite}{white}
\colorlet{popblack}{popink}
\colorlet{popred}{poppink}
\colorlet{popyellow}{popgold}
% Teintes claires (fonds de tableaux/encadrés) — le modèle nomme parfois
% les couleurs avec un suffixe L (ex. popblueL) sans les avoir définies.
\colorlet{poporangeL}{poporange!20}
\colorlet{popdarkL}{popdark!20}
\colorlet{poppurpleL}{poppurple!20}
\colorlet{popgreenL}{popgreen!20}
\colorlet{poppinkL}{poppink!20}
\colorlet{popblueL}{popblue!20}
\colorlet{popgoldL}{popgold!20}
\colorlet{popredL}{popred!20}
\colorlet{popgrayL}{popgray!20}
% Teintes claires pour fonds de tableaux / zones
\colorlet{popblueL}{popblue!10!white}
\colorlet{poporangeL}{poporange!14!white}
\colorlet{popgreenL}{popgreen!12!white}
\colorlet{poppurpleL}{poppurple!12!white}
\colorlet{poppinkL}{poppink!12!white}
\colorlet{popgoldL}{popgold!14!white}

\pagecolor{popcream}

% ============================================================
% Typographie
% ============================================================
\defaultfontfeatures{Ligatures=TeX}
\setmainfont{PlusJakartaSans-Medium.ttf}[Path=../../../fonts/,BoldFont=PlusJakartaSans-Bold.ttf]
% Symboles, API et emojis absents de la police principale : repli sur DejaVu Sans (caractères actifs)
\newfontface\symfont{DejaVuSans.ttf}[Path=../../../fonts/]
\makeatletter
\newcommand\symactive[1]{\catcode#1=13 \begingroup\lccode`\~=#1 \lowercase{\endgroup\def~}{{\symfont\char#1}}}
\newcommand\symblank[1]{\catcode#1=13 \begingroup\lccode`\~=#1 \lowercase{\endgroup\def~}{}}
\newcommand\symhyph[1]{\catcode#1=13 \begingroup\lccode`\~=#1 \lowercase{\endgroup\def~}{\mbox{-}}}
\newcount\symc
\newcommand\symrange[2]{\symc=#1 \@whilenum\symc<#2 \do{\expandafter\symactive\expandafter{\the\symc}\advance\symc by 1 }}
\newcommand\symblankrange[2]{\symc=#1 \@whilenum\symc<#2 \do{\expandafter\symblank\expandafter{\the\symc}\advance\symc by 1 }}
\makeatother
\symrange{"0250}{"02FF}\symactive{"2713}\symactive{"2714}\symactive{"2717}\symactive{"2718}\symactive{"2610}\symactive{"2611}\symactive{"2612}
\symrange{"2600}{"27BF}
\catcode"2705=13 \begingroup\lccode`\~="2705 \lowercase{\endgroup\def~}{{\symfont\char"2713}}\symblank{"26BD}\symblank{"26A1}
\symrange{"2460}{"257F}\symactive{"223C}\symactive{"2248}\symactive{"2260}
\symhyph{"2011}\symhyph{"2E1A}\symblankrange{"1F300}{"1FAFF}\symblank{"FE0F}
% Caractères chinois : WenQuanYi Zen Hei (sous-ensemble GB2312 + pinyin), gras/italique simulés
\setCJKmainfont{WenQuanYiZenHei-subset.ttf}[Path=../../../fonts/,AutoFakeBold=3,AutoFakeSlant=0.2]
\xeCJKsetup{CJKspace=false}
% Pinyin : voyelles à caron (ǎ ǐ ǒ ǔ ǖ ǘ ǚ ǜ) absentes de la police principale -> caractères actifs
\newfontface\pyfont{WenQuanYiZenHei-subset.ttf}[Path=../../../fonts/]
\newcommand\pyactive[1]{\catcode#1=13 \begingroup\lccode`\~=#1 \lowercase{\endgroup\def~}{{\pyfont\char#1}}}
\pyactive{"01CD}\pyactive{"01CE}\pyactive{"01CF}\pyactive{"01D0}\pyactive{"01D1}\pyactive{"01D2}\pyactive{"01D3}\pyactive{"01D4}\pyactive{"01D5}\pyactive{"01D6}\pyactive{"01D7}\pyactive{"01D8}\pyactive{"01D9}\pyactive{"01DA}\pyactive{"01DB}\pyactive{"01DC}
% Maths en sans-serif (Fira Math) avec chiffres et lettres droites de Plus
% Jakarta, pour que les formules se fondent dans le texte.
\usepackage[math-style=ISO,bold-style=ISO]{unicode-math}
\setmathfont{FiraMath-Regular.otf}[Path=../../../fonts/]
\setmathfont{PlusJakartaSans-Medium.ttf}[Path=../../../fonts/,range={up/num,up/{Latin,latin},\mathup}]
\usepackage{icomma} % virgule décimale sans espace en mode math
% Caractères mathématiques Unicode absents des polices de texte (Plus Jakarta,
% Archivo Black) : rendus par leur équivalent mathématique, en texte comme en
% formule — sinon ils s'affichent en carré vide (ex. « Arithmétique dans ℤ »).
\usepackage{newunicodechar}
\newunicodechar{ℝ}{\ensuremath{\mathbb{R}}}
\newunicodechar{ℤ}{\ensuremath{\mathbb{Z}}}
\newunicodechar{ℂ}{\ensuremath{\mathbb{C}}}
\newunicodechar{ℕ}{\ensuremath{\mathbb{N}}}
\newunicodechar{ℚ}{\ensuremath{\mathbb{Q}}}
\newunicodechar{𝔻}{\ensuremath{\mathbb{D}}}
\newunicodechar{α}{\ensuremath{\alpha}}
\newunicodechar{β}{\ensuremath{\beta}}
\newunicodechar{γ}{\ensuremath{\gamma}}
\newunicodechar{δ}{\ensuremath{\delta}}
\newunicodechar{ε}{\ensuremath{\varepsilon}}
\newunicodechar{θ}{\ensuremath{\theta}}
\newunicodechar{λ}{\ensuremath{\lambda}}
\newunicodechar{μ}{\ensuremath{\mu}}
\newunicodechar{σ}{\ensuremath{\sigma}}
\newunicodechar{τ}{\ensuremath{\tau}}
\newunicodechar{φ}{\ensuremath{\varphi}}
\newunicodechar{ψ}{\ensuremath{\psi}}
\newunicodechar{ω}{\ensuremath{\omega}}
\newunicodechar{Σ}{\ensuremath{\Sigma}}
% majuscules produites par \MakeUppercase dans les titres (α-aminés -> Α)
\newunicodechar{Α}{\ensuremath{\alpha}}
\newunicodechar{Β}{\ensuremath{\beta}}
\newunicodechar{Γ}{\ensuremath{\Gamma}}
\newunicodechar{Δ}{\ensuremath{\Delta}}
\newunicodechar{Ε}{\ensuremath{\varepsilon}}
\newunicodechar{Θ}{\ensuremath{\Theta}}
\newunicodechar{Λ}{\ensuremath{\Lambda}}
\newunicodechar{Μ}{\ensuremath{\mu}}
\newunicodechar{Τ}{\ensuremath{\tau}}
\newunicodechar{Φ}{\ensuremath{\Phi}}
\newunicodechar{Ψ}{\ensuremath{\Psi}}
\newunicodechar{Ω}{\ensuremath{\Omega}}
\newunicodechar{↦}{\ensuremath{\mapsto}}
\newunicodechar{∈}{\ensuremath{\in}}
\newunicodechar{∉}{\ensuremath{\notin}}
\newunicodechar{∧}{\ensuremath{\wedge}}
\newunicodechar{⇔}{\ensuremath{\Leftrightarrow}}
\newunicodechar{⟺}{\ensuremath{\Longleftrightarrow}}
\newunicodechar{⇒}{\ensuremath{\Rightarrow}}
\newunicodechar{⟹}{\ensuremath{\Longrightarrow}}
\newunicodechar{∘}{\ensuremath{\circ}}
\newunicodechar{∪}{\ensuremath{\cup}}
\newunicodechar{∩}{\ensuremath{\cap}}
\newunicodechar{≡}{\ensuremath{\equiv}}
\newunicodechar{⊂}{\ensuremath{\subset}}
\newunicodechar{⊥}{\ensuremath{\perp}}
\newunicodechar{∃}{\ensuremath{\exists}}
\newunicodechar{∀}{\ensuremath{\forall}}
\newunicodechar{∛}{\ensuremath{\sqrt[3]{\,}}}
\newunicodechar{∜}{\ensuremath{\sqrt[4]{\,}}}
\newunicodechar{⃗}{\ensuremath{{}^{\to}}}
\newunicodechar{ⁿ}{\ifmmode^{n}\else\textsuperscript{\ensuremath{n}}\fi}
\newunicodechar{ˣ}{\ifmmode^{x}\else\textsuperscript{\ensuremath{x}}\fi}
\newunicodechar{ᵇ}{\ifmmode^{b}\else\textsuperscript{\ensuremath{b}}\fi}
\newunicodechar{ʳ}{\ifmmode^{r}\else\textsuperscript{\ensuremath{r}}\fi}
\newunicodechar{ᵘ}{\ifmmode^{u}\else\textsuperscript{\ensuremath{u}}\fi}
\newunicodechar{ʸ}{\ifmmode^{y}\else\textsuperscript{\ensuremath{y}}\fi}
\newunicodechar{ᵃ}{\ifmmode^{a}\else\textsuperscript{\ensuremath{a}}\fi}
\newunicodechar{ᵉ}{\ifmmode^{e}\else\textsuperscript{\ensuremath{e}}\fi}
\newunicodechar{ᵏ}{\ifmmode^{k}\else\textsuperscript{\ensuremath{k}}\fi}
\newunicodechar{ᵖ}{\ifmmode^{p}\else\textsuperscript{\ensuremath{p}}\fi}
\newunicodechar{ᴺ}{\ifmmode^{N}\else\textsuperscript{\ensuremath{N}}\fi}
\newunicodechar{ᶜ}{\ifmmode^{c}\else\textsuperscript{\ensuremath{c}}\fi}
\newunicodechar{ʰ}{\ifmmode^{h}\else\textsuperscript{\ensuremath{h}}\fi}
\newunicodechar{ᵠ}{\ifmmode^{q}\else\textsuperscript{\ensuremath{q}}\fi}
\newunicodechar{ₙ}{\ifmmode_{n}\else\textsubscript{\ensuremath{n}}\fi}
\newunicodechar{ₐ}{\ifmmode_{a}\else\textsubscript{\ensuremath{a}}\fi}
\newunicodechar{ᵢ}{\ifmmode_{i}\else\textsubscript{\ensuremath{i}}\fi}
\newunicodechar{ᵣ}{\ifmmode_{r}\else\textsubscript{\ensuremath{r}}\fi}
\newunicodechar{ₖ}{\ifmmode_{k}\else\textsubscript{\ensuremath{k}}\fi}
\newunicodechar{ᵦ}{\ifmmode_{\beta}\else\textsubscript{\ensuremath{\beta}}\fi}
\newunicodechar{ₓ}{\ifmmode_{x}\else\textsubscript{\ensuremath{x}}\fi}
\newfontfamily\popdisplay{ArchivoBlack-Regular.ttf}[Path=../../../fonts/]
\newfontfamily\poplogo{SpaceGrotesk-Bold.ttf}[Path=../../../fonts/]
\newfontfamily\popsemi{PlusJakartaSans-SemiBold.ttf}[Path=../../../fonts/]
\newfontfamily\popxbold{PlusJakartaSans-ExtraBold.ttf}[Path=../../../fonts/]
\newfontfamily\popxboldls{PlusJakartaSans-ExtraBold.ttf}[Path=../../../fonts/,LetterSpace=3.0]

\setlist[enumerate]{leftmargin=1.7em,itemsep=5pt,topsep=4pt}
\setlist[itemize]{leftmargin=1.7em,itemsep=4pt,topsep=4pt,label={\color{poporange}\textbullet}}
\setlength{\parindent}{0pt}
\setlength{\parskip}{5pt}
\renewcommand{\arraystretch}{1.25}

% pgfplots : style Pop par défaut
\pgfplotsset{
  every axis/.append style={axis line style={popink,line width=1pt},tick style={popink},
    grid style={popline,line width=0.5pt},label style={font=\small},tick label style={font=\footnotesize}},
  popcurve/.style={poporange,line width=1.8pt,no markers,smooth},
}
% Même style disponible dans un \draw TikZ simple (hors pgfplots)
\tikzset{popcurve/.style={poporange,line width=1.8pt}}
\tikzset{popgrid/.style={popline,very thin}}
\colorlet{popcurve}{poporange} % le nom du style est parfois utilisé comme couleur

% ============================================================
% Pill / squiggle
% ============================================================
\newcommand{\poppill}[3]{%
  \tikz[baseline=-0.55ex]\node[draw=popink,line width=1.2pt,rounded corners=2.6pt,%
    fill=#2,inner xsep=6.5pt,inner ysep=3pt]{%
    \popxboldls\fontsize{7}{7}\selectfont\color{#3}\MakeUppercase{#1}};%
}
\newcommand{\popsquiggle}[1]{%
  \tikz[baseline=-0.4ex]{
    \draw[#1,line width=2.6pt,line cap=round]
      plot [smooth] coordinates {(0,0.09) (0.34,0.01) (0.68,0.21) (1.02,0.01) (1.36,0.10)};
  }%
}
% Mot-clé surligné (marqueur orange sous le texte)
\newcommand{\cle}[1]{{\popxbold\color{popdark}#1}}

% ============================================================
% En-tête / pied
% ============================================================
\pagestyle{fancy}
\fancyhf{}
\renewcommand{\headrulewidth}{0pt}
\renewcommand{\footrulewidth}{0pt}
\lhead{\raisebox{-0.15ex}{\poplogo\fontsize{13}{13}\selectfont\color{popink}OnBuch\color{poporange}+}}
\rhead{\poppill{\DOCMATIERE\ \textbullet\ \DOCNIVEAU\ \textbullet\ Leçon \DOCLECON}{white}{popink}}
\lfoot{\popxbold\fontsize{7.6}{7.6}\selectfont\color{popmuted}\MakeUppercase{Cours signé OnBuch+}\ \textbullet\ {\popsemi\DOCTITRE}}
\rfoot{\tikz[baseline=-0.6ex]\node[rounded corners=8.5pt,fill=popink,minimum height=17pt,minimum width=17pt,inner xsep=5pt,inner sep=0]{\popxbold\fontsize{8}{8}\selectfont\color{popcream}\thepage\,/\,\pageref*{LastPage}};}

% ============================================================
% Titres
% ============================================================
\newcommand{\chaptitre}[2]{%
  \begin{tcolorbox}[enhanced,colback=poporange,colframe=popink,boxrule=2.6pt,arc=11pt,
      drop shadow=popink,left=15pt,right=15pt,top=12pt,bottom=11pt,before skip=3pt,after skip=18pt]
    \poppill{#2}{popink}{popcream}\par\vspace{9pt}
    {\raggedright\hyphenpenalty=10000\exhyphenpenalty=10000\popdisplay\fontsize{22}{24}\selectfont\color{popcream}\MakeUppercase{#1}\par}
  \end{tcolorbox}
}
% Section numérotée : pastille orange avec le numéro + titre Archivo
\newcounter{popsec}
\newcommand{\coursec}[1]{%
  \stepcounter{popsec}\setcounter{popsub}{0}%
  \par\vspace{16pt}\noindent
  \begin{minipage}{\linewidth}
  \tikz[baseline=-0.7ex]\node[draw=popink,line width=1.6pt,fill=poporange,rounded corners=4pt,minimum size=22pt,inner sep=0]{\popdisplay\fontsize{12}{12}\selectfont\color{popcream}\thepopsec};%
  \hspace{8pt}{\popdisplay\fontsize{15}{17}\selectfont\color{popink}\MakeUppercase{#1}}\par
  \vspace{4pt}\noindent{\color{poporange}\rule{36pt}{3.6pt}}
  \end{minipage}\par\nopagebreak\vspace{8pt}
}
\newcounter{popsub}
\newcommand{\courssub}[1]{%
  \stepcounter{popsub}%
  \par\vspace{9pt}\noindent{\popxbold\fontsize{11.5}{13}\selectfont\color{popink}{\color{poporange}\thepopsec.\thepopsub}\hspace{6pt}#1}\par\nopagebreak\vspace{4pt}
}

% ============================================================
% Boîtes pédagogiques — même recette : fond blanc, bordure épaisse colorée,
% ombre dure encre, badge plein. Titre optionnel [#1].
% ============================================================
\tcbset{popbase/.style={breakable,enhanced,colback=white,boxrule=2.3pt,arc=9pt,
  shadow={3pt}{-3pt}{0pt}{popink},left=14pt,right=14pt,top=11pt,bottom=11pt,before skip=11pt,after skip=15pt,
  fontupper=\popsemi\fontsize{10.6}{14.5}\selectfont\color{popink},
  fontlower=\popsemi\fontsize{10.6}{14.5}\selectfont\color{popink}}}
% Coche dessinée (le glyphe ✓ n'existe pas dans les polices du document)
\newcommand{\popcheck}[1]{\tikz[baseline=-0.5ex]\draw[#1,line width=1.6pt,line cap=round,line join=round](-0.13,0.0)--(-0.04,-0.09)--(0.13,0.1);}
\providecommand{\coche}{\popcheck{popgreen}}
\providecommand{\resultat}[1]{\par\smallskip\noindent\fcolorbox{popgreen}{popgreenL}{\parbox{\dimexpr\linewidth-2\fboxsep-2\fboxrule\relax}{\textbf{Résultat.} #1}}\par\smallskip}
\newcommand{\popboxhead}[3]{\poppill{#1}{#2}{white}\ifx\relax#3\relax\else\hspace{6pt}{\popxbold\fontsize{10.6}{12}\selectfont\color{popink}#3}\fi\par\vspace{7pt}}

\newtcolorbox{aretenir}[1][]{popbase,colframe=popgreen,before upper={\popboxhead{À retenir}{popgreen}{#1}}}
% Texte en chinois (document de lecture, dialogue, extrait) : cadre
% d'étude. Un titre optionnel (source, auteur) s'affiche dans l'en-tête.
\newtcolorbox{textebox}[1][]{popbase,colframe=popblue,before upper={\popboxhead{文本 · Texte}{popblue}{#1}},after upper={}}
% Marques ✓ / ✗ (forme correcte / fautive) souvent utilisées par les modèles
\providecommand{\cross}{\textcolor{poppink}{\ensuremath{\times}}}
\providecommand{\xmark}{\cross}\providecommand{\crossmark}{\cross}\providecommand{\cmark}{\ensuremath{\checkmark}}
% Ligne de réponse / justification à compléter (inventée par les modèles)
\providecommand{\justification}[1]{\par\noindent\textit{Justification~:} \dotfill\par}
% \textipa (package tipa non chargé) : la police ne couvre pas l'API -> simple passage
\providecommand{\textipa}[1]{#1}
% Soulignement (ulem non chargé) : \uline, \ul
\providecommand{\uline}[1]{\underline{#1}}\providecommand{\ul}[1]{\underline{#1}}
% \glossaire{\item ...} inventée par les modèles : liste de glossaire
\providecommand{\glossaire}[1]{\begin{itemize}#1\end{itemize}}
% \alert (beamer) inventée par les modèles : texte mis en évidence
\providecommand{\alert}[1]{\textbf{\textcolor{poppink}{#1}}}
% variantes ASCII d'ordinaux écrites en macro
\providecommand{\iere}{\textsuperscript{re}}\providecommand{\ieme}{\textsuperscript{e}}\providecommand{\ere}{\textsuperscript{re}}\providecommand{\eme}{\textsuperscript{e}}\providecommand{\er}{\textsuperscript{er}}\providecommand{\ie}{\textsuperscript{e}}
% Ordinaux français écrits en macro par les modèles : 1\ière, 3\ième
\newcommand{\ière}{\textsuperscript{re}}\newcommand{\ième}{\textsuperscript{e}}
% \exemple (inventée par les modèles) : étiquette « Exemple : »
\providecommand{\exemple}{\textbf{Exemple~:}\ }
% \square utilisable aussi hors mode math (cases à cocher)
\AtBeginDocument{\let\squareMath\square\renewcommand{\square}{\ensuremath{\squareMath}}}
% Mot ou phrase en chinois à l'intérieur d'un paragraphe français
\newcommand{\zh}[1]{\ifmmode\text{#1}\else{#1}\fi}
\let\eng\zh \let\esp\zh \let\all\zh % alias tolérés
% flèches d'intonation inventées par les modèles
\providecommand{\textsearrow}{}\renewcommand{\textsearrow}{\ensuremath{\searrow}}\providecommand{\textnearrow}{}\renewcommand{\textnearrow}{\ensuremath{\nearrow}}
% \fr{..} : mot français (inventé par les modèles)
\providecommand{\fr}[1]{\textit{#1}}
% zhbox : encadré de texte chinois inventé par les modèles
\newenvironment{zhbox}{\begin{quote}}{\end{quote}}
% \vs : « versus » inventé par les modèles
\providecommand{\vs}{\textit{vs}\ }
% \blank : case à compléter inventée par les modèles
\providecommand{\blank}[1][]{\underline{\hspace{1.6em}}}
\newtcolorbox{exemplebox}[1][]{popbase,colframe=poporange,before upper={\popboxhead{Exemple}{poporange}{#1}}}
\newtcolorbox{definition}[1][]{popbase,colframe=popblue,before upper={\popboxhead{Définition}{popblue}{#1}}}
\newtcolorbox{propriete}[1][]{popbase,colframe=poppurple,before upper={\popboxhead{Propriété}{poppurple}{#1}}}
\newtcolorbox{methode}[1][]{popbase,colframe=popgold,before upper={\popboxhead{Méthode}{popgold}{#1}}}
\newtcolorbox{attention}[1][]{popbase,colframe=poppink,before upper={\popboxhead{Attention}{poppink}{#1}}}
\newtcolorbox{experience}[1][]{popbase,colframe=popblue,colback=popblueL,before upper={\popboxhead{Expérience}{popblue}{#1}}}
% « Le savais-tu ? » : carte violette pleine (contexte, culture, Cameroun)
\newtcolorbox{savaistu}[1][]{popbase,colback=poppurple,colframe=popink,
  fontupper=\popsemi\fontsize{10.4}{14.2}\selectfont\color{white},
  before upper={\poppill{Le savais-tu\,?}{white}{poppurple}\ifx\relax#1\relax\else\hspace{6pt}{\popxbold\color{white}#1}\fi\par\vspace{7pt}}}
% Exercice résolu : énoncé en haut, \tcblower puis la solution sur fond crème
\newcounter{popexo}\newcounter{popexr}
\newtcolorbox{exoresolu}[1][]{popbase,colframe=poporange,colbacklower=popcream,
  segmentation style={popink,line width=1.2pt,dashed},
  before upper={\stepcounter{popexr}\popboxhead{Exercice résolu \thepopexr}{poporange}{#1}},
  before lower={\poppill{Solution}{popink}{popcream}\par\vspace{6pt}}}
% Exercice (non résolu en place) — la correction est dans la partie Corrigés
\newtcolorbox{exercice}[1][]{popbase,colframe=popink,
  before upper={\stepcounter{popexo}\popboxhead{Exercice \thepopexo}{popink}{#1}}}
% Corrigé
\newtcolorbox{corrige}[1][]{popbase,colframe=popgreen,colback=popgreenL,
  before upper={\popboxhead{Corrigé}{popgreen}{#1}}}
% Objectifs / prérequis / bilan
\newtcolorbox{objectifs}{popbase,colframe=popink,colback=poporangeL,
  before upper={\popboxhead{Objectifs de la leçon}{popink}{}}}
\newtcolorbox{prerequis}{popbase,colframe=popink,colback=popblueL,
  before upper={\popboxhead{Prérequis}{popblue}{}}}
\newtcolorbox{bilan}{popbase,colframe=popink,colback=white,boxrule=2.8pt,
  before upper={\popboxhead{Fiche bilan}{poporange}{L'essentiel à connaître pour l'examen}}}
% Figure encadrée avec légende
\newtcolorbox{popfigure}[1][]{enhanced,breakable=false,colback=white,colframe=popline,boxrule=1.4pt,arc=8pt,
  left=10pt,right=10pt,top=10pt,bottom=8pt,before skip=10pt,after skip=12pt,halign=center,
  lower separated=false,halign lower=center,
  fontlower=\popsemi\fontsize{9}{11.5}\selectfont\color{popmuted},
  #1}
\newcommand{\legende}[1]{\par\vspace{6pt}{\popsemi\fontsize{9}{11.5}\selectfont\color{popmuted}#1}}

% ============================================================
% Signature OnBuch+
% ============================================================
\newcommand{\poplogoinline}[1]{{\poplogo\fontsize{#1}{#1}\selectfont\color{popink}OnBuch\color{poporange}+}}
\newcommand{\signatureonbuch}{%
  \begin{tcolorbox}[enhanced,colback=popink,colframe=popink,boxrule=2.2pt,arc=10pt,
      drop shadow=poporange,left=14pt,right=14pt,top=10pt,bottom=10pt,before skip=0pt,after skip=0pt]
    \tikz[baseline=-0.7ex]\node[circle,fill=popgreen,draw=popcream,line width=1.3pt,minimum size=22pt,inner sep=0]{\popcheck{white}};%
    \hspace{9pt}\begin{minipage}[c]{0.62\linewidth}
      {\popxbold\fontsize{12}{13}\selectfont\color{popcream}Signé\ \textbullet\ {\poplogo OnBuch\color{poporange}+}}\par\vspace{2pt}
      {\raggedright\popsemi\fontsize{8}{10.5}\selectfont\color{popline}\MakeUppercase{Équipe pédagogique OnBuch+}\\\MakeUppercase{Conforme au programme officiel MINESEC}\par}
    \end{minipage}\hfill
    {\poplogo\fontsize{22}{22}\selectfont\color{popcream}OnBuch\color{poporange}+}
  \end{tcolorbox}%
}

% ============================================================
% Page de garde
% #1 titre · #2 matière · #3 niveau · #4 module · #5 n° leçon · #6 durée ·
% #7 description / objectifs (texte ou itemize)
% ============================================================
\newcommand{\pagegarde}[7]{%
  \thispagestyle{empty}
  \newgeometry{a4paper,margin=1.7cm,top=1.6cm,bottom=1.4cm}%
  \begin{tikzpicture}[remember picture,overlay]
    \fill[poporange!18] ([xshift=-3.2cm,yshift=-3.6cm]current page.north east) circle (4.6cm);
    \fill[poppurple!14] ([xshift=2.4cm,yshift=7.6cm]current page.south west) circle (3.8cm);
  \end{tikzpicture}%
  {\poplogo\fontsize{30}{30}\selectfont\color{popink}OnBuch\color{poporange}+}\hfill
  \raisebox{0.6ex}{\poppill{Cours signé \textbullet\ Premium}{popink}{popcream}}\par
  \vspace{1pt}\popsquiggle{poppurple}\par
  \vspace{8pt}
  \begin{tcolorbox}[enhanced,colback=poporange,colframe=popink,boxrule=2.8pt,arc=13pt,
      drop shadow=popink,left=18pt,right=18pt,top=12pt,bottom=13pt,after skip=10pt]
    \poppill{#2 \textbullet\ #3}{popink}{popcream}\hspace{5pt}\poppill{#4}{white}{popink}\par\vspace{8pt}
    {\popdisplay\fontsize{14}{14}\selectfont\color{popink}LEÇON #5}\par\vspace{4pt}
    {\raggedright\hyphenpenalty=10000\exhyphenpenalty=10000\popdisplay\fontsize{25}{27}\selectfont\color{popcream}\MakeUppercase{#1}\par}
    \vspace{8pt}
    {\popxbold\fontsize{9.5}{9.5}\selectfont\color{popink}\MakeUppercase{Durée conseillée : #6}}
  \end{tcolorbox}
  \begin{tcolorbox}[enhanced,colback=white,colframe=popink,boxrule=2.2pt,arc=10pt,
      drop shadow=popink,left=16pt,right=16pt,top=10pt,bottom=10pt,after skip=8pt]
    \poppill{Dans cette leçon}{poppurple}{white}\par\vspace{5pt}
    \popsemi\fontsize{9.3}{12.6}\selectfont\color{popink}#7
  \end{tcolorbox}
  \noindent\begin{minipage}[t]{0.487\linewidth}
    \begin{tcolorbox}[enhanced,colback=popgreen,colframe=popink,boxrule=2.2pt,arc=10pt,
        drop shadow=popink,left=11pt,right=11pt,top=10pt,bottom=10pt,height=3.1cm,valign=center]
      \tcbox[colback=white,colframe=popink,boxrule=1.3pt,arc=5pt,boxsep=0pt,nobeforeafter,
        left=3pt,right=3pt,top=3pt,bottom=3pt,valign=center]{\qrcode[height=1.9cm]{https://play.google.com/store/apps/details?id=com.luvvix.onbuch&hl=fr}}%
      \hspace{8pt}\begin{minipage}[c]{0.5\linewidth}
        {\popdisplay\fontsize{11}{12.5}\selectfont\color{white}\MakeUppercase{Télécharge OnBuch}}\par\vspace{4pt}
        {\raggedright\popsemi\fontsize{8.4}{11}\selectfont\color{white}Gratuit \textbullet\ Google Play\\Tuteur IA, annales, concours, résultats\par}
      \end{minipage}
    \end{tcolorbox}
  \end{minipage}\hfill
  \begin{minipage}[t]{0.487\linewidth}
    \begin{tcolorbox}[enhanced,colback=poppurple,colframe=popink,boxrule=2.2pt,arc=10pt,
        drop shadow=popink,left=11pt,right=11pt,top=10pt,bottom=10pt,height=3.1cm,valign=center]
      \tikz[baseline=-0.7ex]\node[circle,fill=white,draw=popink,line width=1.3pt,minimum size=1.6cm,inner sep=0]{\popdisplay\fontsize{24}{24}\selectfont\color{poppurple}+};%
      \hspace{8pt}\begin{minipage}[c]{0.6\linewidth}
        {\popdisplay\fontsize{11}{12.5}\selectfont\color{white}\MakeUppercase{Passe à OnBuch+}}\par\vspace{4pt}
        {\raggedright\popsemi\fontsize{8.4}{11}\selectfont\color{white}Tous les cours signés, Tuteur IA illimité, fiches PDF — onglet OnBuch+ de l'app\par}
      \end{minipage}
    \end{tcolorbox}
  \end{minipage}
  \vspace{8pt}
  \popcontenu
  \vfill
  \signatureonbuch
  \restoregeometry
  \newpage
}

% Rangée « Ce cours contient » (page de garde)
\newcommand{\poptile}[3]{% #1 couleur #2 gros texte #3 légende
  \begin{tcolorbox}[enhanced,colback=#1,colframe=popink,boxrule=2pt,arc=9pt,shadow={2.5pt}{-2.5pt}{0pt}{popink},
      left=6pt,right=6pt,top=9pt,bottom=9pt,halign=center,valign=center,height=1.75cm,width=\linewidth,nobeforeafter]
    {\popdisplay\fontsize{12.5}{13}\selectfont\color{white}\MakeUppercase{#2}}\par\vspace{3pt}
    {\centering\popsemi\fontsize{7.8}{9.5}\selectfont\color{white}#3\par}
  \end{tcolorbox}}
\newcommand{\popcontenu}{%
  \par\noindent{\popxboldls\fontsize{7.5}{7.5}\selectfont\color{popmuted}CE COURS SIGNÉ CONTIENT}\par\vspace{7pt}
  \noindent\begin{minipage}[t]{0.235\linewidth}\poptile{popblue}{Cours}{complet, illustré, pas à pas}\end{minipage}\hfill
  \begin{minipage}[t]{0.235\linewidth}\poptile{poporange}{Méthodes}{et exercices résolus}\end{minipage}\hfill
  \begin{minipage}[t]{0.235\linewidth}\poptile{poppink}{Exercices}{type examen + corrigés}\end{minipage}\hfill
  \begin{minipage}[t]{0.235\linewidth}\poptile{popgreen}{Fiche bilan}{l'essentiel à retenir}\end{minipage}\par}

% Sommaire
\newtcolorbox{sommaire}{enhanced,colback=white,colframe=popink,boxrule=2.2pt,arc=10pt,
  drop shadow=popink,left=18pt,right=18pt,top=13pt,bottom=13pt,before skip=6pt,after skip=20pt,
  before upper={\poppill{Sommaire}{poppurple}{white}\par\vspace{9pt}\popsemi\fontsize{10.6}{14.5}\selectfont\color{popink}}}

% Signature de fin de cours — poussée tout en bas de la dernière page
\newcommand{\findecours}{%
  \par\vfill\nopagebreak
  \begin{center}\popsquiggle{poporange}\end{center}\vspace{4pt}
  \signatureonbuch
}
\AtBeginDocument{\def\llbracket{\mathopen{[\mkern-3.5mu[}}\def\rrbracket{\mathclose{]\mkern-3.5mu]}}} % intervalles d'entiers (glyphe ⟦ absent de la police)
% Glyphes absents de Fira Math : redéfinis à partir de glyphes disponibles
\AtBeginDocument{%
  \def\setminus{\mathbin{\backslash}}\let\smallsetminus\setminus
  \def\vdots{\mathord{\vbox{\baselineskip3.2pt\lineskiplimit0pt\kern4pt\hbox{.}\hbox{.}\hbox{.}}}}%
  \def\ddots{\mathinner{\mkern1mu\raise6.5pt\hbox{.}\mkern2mu\raise3.6pt\hbox{.}\mkern2mu\raise0.7pt\hbox{.}\mkern1mu}}%
  \def\triangle{\mathord{\scalebox{0.9}{$\symup{\Delta}$}}}%
  \def\nabla{\mathord{\raisebox{\depth}{\scalebox{1}[-1]{$\symup{\Delta}$}}}}%
  \def\checkmark{\popcheck{popdark}}%
  \def\star{\mathbin{\ast}}\def\bigstar{\mathord{\ast}}%
  \def\diamond{\mathbin{\scalebox{0.6}{$\Diamond$}}}%
  \def\bigcup{\mathop{\mathchoice{\raisebox{-0.3em}{\scalebox{1.7}{$\cup$}}}{\scalebox{1.25}{$\cup$}}{\cup}{\cup}}\displaylimits}%
  \def\bigcap{\mathop{\mathchoice{\raisebox{-0.3em}{\scalebox{1.7}{$\cap$}}}{\scalebox{1.25}{$\cap$}}{\cap}{\cap}}\displaylimits}%
  \def\lesseqqgtr{\mathrel{\lessgtr}}\def\bigoplus{\mathop{\oplus}\displaylimits}%
}
\providecommand{\ssi}{si et seulement si} % abréviation fréquente
\AtBeginDocument{\providecommand{\wideparen}{\overparen}} % arc de cercle

%% FIGFIX-BEGIN v5 — correctifs globaux des figures (docs/onbuch-generation/tools/figfix)
\usepackage{adjustbox}\usepackage{etoolbox}\tcbuselibrary{raster}
% toute figure TikZ/pgfplots plus large que la ligne est réduite à la largeur disponible
\BeforeBeginEnvironment{tikzpicture}{\begin{adjustbox}{max width=\linewidth}}
\AfterEndEnvironment{tikzpicture}{\end{adjustbox}}
% légendes pgfplots lisibles par-dessus les courbes
\pgfplotsset{every axis/.append style={legend style={fill=white,fill opacity=0.92,text opacity=1,draw=black!15,inner sep=3pt}}}
% étiquettes d'arêtes (syntaxe edge["..."]) avec fond blanc
\tikzset{every edge quotes/.append style={fill=white,inner sep=1.2pt}}
%% FIGFIX-END
\begin{document}

\pagegarde{Leçon 8 — Éléments socioculturels : processus d'hérédité au Cameroun et en Chine (suite : comparaison et exposé)}{Chinois}{Tle A}{Module 1 : Vie familiale et intégration sociale}{ZH08}{2 h}{Cette leçon socioculturelle prolonge la leçon 4 en proposant une comparaison factuelle et respectueuse des processus d'hérédité au Cameroun et en Chine. Elle fournit le vocabulaire chinois du thème et prépare les élèves à un mini-exposé comparatif oral.
\vspace{4pt}
\begin{multicols}{2}\small
\begin{itemize}
\item À la fin de cette leçon, je sais décrire les grandes étapes du processus d'hérédité au Cameroun (succession coutumière et droit moderne)
\item À la fin de cette leçon, je sais décrire le processus d'hérédité en Chine (loi sur la succession, testament, partage entre enfants)
\item À la fin de cette leçon, je sais comparer les deux systèmes de façon factuelle et respectueuse
\item À la fin de cette leçon, je sais employer le vocabulaire socioculturel du thème en caractères et en pinyin
\item À la fin de cette leçon, je sais utiliser les structures de comparaison 比 et 和……一样/不一样
\item À la fin de cette leçon, je sais exprimer un point de vue simple avec 我觉得 / 我认为
\end{itemize}\end{multicols}}

\begin{sommaire}
\begin{multicols}{2}
\begin{enumerate}
\item Rappel ciblé et mise en route
\item L'hérédité au Cameroun : faits
\item L'hérédité en Chine : faits
\item Vocabulaire du thème
\item Comparer et donner son avis en chinois
\item Questions de réflexion et préparation du mini-exposé
\item Activité d'intégration
\item Méthodes et exercices résolus
\item Exercices
\item Corrigés des exercices
\item Fiche bilan
\end{enumerate}
\end{multicols}
\end{sommaire}

\begin{objectifs}
\begin{itemize}
\item À la fin de cette leçon, je sais décrire les grandes étapes du processus d'hérédité au Cameroun (succession coutumière et droit moderne)
\item À la fin de cette leçon, je sais décrire le processus d'hérédité en Chine (loi sur la succession, testament, partage entre enfants)
\item À la fin de cette leçon, je sais comparer les deux systèmes de façon factuelle et respectueuse
\item À la fin de cette leçon, je sais employer le vocabulaire socioculturel du thème en caractères et en pinyin
\item À la fin de cette leçon, je sais utiliser les structures de comparaison 比 et 和……一样/不一样
\item À la fin de cette leçon, je sais exprimer un point de vue simple avec 我觉得 / 我认为
\item À la fin de cette leçon, je sais répondre à des questions de réflexion sur les faits culturels
\item À la fin de cette leçon, je sais préparer et présenter un mini-exposé comparatif en chinois
\end{itemize}
\end{objectifs}

\begin{prerequis}
\begin{itemize}
\item Leçon 4 : notions de base sur la famille et l'hérédité (vocabulaire familial déjà acquis)
\item Structure comparative 比 (bǐ) vue en classe de Première
\item Vocabulaire de la famille : 父亲, 母亲, 儿子, 女儿, 家
\item Expression de l'opinion simple : 我觉得, 因为……所以……
\item Nombres et possessifs (的) pour décrire le partage des biens
\end{itemize}
\end{prerequis}

\newpage

\coursec{Rappel ciblé et mise en route}

\courssub{Situation d'accroche : deux mondes, une même question}

Imaginons deux scènes. À Yaoundé, dans le quartier Mvolyé, un chef de famille vient de décéder. Le soir même, la maison se remplit : oncles, tantes, notables du village, voisins. On ne parle pas seulement de tristesse ; on parle aussi d'avenir. Qui héritera de la maison familiale ? Qui recevra les champs de cacao ? Qui portera le titre et les responsabilités du défunt ? Le fils aîné ? Le frère du défunt ? Parfois même la veuve ? Toute la communauté participe à la décision.

Des milliers de kilomètres plus loin, à Pékin, un jeune étudiant camerounais en échange universitaire discute avec ses camarades chinois. Il découvre avec surprise qu'en Chine, on parle d'un \cle{testament écrit}, d'une \cle{loi nationale sur l'héritage} votée en 1985, et que filles et garçons ont, en principe, les mêmes droits devant la loi.

\begin{textebox}[Dialogue : un échange entre étudiants]
\textbf{李明：} 阿里，在你们国家，人去世以后，谁继承他的房子？\\
Lǐ Míng : Ālǐ, zài nǐmen guójiā, rén qùshì yǐhòu, shéi jìchéng tā de fángzi?\\
« Li Ming : Ali, dans ton pays, quand quelqu'un meurt, qui hérite de sa maison ? »\\[4pt]
\textbf{阿里：} 在我们那儿，家人和酋长一起决定。常常是长子继承。\\
Ālǐ : Zài wǒmen nàr, jiārén hé qiúzhǎng yìqǐ juédìng. Chángcháng shì zhǎngzǐ jìchéng.\\
« Ali : Chez nous, la famille et le chef décident ensemble. C'est souvent le fils aîné qui hérite. »\\[4pt]
\textbf{李明：} 在中国，我们有继承法。儿子和女儿都可以继承。\\
Lǐ Míng : Zài Zhōngguó, wǒmen yǒu jìchéngfǎ. Érzi hé nǚ'ér dōu kěyǐ jìchéng.\\
« Li Ming : En Chine, nous avons une loi sur l'héritage. Les fils et les filles peuvent tous hériter. »\\[4pt]
\textbf{阿里：} 真有意思！两个国家不一样，但是都很有意思。\\
Ālǐ : Zhēn yǒu yìsi! Liǎng ge guójiā bù yīyàng, dànshì dōu hěn yǒu yìsi.\\
« Ali : C'est très intéressant ! Les deux pays sont différents, mais tous les deux sont intéressants. »
\end{textebox}

\textbf{Problématique de la leçon.} Comment deux grandes cultures, le Cameroun et la Chine, organisent-elles la transmission des biens, des noms et des responsabilités d'une génération à l'autre ? Et surtout : comment en parler en chinois, de façon correcte, factuelle et respectueuse ? Aujourd'hui, nous comparons et nous apprenons à exprimer une comparaison et un avis personnel.

\courssub{Réactivation rapide de la leçon 4 (5 minutes)}

Avant d'aller plus loin, rappelons en quelques mots ce que nous avons découvert dans la leçon 4. Nous ne répétons pas le contenu : nous le réactivons pour construire dessus.

\begin{aretenir}[Ce que nous retenons de la leçon 4]
\begin{itemize}
\item Au Cameroun, la \cle{famille élargie} (\zh{大家庭} \emph{dà jiātíng}) joue un rôle central : oncles, tantes, cousins font partie du cercle décisionnel.
\item La \cle{chefferie} (\zh{酋长} \emph{qiúzhǎng}, « le chef traditionnel ») arbitre souvent les questions de succession et de transmission.
\item En Chine aussi, la famille a une place forte, mais l'État et la loi encadrent davantage les pratiques.
\end{itemize}
\end{aretenir}

\textbf{Question flash à la classe.} Qui peut rappeler, en une phrase en chinois, le rôle du chef traditionnel ? Réponse attendue, par exemple : \zh{酋长帮助家人决定。} \emph{Qiúzhǎng bāngzhù jiārén juédìng.} « Le chef aide la famille à décider. »

\courssub{Brainstorming : que se passe-t-il chez vous ?}

\begin{experience}[Brainstorming en français (10 minutes)]
En français, répondez librement à cette question : \textbf{que se passe-t-il dans votre famille ou votre village quand un grand-parent décède ?}
\begin{enumerate}
\item Qui se réunit ? Où ? Pendant combien de temps ?
\item Qui décide de la répartition des biens (maison, champs, argent, objets) ?
\item Y a-t-il des cérémonies ? Lesquelles ?
\item Les filles héritent-elles comme les garçons dans votre expérience ?
\end{enumerate}
Le professeur note au tableau les mots-clés qui reviennent : \emph{famille, chef, réunion, partage, cérémonie, fils aîné, veuve}. Ces mots seront traduits en chinois dans la section vocabulaire.
\end{experience}

Ce brainstorming a un objectif précis : partir de \textbf{votre expérience réelle} avant de découvrir les faits chinois. En didactique des langues, on appelle cela activer les \cle{connaissances préalables} : on apprend mieux une culture étrangère quand on a d'abord mis des mots sur la sienne.

\courssub{Le mot-clé du jour : 继承}

\begin{savaistu}[D'où vient le mot « hériter » en chinois ?]
En chinois, « hériter » se dit \zh{继承} \emph{jìchéng}. Ce mot est composé de deux caractères : \zh{继} \emph{jì}, « continuer, succéder », et \zh{承} \emph{chéng}, « recevoir, porter, assumer ». Littéralement, hériter, c'est donc « \cle{continuer et recevoir} » : on reçoit un bien, mais on reçoit aussi une responsabilité à faire vivre. Belle idée, non ? Le nom correspondant est \zh{继承} ou \zh{遗产} \emph{yíchǎn} pour « l'héritage, le patrimoine » (les biens transmis), et \zh{继承人} \emph{jìchéngrén} désigne « l'héritier ».
\end{savaistu}

\begin{center}
\begin{tabularx}{\linewidth}{|X|X|X|X|}
\hline
\rowcolor{popblueL} Caractères & Pinyin & Nature & Traduction \\
\hline
继承 & jìchéng & verbe & hériter, succéder \\
\hline
遗产 & yíchǎn & nom & héritage, patrimoine (biens) \\
\hline
继承人 & jìchéngrén & nom & héritier, héritière \\
\hline
去世 & qùshì & verbe & décéder (registre respectueux) \\
\hline
决定 & juédìng & verbe & décider \\
\hline
不一样 & bù yīyàng & expression & différent, pas pareil \\
\hline
\end{tabularx}
\end{center}

\begin{attention}[Décès : choisir le bon mot]
En français, « mourir » et « décéder » n'ont pas le même registre. En chinois, c'est pareil : \zh{死} \emph{sǐ} est direct, voire brutal ; dans un contexte familial respectueux, on préfère \zh{去世} \emph{qùshì}, littéralement « quitter le monde ». Dans votre exposé, dites : \zh{爷爷去世以后…} \emph{Yéye qùshì yǐhòu…} « Après le décès de grand-père… ». Autre piège : le chiffre 4 (\zh{四} \emph{sì}) est évité dans les cadeaux et les étages en Chine car il sonne presque comme \zh{死} \emph{sǐ} — un détail culturel à retenir !
\end{attention}

\courssub{La règle d'or : comparer sans juger}

Avant toute comparaison culturelle, fixons ensemble une règle de respect. En classe de chinois, nous ne disons jamais qu'un système est « meilleur » ou « moins bon » qu'un autre. Nous disons qu'il est \cle{différent}.

\begin{attention}[Le mot magique : 不一样]
La phrase-clé de toute cette leçon est :\\
\zh{两个国家不一样。} \emph{Liǎng ge guójiā bù yīyàng.} « Les deux pays sont différents. »\\
Structure : \textbf{A 和 B 不一样。} \emph{A hé B bù yīyàng.} « A et B ne sont pas pareils. »\\
Interdit en classe : les formulations du type « leur système est mauvais / meilleur ». Obligatoire : décrire les faits d'abord, puis donner un avis personnel argumenté et respectueux.
\end{attention}

\begin{propriete}[Structure : comparer avec 一样 / 不一样]
\begin{itemize}
\item \textbf{Égalité :} A 和 B 一样 (+ adjectif)。\emph{A hé B yīyàng.} « A et B sont pareils. »
\item \textbf{Différence :} A 和 B 不一样。\emph{A hé B bù yīyàng.} « A et B sont différents. »
\item \textbf{Comparaison de supériorité (rappel) :} A 比 B + adjectif。\emph{A bǐ B…} « A est plus… que B. »
\end{itemize}
Exemples :\\
\zh{中国和喀麦隆的家庭不一样。} \emph{Zhōngguó hé Kāmàilóng de jiātíng bù yīyàng.} « Les familles chinoises et camerounaises sont différentes. »\\
\zh{两个国家都很重视家庭。} \emph{Liǎng ge guójiā dōu hěn zhòngshì jiātíng.} « Les deux pays attachent tous les deux beaucoup d'importance à la famille. »
\end{propriete}

\courssub{Méthode : mener une comparaison culturelle en quatre étapes}

\begin{methode}[Observer → Décrire → Comparer → Donner un avis]
Pour réussir votre mini-exposé de fin de leçon, suivez toujours ces quatre étapes :
\begin{enumerate}
\item \textbf{Observer} : lire le texte, écouter le dialogue, repérer les faits (qui ? quoi ? comment ?).
\item \textbf{Décrire} : formuler les faits en phrases simples, sans jugement. Exemple : \zh{在喀麦隆，家人和酋长一起决定。} « Au Cameroun, la famille et le chef décident ensemble. »
\item \textbf{Comparer} : utiliser \zh{一样} / \zh{不一样} / \zh{比}. Exemple : \zh{中国用继承法，喀麦隆常常用传统。} « La Chine utilise une loi sur l'héritage, le Cameroun utilise souvent la tradition. »
\item \textbf{Donner un avis argumenté} : \zh{我觉得…，因为…} \emph{Wǒ juéde…, yīnwèi…} « Je pense que…, parce que… ». Exemple : \zh{我觉得两个系统都很有意思，因为两个系统都保护家庭。} « Je trouve les deux systèmes intéressants, parce que tous les deux protègent la famille. »
\end{enumerate}
\end{methode}

\courssub{Carte mentale : le thème de la leçon}

Voici la carte mentale qui organisera toute notre leçon : un concept central, deux branches à explorer.

\begin{popfigure}
\begin{center}\tcbox[colback=popblue,colframe=popblue,coltext=white,arc=9pt,boxsep=4pt,left=6pt,right=6pt]{\bfseries\small 继承 \emph{jìchéng} héritage}\end{center}
\vspace{-2pt}
\begin{tcbraster}[raster columns=2,raster equal height=rows,raster column skip=6pt,raster row skip=6pt,enhanced,blankest]
\begin{tcolorbox}[enhanced,colback=popblue,colframe=popblue,coltext=white,boxrule=0.9pt,arc=3pt,left=4pt,right=4pt,top=3pt,bottom=3pt,boxsep=1pt,halign=left]\bfseries\scriptsize 喀麦隆 \emph{Kāmàilóng} Cameroun famille élargie chefferie, tradition\end{tcolorbox}
\begin{tcolorbox}[enhanced,colback=poporange,colframe=poporange,coltext=white,boxrule=0.9pt,arc=3pt,left=4pt,right=4pt,top=3pt,bottom=3pt,boxsep=1pt,halign=left]\bfseries\scriptsize 中国 \emph{Zhōngguó} Chine loi sur l'héritage testament écrit\end{tcolorbox}
\end{tcbraster}

\legende{Carte mentale du thème : l'héritage au Cameroun et en Chine.}
\end{popfigure}

\courssub{Annonce du plan et des règles de travail}

\begin{aretenir}[Plan de la leçon]
\begin{enumerate}
\item \textbf{Rappel ciblé et mise en route} (nous y sommes).
\item \textbf{L'hérédité au Cameroun} : les faits, sans répéter la leçon 4.
\item \textbf{L'hérédité en Chine} : la loi, le testament, les pratiques.
\item \textbf{Vocabulaire du thème} : les mots pour en parler.
\item \textbf{Comparer et donner son avis en chinois} : les structures.
\item \textbf{Questions de réflexion et préparation du mini-exposé.}
\end{enumerate}
Règles de respect : on décrit des faits, on ne classe pas les cultures, on utilise \zh{不一样} et on argumente avec \zh{因为}.
\end{aretenir}

\courssub{Exercice résolu de mise en route}

\begin{exoresolu}[Première comparaison guidée]
Énoncé. À partir du dialogue entre Li Ming et Ali, complétez en chinois :\\
1. 在喀麦隆，\_\_\_\_ 和 \_\_\_\_ 一起决定。(Au Cameroun, \_\_\_\_ et \_\_\_\_ décident ensemble.)\\
2. 在中国，儿子和女儿 \_\_\_\_ 继承。(En Chine, fils et filles \_\_\_\_ hériter.)\\
3. Traduisez : « Les deux pays sont différents. »
\tcblower
Solution détaillée.\\
1. \zh{在喀麦隆，家人和酋长一起决定。} \emph{Zài Kāmàilóng, jiārén hé qiúzhǎng yìqǐ juédìng.} On relie deux sujets avec \zh{和} \emph{hé} « et », et \zh{一起} \emph{yìqǐ} « ensemble » se place avant le verbe.\\
2. \zh{在中国，儿子和女儿都可以继承。} \emph{Zài Zhōngguó, érzi hé nǚ'ér dōu kěyǐ jìchéng.} \zh{都} \emph{dōu} « tous les deux » se place après les deux sujets et avant \zh{可以} \emph{kěyǐ} « pouvoir ».\\
3. \zh{两个国家不一样。} \emph{Liǎng ge guójiā bù yīyàng.} Attention : avec \zh{两} « deux », le classificateur \zh{个} est obligatoire ; on ne dit jamais \zh{二个国家}.
\end{exoresolu}

\begin{attention}[Pièges fréquents des francophones]
\begin{itemize}
\item \textbf{Ordre des mots :} en français on dit « décider ensemble », en chinois \zh{一起} se place \emph{avant} le verbe : \zh{一起决定}, jamais \zh{决定一起} dans ce sens.
\item \textbf{和 ne relie que des noms :} on peut dire \zh{家人和酋长} (deux noms), mais pour relier deux phrases, le français « et » ne se traduit pas par \zh{和} ; on juxtapose simplement avec une virgule chinoise ，.
\item \textbf{两 et non 二 devant un classificateur :} \zh{两个国家} « deux pays », jamais \zh{二个}.
\item \textbf{Ne pas traduire « différent de » mot à mot :} « différent de » se dit avec la structure A 和 B 不一样, pas avec une préposition.
\end{itemize}
\end{attention}

Nous sommes maintenant prêts : le mot-clé \zh{继承} est acquis, la structure de comparaison est en place, et les règles de respect sont fixées. Passons aux faits : d'abord le Cameroun, ensuite la Chine.

\coursec{L'hérédité au Cameroun : faits}

Après le rappel ciblé de la section précédente — où nous avons réactivé le vocabulaire de la famille et les structures de comparaison vues à la leçon 4 — nous entrons maintenant dans le cœur du sujet : comment l'hérédité fonctionne-t-elle concrètement au Cameroun ? Cette section présente les faits, de façon rigoureuse et respectueuse, en vous donnant à chaque étape les mots chinois pour en parler. La Chine sera étudiée dans la section suivante ; ici, observons d'abord notre propre contexte, celui que vous connaissez le mieux.

\courssub{Observer d'abord : un texte sur la succession coutumière}

Avant toute explication, lisons un texte chinois original qui décrit, de manière simple et factuelle, ce qui se passe dans de nombreuses familles camerounaises après un décès. Lisez-le deux fois : une première fois pour le sens général, une deuxième fois en suivant le pinyin.

\begin{textebox}[在喀麦隆，人去世以后]
在喀麦隆，人去世以后，家人开会选择继承人。\\[4pt]
\emph{Zài Kāmàilóng, rén qùshì yǐhòu, jiārén kāihuì xuǎnzé jìchéngrén.}\\
« Au Cameroun, après le décès d'une personne, la famille se réunit pour choisir l'héritier. »\\[8pt]
在很多地方，家里的长辈和有名望的人一起商量。他们决定谁继承死者的名字、头衔和财产。财产可以是房子、田地，也可以是牲畜。\\
\emph{Zài hěn duō dìfang, jiā lǐ de zhǎngbèi hé yǒu míngwàng de rén yìqǐ shāngliang. Tāmen juédìng shéi jìchéng sǐzhě de míngzi, tóuxián hé cáichǎn. Cáichǎn kěyǐ shì fángzi, tiándì, yě kěyǐ shì shēngchù.}\\
« Dans beaucoup d'endroits, les aînés de la famille et les personnes respectées délibèrent ensemble. Ils décident qui hérite du nom, du titre et des biens du défunt. Les biens peuvent être une maison, des champs, ou encore du bétail. »\\[8pt]
在很多传统里，长子有很重要的地位：他常常继承父亲的头衔，也要照顾母亲和弟弟妹妹。\\
\emph{Zài hěn duō chuántǒng lǐ, zhǎngzǐ yǒu hěn zhòngyào de dìwèi : tā chángcháng jìchéng fùqin de tóuxián, yě yào zhàogù mǔqin hé dìdi mèimei.}\\
« Dans de nombreuses traditions, le fils aîné occupe une place très importante : il hérite souvent du titre du père et doit aussi prendre soin de sa mère et de ses cadets. »\\[8pt]
不过，喀麦隆的每个地区、每个民族的做法不完全一样。现在也有现代的办法：人们可以写遗嘱，也可以去公证处。所以，有时候一个家庭同时用传统的办法和现代的法律。\\
\emph{Búguò, Kāmàilóng de měi ge dìqū, měi ge mínzú de zuòfǎ bù wánquán yíyàng. Xiànzài yě yǒu xiàndài de bànfǎ : rénmen kěyǐ xiě yízhǔ, yě kěyǐ qù gōngzhèngchù. Suǒyǐ, yǒu shíhou yí ge jiātíng tóngshí yòng chuántǒng de bànfǎ hé xiàndài de fǎlǜ.}\\
« Cependant, les pratiques ne sont pas exactement les mêmes dans chaque région et chaque ethnie du Cameroun. Aujourd'hui, il existe aussi des moyens modernes : on peut rédiger un testament et se rendre chez le notaire. C'est pourquoi une famille utilise parfois à la fois la coutume et le droit moderne. »
\end{textebox}

\textbf{Questions de compréhension (répondez en français, puis essayez en chinois) :}
\begin{enumerate}
\item Qui participe à la réunion de famille après un décès ?
\item Quels biens sont mentionnés dans le texte ?
\item Quel rôle joue le fils aîné dans de nombreuses traditions ?
\item Quels sont les deux « systèmes » qui coexistent aujourd'hui au Cameroun ?
\end{enumerate}

\courssub{Le vocabulaire clé du texte}

\begin{definition}[继承人 — l'héritier]
\zh{继承人} (\emph{jìchéngrén}), nom : l'héritier, la personne qui reçoit les biens, le nom ou le titre d'une personne décédée. Le verbe correspondant est \zh{继承} (\emph{jìchéng}) : « hériter de, succéder à ». Structure : \cle{hériter de + nom} = \zh{继承 + nom}, sans préposition, contrairement au français « hériter \textbf{de} ».
\end{definition}

\begin{center}
\begin{tabularx}{\linewidth}{|X|X|X|X|}
\hline
\rowcolor{popblueL} Caractères & Pinyin & Nature & Traduction \\
\hline
去世 & qùshì & verbe & décéder (terme respectueux) \\
\hline
继承人 & jìchéngrén & nom & héritier, héritière \\
\hline
继承 & jìchéng & verbe & hériter (de), succéder (à) \\
\hline
长辈 & zhǎngbèi & nom & les aînés (de la famille) \\
\hline
商量 & shāngliang & verbe & délibérer, discuter ensemble \\
\hline
头衔 & tóuxián & nom & titre (chef, notable) \\
\hline
财产 & cáichǎn & nom & biens, patrimoine \\
\hline
田地 & tiándì & nom & champs, terres agricoles \\
\hline
牲畜 & shēngchù & nom & bétail, animaux d'élevage \\
\hline
长子 & zhǎngzǐ & nom & fils aîné \\
\hline
遗嘱 & yízhǔ & nom & testament \\
\hline
公证处 & gōngzhèngchù & nom & étude de notaire, office notarial \\
\hline
法律 & fǎlǜ & nom & loi, droit \\
\hline
传统 & chuántǒng & nom / adj. & tradition ; traditionnel \\
\hline
\end{tabularx}
\end{center}

\begin{attention}[Faux amis et pièges de vocabulaire]
\begin{itemize}
\item \zh{去世} (\emph{qùshì}) signifie « décéder » de façon respectueuse ; ne le confondez pas avec \zh{去} (\emph{qù}, « aller ») seul. On ne dit jamais \zh{他去了} pour « il est mort » dans un texte soigné : cela voudrait dire « il est parti ».
\item \zh{长子} (\zh{zhǎngzǐ}) se prononce avec \emph{zhǎng} (3\textsuperscript{e} ton, « aîné »), pas \emph{cháng} (« long »). Le même caractère 长 a deux lectures : piège classique pour les francophones !
\item \zh{继承} se construit \textbf{sans préposition} : on dit \zh{继承父亲的头衔} (litt. « hériter le titre du père »), jamais \zh{继承} + \zh{的} + titre.
\end{itemize}
\end{attention}

\courssub{La succession coutumière : la famille élargie et les notables}

Dans la plupart des traditions camerounaises, la succession n'est pas une affaire individuelle : c'est une affaire de \textbf{famille élargie}. Après le décès, la famille se réunit — on parle souvent de « réunion de famille » — avec les aînés (\zh{长辈}) et les notables, c'est-à-dire les personnes respectées du village ou du quartier. Ensemble, ils délibèrent (\zh{商量}) et désignent l'héritier (\zh{继承人}).

Ce que l'on transmet n'est pas seulement matériel. Trois types de choses sont transmis :
\begin{itemize}
\item le \textbf{nom} (\zh{名字}) : dans certaines traditions, un enfant de la famille reçoit le nom du défunt pour perpétuer sa mémoire ;
\item le \textbf{titre} (\zh{头衔}) : chef de famille, notable, parfois chef de village ou de quartier ;
\item les \textbf{biens} (\zh{财产}) : la maison (\zh{房子}), les champs (\zh{田地}), le bétail (\zh{牲畜}).
\end{itemize}

\begin{exemplebox}[Exemples traduits]
\begin{itemize}
\item \zh{长子继承父亲的头衔。}\\
\emph{Zhǎngzǐ jìchéng fùqin de tóuxián.}\\
« Le fils aîné hérite du titre du père. »
\item \zh{家人开会选择继承人。}\\
\emph{Jiārén kāihuì xuǎnzé jìchéngrén.}\\
« La famille se réunit pour choisir l'héritier. »
\item \zh{他继承了父亲的房子和田地。}\\
\emph{Tā jìchéng le fùqin de fángzi hé tiándì.}\\
« Il a hérité de la maison et des champs de son père. »
\item \zh{长辈们一起商量谁继承牲畜。}\\
\emph{Zhǎngbèimen yìqǐ shāngliang shéi jìchéng shēngchù.}\\
« Les aînés délibèrent ensemble pour décider qui hérite du bétail. »
\end{itemize}
\end{exemplebox}

\begin{propriete}[La structure « 继承 + A + 的 + B »]
Pour dire « X hérite du B de A », le chinois utilise : \cle{Sujet + 继承 + A + 的 + B}.\\
Le mot-outil \zh{的} (\emph{de}) marque la possession : \zh{父亲的头衔} = « le titre \textbf{du} père ».\\
Exemple : \zh{女儿继承母亲的房子。} (\emph{Nǚ'ér jìchéng mǔqin de fángzi.}) « La fille hérite de la maison de sa mère. »\\
Attention : en français on dit « hériter \textbf{de} quelque chose » ; en chinois, \zh{继承} est directement suivi du complément, sans mot entre les deux. Le particle d'aspect \zh{了} (\emph{le}) peut suivre le verbe pour marquer l'action accomplie (ex. \zh{继承了}).
\end{propriete}

\courssub{Le processus de succession : un organigramme}

La figure suivante résume les grandes étapes du processus coutumier observé dans de nombreuses communautés camerounaises.

\begin{popfigure}
\begin{center}
\begin{tikzpicture}[
  node distance=0.9cm,
  etape/.style={rectangle, rounded corners=4pt, draw=popblue, fill=popblueL, thick, align=center, minimum width=3.4cm, minimum height=1.1cm, font=\small},
  fleche/.style={-{Stealth[length=3mm]}, thick, draw=poporange}
]
\node[etape, align=center] (deces){Décès\\ \zh{去世} \emph{qùshì}};
\node[etape, right=of deces, align=center] (reunion){Réunion de famille\\ \zh{家人开会} \emph{jiārén kāihuì}};
\node[etape, right=of reunion, align=center] (designation){Désignation de l'héritier\\ \zh{选择继承人} \emph{xuǎnzé jìchéngrén}};
\node[etape, right=of designation, align=center] (partage){Partage des biens\\ \zh{分财产} \emph{fēn cáichǎn}};
\draw[fleche] (deces) -- (reunion);
\draw[fleche] (reunion) -- (designation);
\draw[fleche] (designation) -- (partage);
\end{tikzpicture}
\end{center}
\legende{Les quatre grandes étapes de la succession coutumière au Cameroun : du décès au partage des biens.}
\end{popfigure}

\courssub{La place du fils aîné et la diversité des pratiques}

Dans de nombreuses traditions camerounaises, le \textbf{fils aîné} (\zh{长子}) occupe une place centrale : il hérite souvent du titre du père, devient le nouveau « pilier » de la famille et a le devoir de s'occuper de sa mère et de ses cadets. Ce n'est pas seulement un privilège : c'est aussi une \textbf{responsabilité} (\zh{责任}, \emph{zérèn}).

Mais attention à ne pas généraliser : le Cameroun compte plus de deux cents groupes ethniques, et les pratiques varient fortement selon les régions et les ethnies. Dans certaines communautés, l'héritage passe par la lignée du père (système patrilinéaire) ; dans d'autres, par la lignée de la mère (système matrilinéaire) — par exemple, un neveu, fils de la sœur du défunt, peut être désigné héritier. Les biens, le titre et le nom ne suivent pas toujours le même chemin.

\begin{exemplebox}[Exprimer la diversité en chinois]
\begin{itemize}
\item \zh{每个民族的做法不一样。}\\
\emph{Měi ge mínzú de zuòfǎ bù yíyàng.}\\
« Les pratiques de chaque ethnie ne sont pas les mêmes. »
\item \zh{有的地方，侄子继承叔叔的财产。}\\
\emph{Yǒu de dìfang, zhízi jìchéng shūshu de cáichǎn.}\\
« Dans certains endroits, le neveu hérite des biens de son oncle. »
\item \zh{长子不但要继承头衔，而且要照顾家人。}\\
\emph{Zhǎngzǐ búdàn yào jìchéng tóuxián, érqiě yào zhàogù jiārén.}\\
« Le fils aîné doit non seulement hériter du titre, mais aussi prendre soin de la famille. »
\end{itemize}
\end{exemplebox}

\begin{propriete}[La structure « 不但……而且…… »]
Pour dire « non seulement... mais aussi... » : \cle{Sujet + 不但 + A + 而且 + B}.\\
Exemple : \zh{他不但继承房子，而且继承田地。} (\emph{Tā búdàn jìchéng fángzi, érqiě jìchéng tiándì.}) « Il hérite non seulement de la maison, mais aussi des champs. »\\
Piège : ne traduisez pas « mais » par \zh{但是} dans cette structure ; c'est toujours \zh{而且}.
\end{propriete}

\courssub{Le droit moderne et la coexistence des deux systèmes}

À côté de la coutume, le Cameroun dispose d'un droit moderne, hérité de l'organisation étatique. Trois outils sont essentiels :
\begin{itemize}
\item le \textbf{mariage civil} (\zh{民事结婚}, \emph{mínshì jiéhūn}), célébré à l'état civil, qui donne des droits reconnus par la loi aux époux et aux enfants ;
\item le \textbf{testament} (\zh{遗嘱}, \emph{yízhǔ}), document écrit par lequel une personne décide elle-même de la répartition de ses biens ;
\item le \textbf{notaire} (\zh{公证人}, \emph{gōngzhèngrén}), officier public qui authentifie les actes à l'étude notariale (\zh{公证处}).
\end{itemize}

La réalité camerounaise est souvent une \textbf{coexistence} des deux systèmes : une même succession peut être réglée en partie par la coutume (réunion de famille, désignation de l'héritier coutumier) et en partie par le droit moderne (tribunal, acte notarié). On parle alors de « doubles successions » ou de pluralisme juridique.

\begin{savaistu}[Le pluralisme juridique au Cameroun]
Au Cameroun, certaines successions relèvent à la fois de la coutume et du tribunal : on parle de \textbf{pluralisme juridique}. Par exemple, une famille peut désigner un héritier coutumier pour le titre et la maison familiale, tandis qu'un tribunal ou un notaire règle officiellement le partage des biens enregistrés. Ce n'est pas un désordre : c'est la rencontre de deux logiques, la tradition communautaire et le droit écrit, qui coexistent dans la vie quotidienne.
\end{savaistu}

La situation des \textbf{veuves} (\zh{寡妇}, \emph{guǎfu}) et des \textbf{filles} (\zh{女儿}) varie beaucoup selon les contextes. Dans certaines traditions, la veuve est prise en charge par la famille de son mari ; dans d'autres, elle peut rencontrer des difficultés à conserver les biens. Les filles, dans certaines coutumes, n'héritent pas des terres, alors que le droit moderne reconnaît l'égalité entre enfants. Aujourd'hui, de nombreuses familles et de nombreux tribunaux cherchent un équilibre respectueux des personnes.

\begin{exemplebox}[Exemples sur le droit moderne]
\begin{itemize}
\item \zh{他写了遗嘱，把房子给了女儿。}\\
\emph{Tā xiě le yízhǔ, bǎ fángzi gěi le nǚ'ér.}\\
« Il a rédigé un testament et a donné la maison à sa fille. »
\item \zh{他们去公证处办手续。}\\
\emph{Tāmen qù gōngzhèngchù bàn shǒuxù.}\\
« Ils vont chez le notaire pour accomplir les formalités. »
\item \zh{这个家庭同时用传统和法律。}\\
\emph{Zhè ge jiātíng tóngshí yòng chuántǒng hé fǎlǜ.}\\
« Cette famille utilise à la fois la tradition et la loi. »
\end{itemize}
\end{exemplebox}

\courssub{Tableau récapitulatif : les deux systèmes en regard}

\begin{center}
\begin{tabularx}{\linewidth}{|X|X|X|}
\hline
\rowcolor{popblueL} Aspect & Coutume (传统) & Droit moderne (现代法律) \\
\hline
Qui décide ? & Famille élargie, aînés, notables (\zh{长辈、家人}) & Tribunaux, notaire (\zh{公证人}) \\
\hline
Quel document ? & Parole, réunion de famille (\zh{开会}) & Testament, actes (\zh{遗嘱}) \\
\hline
Qui hérite souvent ? & Fils aîné ou héritier désigné (\zh{长子}) & Conjoint et enfants selon la loi \\
\hline
Qu'est-ce qui est transmis ? & Nom, titre, maison, champs, bétail & Biens enregistrés, comptes, maisons \\
\hline
\end{tabularx}
\end{center}

\courssub{Exercice résolu et piège à éviter}

\begin{exoresolu}[Traduisez en chinois]
Énoncé : Traduisez la phrase suivante en chinois (caractères + pinyin) : « Après le décès du père, la famille se réunit et choisit le fils aîné comme héritier. »
\tcblower
Solution détaillée, étape par étape :
\begin{enumerate}
\item « Après le décès du père » : on utilise la structure \zh{……以后} (« après que... ») : \zh{父亲去世以后} (\emph{fùqin qùshì yǐhòu}). Le complément de temps se place \textbf{avant} le verbe principal, contrairement au français.
\item « la famille se réunit » : \zh{家人开会} (\emph{jiārén kāihuì}).
\item « choisit le fils aîné comme héritier » : \zh{选择长子当继承人} (\emph{xuǎnzé zhǎngzǐ dāng jìchéngrén}). Le verbe \zh{当} (\emph{dāng}) signifie ici « comme, en tant que ».
\end{enumerate}
Phrase complète : \zh{父亲去世以后，家人开会，选择长子当继承人。}\\
\emph{Fùqin qùshì yǐhòu, jiārén kāihuì, xuǎnzé zhǎngzǐ dāng jìchéngrén.}
\end{exoresolu}

\begin{attention}[Erreurs fréquentes des francophones]
\begin{itemize}
\item \textbf{Ordre des mots} : ne dites pas \zh{家人开会以后父亲去世} en voulant dire « après le décès du père, la famille se réunit ». En chinois, l'action qui précède vient d'abord : \zh{父亲去世以后，家人开会}. La chronologie chinoise suit l'ordre réel des événements.
\item \textbf{Fausse préposition} : ne traduisez pas « hériter de » par \zh{继承} + \zh{从} ou \zh{的}. On dit directement \zh{继承财产} (« hériter des biens »).
\item \textbf{Ton de 长} : dans \zh{长子} (fils aîné), prononcez \emph{zhǎngzǐ} ; \emph{chángzi} n'existe pas avec ce sens et ferait sourire.
\item \textbf{Respect du vocabulaire} : pour parler d'un décès, préférez toujours \zh{去世} à une formulation trop directe ; c'est une question de politesse socioculturelle, au Cameroun comme en Chine.
\end{itemize}
\end{attention}

\begin{exercice}[À vous de jouer]
\begin{enumerate}
\item Traduisez en français : \zh{在很多传统里，长子继承父亲的头衔和房子。}
\item Traduisez en chinois : « Les aînés de la famille décident qui hérite des champs. »
\item Complétez avec \zh{继承}, \zh{继承人} ou \zh{遗嘱} : \zh{他写了……，女儿是……，她……母亲的财产。}
\item Question de réflexion (pour préparer l'exposé de la section 6) : selon vous, quels sont les avantages de la réunion de famille coutumière ? Et ceux du testament moderne ? Notez deux idées en français, puis essayez d'en formuler une en chinois avec \zh{我觉得……} (« je pense que... »).
\end{enumerate}
\end{exercice}

\begin{corrige}[Exercice « À vous de jouer »]
\begin{enumerate}
\item « Dans de nombreuses traditions, le fils aîné hérite du titre et de la maison du père. »
\item \zh{家里的长辈决定谁继承田地。} (\emph{Jiā lǐ de zhǎngbèi juédìng shéi jìchéng tiándì.})
\item \zh{他写了遗嘱，女儿是继承人，她继承母亲的财产。}
\item Réponse libre. Modèle de formulation : \zh{我觉得家人开会很好，因为大家一起商量。} (\emph{Wǒ juéde jiārén kāihuì hěn hǎo, yīnwèi dàjiā yìqǐ shāngliang.}) « Je pense que la réunion de famille est une bonne chose, parce que tout le monde délibère ensemble. »
\end{enumerate}
\end{corrige}

\begin{aretenir}[L'essentiel de la section]
\begin{itemize}
\item Au Cameroun, la succession coutumière repose sur la \textbf{famille élargie} : après le décès (\zh{去世}), la famille se réunit (\zh{开会}) et désigne l'héritier (\zh{继承人}).
\item On transmet le nom, le titre (\zh{头衔}) et les biens (\zh{财产} : maison, champs, bétail).
\item Le fils aîné (\zh{长子}) joue souvent un rôle central, mais les pratiques varient selon les régions et les ethnies (systèmes patrilinéaires et matrilinéaires).
\item Le droit moderne (mariage civil, testament \zh{遗嘱}, notaire) coexiste avec la coutume : c'est le pluralisme juridique.
\item Structures clés : \zh{继承 + A + 的 + B} ; \zh{……以后} ; \zh{不但……而且……}.
\end{itemize}
\end{aretenir}

Nous disposons maintenant d'une vision claire et factuelle des pratiques camerounaises. La section suivante appliquera la même démarche à la Chine, avant que nous ne passions à la comparaison proprement dite et à la préparation de votre mini-exposé.

\coursec{Rappel ciblé et mise en route : l'hérédité au Cameroun (leçon 4)}

\courssub{Contexte : pluralisme juridique et coutumes}

Avant d'étudier le système chinois, rappelons les points clés de la leçon 4 sur le Cameroun. Le Cameroun se caractérise par un \cle{pluralisme juridique} : le droit moderne (Code civil, Code de la famille) coexiste avec les \cle{droits coutumiers} qui varient selon les groupes ethniques et les régions. Cette dualité n'existe pas en Chine, où la loi nationale est unique.

\begin{textebox}[Rappel : l'hérédité au Cameroun — faits principaux]
Au Cameroun, la succession dépend souvent du \cle{régime matrimonial} et de la \cle{coutume} du défunt.\\
\emph{Au Kámélún, jìchéng yībān yīlài de qīnqī guānxì hé gùxí.}\\
« Au Cameroun, l'héritage dépend souvent des liens de parenté et des coutumes. »\\[0.4em]
Dans de nombreuses traditions (Bamiléké, Bassa, Douala, etc.), la famille élargie joue un rôle central. Le \cle{chef de famille} ou le \cle{notable} supervise le partage.\\
\emph{Zài xǔduō chuántǒng zhōng, jiātíng kuòdà juéjué zhōngyāng zuòyòng. Jiāzhǎng huò zhěnmíng jiāndū fēnpèi.}\\
« Dans de nombreuses traditions, la famille élargie joue un rôle central. Le chef de famille ou le notable supervise le partage. »\\[0.4em]
La transmission est souvent \cle{patrilinéaire} (biens aux fils, surtout l'aîné), mais des coutumes \cle{matrilinéaires} existent (héritage par les neveux du côté maternel).\\
\emph{Chuánchéng yībān shì pǔxì (gěi érzi, tèbié dà érzi), kěshì yě yǒu mǔxì gùxí (wàishēng zǐnǚ jìchéng).}\\
« La transmission est souvent patrilinéaire (aux fils, surtout l'aîné), mais il existe des coutumes matrilinéaires (héritage par les neveux maternels). »\\[0.4em]
Le droit moderne protège le \cle{conjoint survivant} et les \cle{enfants} (filles et garçons), mais l'application reste complexe entre loi et coutume.\\
\emph{Xiàndài fǎlǜ bǎohù cúnshēng pèi'ǒu hé zǐnǚ (nǚ'ér hé érzi), dàn shíshī fùzá.}\\
« La loi moderne protège le conjoint survivant et les enfants (filles et garçons), mais l'application reste complexe. »
\end{textebox}

\begin{center}
\begin{tabularx}{\linewidth}{|X|X|X|X|}
\hline
\rowcolor{popblueL} Caractères & Pinyin & Nature & Traduction \\
\hline
习俗 & xísú & n. & coutume, usage \\
\hline
家族 & jiāzú & n. & clan, famille élargie \\
\hline
族长 & zúzhǎng & n. & chef de clan / de famille \\
\hline
父系 / 母系 & fùxì / mǔxì & adj. & patrilinéaire / matrilinéaire \\
\hline
继承人 & jìchéngrén & n. & héritier \\
\hline
遗产 & yíchǎn & n. & héritage, succession \\
\hline
分配 & fēnpèi & v. & répartir, partager \\
\hline
配偶 & pèi'ǒu & n. & conjoint(e) \\
\hline
\end{tabularx}
\end{center}

\begin{experience}[Mise en route orale : « Chez moi, comment ça se passe ? »]
En binôme, chaque élève décrit en 2 phrases (chinois + français) une coutume d'héritage connue dans sa famille ou sa région (ex. : \zh{我们家按父系继承。} \emph{Wǒmen jiā àn fùxì jìchéng.} / \zh{外婆的财产给舅舅。} \emph{Wàipó de cáichǎn gěi jiùjiu.}). L'enseignant note au tableau les mots-clés : \zh{父系}、\zh{母系}、\zh{长子}、\zh{女儿也有份}。 Cela prépare le vocabulaire comparatif.
\end{experience}

\coursec{L'hérédité au Cameroun : faits structurés (synthèse pour la comparaison)}

Pour comparer rigoureusement, fixons trois piliers du système camerounais :

\begin{propriete}[Trois piliers de l'hérédité camerounaise]
\begin{enumerate}
\item \textbf{Pluralisme des sources} : Loi moderne (nationale) \textbf{ET} droits coutumiers (locaux). Un même décès peut relever des deux.
\item \textbf{Place de la famille élargie} : Le clan (\zh{家族} \emph{jiāzú}) et son chef (\zh{族长} \emph{zúzhǎng}) interviennent dans le partage, au-delà du cercle nucléaire (père/mère/enfants).
\item \textbf{Diversité des règles de dévolution} : Majoritairement \textbf{patrilinéaire} (\zh{父系} \emph{fùxì} : fils aîné privilégié), parfois \textbf{matrilinéaire} (\zh{母系} \emph{mǔxì} : neveux maternels). La loi moderne impose l'égalité filles/garçons, mais la coutume résiste souvent.
\end{enumerate}
\end{propriete}

\begin{attention}[Ne pas confondre : « Coutume » et « Habitude »]
En français, « coutume » (juridique, contraignante) $\neq$ « habitude » (comportement répété). En chinois : \zh{习俗} (\emph{xísú}, coutume sociale/juridique) vs \zh{习惯} (\emph{xíguàn}, habitude personnelle). Pour l'héritage, on utilise \zh{习俗} ou \zh{ customary law} (\zh{习惯法} \emph{xíguànfǎ}).
\end{attention}

\coursec{L'hérédité en Chine : faits}

Après avoir observé les pratiques d'héritage au Cameroun — pluralité des droits, rôle de la famille élargie, place du chef de famille — tournons-nous maintenant vers la Chine. Le contraste est frappant : là où le Cameroun connaît une coexistence de plusieurs systèmes (droit moderne, droits coutumiers), la Chine applique un \cle{cadre national unique}, une loi écrite qui s'impose à tous. Mais derrière cette loi moderne se cache une tradition millénaire : le respect des ancêtres, la place du fils aîné, la transmission du nom de famille. Observons d'abord un texte, puis dégageons les faits.

\courssub{Observation : un texte sur la succession en Chine}

\begin{textebox}[中国的继承制度 — Le système successoral chinois]
中国有一部全国统一的法律，叫《继承法》。\\
\emph{Zhōngguó yǒu yī bù quánguó tǒngyī de fǎlǜ, jiào 《Jìchéngfǎ》.}\\
« La Chine possède une loi nationale unifiée, appelée la Loi sur la succession. »\\[0.4em]
这部法律规定：人去世以后，他的财产由谁继承。第一顺序的继承人是配偶、子女和父母。\\
\emph{Zhè bù fǎlǜ guīdìng: rén qùshì yǐhòu, tā de cáichǎn yóu shéi jìchéng. Dì-yī shùnxù de jìchéngrén shì pèi'ǒu, zǐnǚ hé fùmǔ.}\\
« Cette loi stipule : après le décès d'une personne, qui hérite de ses biens. Les héritiers de premier rang sont le conjoint, les enfants et les parents. »\\[0.4em]
中国的继承法规定，儿子和女儿有一样的继承权。这一点和古代很不一样：以前，财产主要传给儿子，特别是大儿子。\\
\emph{Zhōngguó de jìchéngfǎ guīdìng, érzi hé nǚ'ér yǒu yīyàng de jìchéngquán. Zhè yī diǎn hé gǔdài hěn bù yīyàng: yǐqián, cáichǎn zhǔyào chuán gěi érzi, tèbié shì dà érzi.}\\
« La loi chinoise sur la succession stipule que les fils et les filles ont les mêmes droits d'héritage. Sur ce point, c'est très différent de l'Antiquité : autrefois, les biens étaient transmis principalement aux fils, surtout au fils aîné. »\\[0.4em]
现在，越来越多的城里人写遗嘱。遗嘱是一个人去世以前写的文件，说明他想把财产给谁。有了遗嘱，就要按照遗嘱来分配财产。\\
\emph{Xiànzài, yuè lái yuè duō de chénglǐ rén xiě yízhǔ. Yízhǔ shì yī gè rén qùshì yǐqián xiě de wénjiàn, shuōmíng tā xiǎng bǎ cáichǎn gěi shéi. Yǒu le yízhǔ, jiù yào ànzhào yízhǔ lái fēnpèi cáichǎn.}\\
« Aujourd'hui, de plus en plus de citadins rédigent un testament. Le testament est un document écrit avant le décès, qui précise à qui la personne veut laisser ses biens. Quand il y a un testament, on doit répartir les biens selon le testament. »\\[0.4em]
在中国，孩子一般跟父亲姓。不过现在，也有一些孩子跟母亲姓。人去世以后，家人要守孝，穿深色的衣服。每年清明节，全家人一起去扫墓，纪念去世的亲人。对中国人来说，继承不只是财产，也是对祖先的记忆。\\
\emph{Zài Zhōngguó, háizi yībān gēn fùqīn xìng. Bùguò xiànzài, yě yǒu yīxiē háizi gēn mǔqīn xìng. Rén qùshì yǐhòu, jiārén yào shǒuxiào, chuān shēnsè de yīfu. Měi nián Qīngmíng jié, quánjiārén yīqǐ qù sǎomù, jìniàn qùshì de qīnrén. Duì Zhōngguórén lái shuō, jìchéng bù zhǐ shì cáichǎn, yě shì duì zǔxiān de jìyì.}\\
« En Chine, les enfants portent en général le nom du père. Mais aujourd'hui, certains enfants portent aussi le nom de la mère. Après un décès, la famille observe le deuil et porte des vêtements sombres. Chaque année, à la fête de Qingming, toute la famille va ensemble nettoyer les tombes et commémorer les proches disparus. Pour les Chinois, hériter, ce n'est pas seulement une question de biens, c'est aussi la mémoire des ancêtres. »
\end{textebox}

\textbf{Glossaire du texte}\par

\begin{center}
\begin{tabularx}{\linewidth}{|X|X|X|X|}
\hline
\rowcolor{popblueL} Caractères & Pinyin & Nature & Traduction \\
\hline
继承法 & jìchéngfǎ & n. & loi sur la succession \\
\hline
统一 & tǒngyī & adj./v. & unifié ; unifier \\
\hline
法律 & fǎlǜ & n. & loi, législation \\
\hline
规定 & guīdìng & v./n. & stipuler ; disposition \\
\hline
去世 & qùshì & v. & décéder (registre soutenu) \\
\hline
财产 & cáichǎn & n. & biens, patrimoine \\
\hline
顺序 & shùnxù & n. & ordre, rang \\
\hline
继承人 & jìchéngrén & n. & héritier \\
\hline
继承权 & jìchéngquán & n. & droit d'héritage \\
\hline
遗嘱 & yízhǔ & n. & testament \\
\hline
文件 & wénjiàn & n. & document \\
\hline
分配 & fēnpèi & v. & répartir, distribuer \\
\hline
跟……姓 & gēn... xìng & expr. & porter le nom de... \\
\hline
守孝 & shǒuxiào & v. & observer le deuil \\
\hline
扫墓 & sǎomù & v. & nettoyer les tombes \\
\hline
祖先 & zǔxiān & n. & ancêtres \\
\hline
\end{tabularx}
\end{center}

\textbf{Questions de compréhension}\par
\begin{enumerate}
\item Qui sont les héritiers de premier rang en Chine ? (\emph{Utilisez} 第一顺序的继承人是……)
\item Les filles ont-elles les mêmes droits que les fils ? Répondez par une phrase complète.
\item Que se passe-t-il quand le défunt a écrit un testament ?
\item Que font les familles chinoises à la fête de Qingming ?
\end{enumerate}

\courssub{Un cadre national unique : la loi sur la succession}

\begin{definition}[La loi sur la succession 继承法]
En Chine, la succession est réglée par un texte de loi national, la \cle{继承法} (\emph{jìchéngfǎ}), aujourd'hui intégrée dans le Code civil chinois (\zh{民法典} \emph{Mínfǎdiǎn}). Ce cadre est \textbf{unique} : il s'applique de la même façon dans toutes les régions, en ville comme à la campagne, sans droits coutumiers parallèles. C'est une différence majeure avec le Cameroun, où droit moderne et coutumes coexistent.
\end{definition}

\begin{exemplebox}[Observer la structure « 由 + personne + verbe »]
\zh{他的财产由配偶、子女和父母继承。}\\
\emph{Tā de cáichǎn yóu pèi'ǒu, zǐnǚ hé fùmǔ jìchéng.}\\
« Ses biens sont hérités par le conjoint, les enfants et les parents. »
\end{exemplebox}

\begin{propriete}[La structure passive avec 由]
Structure : \textbf{Nom + 由 (yóu) + agent + verbe}. Elle signifie « être ... par ... » et insiste sur \emph{qui fait l'action} (responsabilité de l'agent).\\
\zh{遗产由第一顺序的继承人继承。} \emph{Yíchǎn yóu dì-yī shùnxù de jìchéngrén jìchéng.} « L'héritage est transmis par les héritiers de premier rang. »\\
\zh{这个决定由法院做出。} \emph{Zhège juédìng yóu fǎyuàn zuòchū.} « Cette décision est prise par le tribunal. »\\
Attention : 由 n'est pas interchangeable avec 被 (\emph{bèi}) : 被 marque le passif « subi » (souvent négatif), 由 marque la responsabilité de l'agent (souvent neutre ou formel).
\end{propriete}

\courssub{L'ordre des héritiers légaux}

La loi chinoise classe les héritiers par \cle{rangs} (顺序, \emph{shùnxù}). Les héritiers du \textbf{premier rang} (第一顺序) héritent en priorité et se partagent les biens à parts égales ; ceux du deuxième rang n'héritent que s'il n'y a \textbf{personne} au premier rang.

\begin{center}
\begin{tabularx}{\linewidth}{|X|X|X|}
\hline
\rowcolor{popblueL} Rang & Personnes (caractères + pinyin) & Traduction \\
\hline
Premier rang (第一顺序) & 配偶、子女、父母 \emph{pèi'ǒu, zǐnǚ, fùmǔ} & conjoint, enfants, parents \\
\hline
Deuxième rang (第二顺序) & 兄弟姐妹、祖父母 \emph{xiōngdì jiěmèi, zǔfùmǔ} & frères et sœurs, grands-parents \\
\hline
\end{tabularx}
\end{center}

\begin{aretenir}[L'égalité filles / fils — point d'examen]
Point essentiel : la loi chinoise garantit l'\cle{égalité} entre fils et filles (\zh{儿子和女儿有一样的继承权}). C'est une règle moderne qui rompt avec la tradition ancienne, où les biens passaient surtout aux fils. La formule à retenir :\\
\zh{儿子和女儿有一样的继承权。}\\
\emph{Érzi hé nǚ'ér yǒu yīyàng de jìchéngquán.}\\
« Les fils et les filles ont les mêmes droits d'héritage. »
\end{aretenir}

\begin{popfigure}
\begin{tikzpicture}[scale=1, every node/.style={font=\small}]
% Testament top
\fill[popgold!80] (0,1.2) -- (2.5,1.2) -- (1.25,2.2) -- cycle;
\node[align=center, text=popdark] at (1.25,1.6) {\zh{遗嘱}\\ testament};
% 1st rank
\fill[popblue!40] (-2,0) rectangle (4.5,1.1);
\node[align=center, text=popdark] at (1.25,0.55) {\zh{配偶、子女、父母}\\ conjoint, enfants, parents (1er rang)};
% 2nd rank
\fill[popgreen!35] (-3,-1.5) rectangle (5.5,-0.1);
\node[align=center, text=popdark] at (1.25,-0.8) {\zh{兄弟姐妹、祖父母等}\\ frères/sœurs, grands-parents (2e rang)};
% Arrows
\draw[popdark, thick, ->] (1.25,2.2) -- (1.25,1.1);
\draw[popdark, thick, ->] (1.25,0) -- (1.25,-0.1);
% Priority text
\node[align=center, text=popmuted] at (1.25,-2) {Priorité : Testament $\rightarrow$ 1er rang $\rightarrow$ 2e rang};
\end{tikzpicture}
\legende{Pyramide de la transmission successorale en Chine : le testament prime, puis les héritiers de premier rang, puis les autres parents.}
\end{popfigure}

\courssub{Le testament écrit : une pratique en développement}

\begin{definition}[遗嘱 yízhǔ]
Le \cle{testament} : document écrit dans lequel une personne exprime ses \textbf{dernières volontés}, c'est-à-dire à qui elle veut laisser ses biens après sa mort. Caractères : 遗 (laisser après soi, clé 辶 « mouvement ») + 嘱 (ordonner, recommander, clé 口 « bouche »). Autrefois rare en Chine (on évitait de parler de la mort), le testament devient \textbf{de plus en plus courant en ville}, car il évite les conflits familiaux.
\end{definition}

\begin{exemplebox}[Exemples traduits]
\zh{他写了遗嘱，把房子给了两个孩子。}\\
\emph{Tā xiě le yízhǔ, bǎ fángzi gěi le liǎng gè háizi.}\\
« Il a écrit un testament et a laissé la maison à ses deux enfants. »\\[0.4em]
\zh{有了遗嘱，就要按照遗嘱分配财产。}\\
\emph{Yǒu le yízhǔ, jiù yào ànzhào yízhǔ fēnpèi cáichǎn.}\\
« Quand il y a un testament, il faut répartir les biens selon le testament. »\\[0.4em]
\zh{现在，很多城里人愿意写遗嘱。}\\
\emph{Xiànzài, hěn duō chénglǐ rén yuànyì xiě yízhǔ.}\\
« Aujourd'hui, beaucoup de citadins acceptent de rédiger un testament. »
\end{exemplebox}

\begin{propriete}[La construction 把 (bǎ) — transmission et disposition]
Structure : \textbf{Sujet + 把 + objet + verbe + complément (souvent 给 + destinataire)}. Elle déplace l'objet avant le verbe quand l'action « dispose » de l'objet (donner, laisser, confier). Très utile pour parler de transmission successorale :\\
\zh{爷爷把土地给了大儿子。} \emph{Yéye bǎ tǔdì gěi le dà érzi.} « Le grand-père a donné la terre au fils aîné. »\\
\zh{她把遗产留给了女儿。} \emph{Tā bǎ yíchǎn liú gěi le nǚ'ér.} « Elle a laissé l'héritage à sa fille. »\\
En français, on dit simplement « donner la terre au fils » ; en chinois, \textbf{把 + objet + 给} est la tournure la plus naturelle pour insister sur le sort de l'objet.
\end{propriete}

\courssub{Tradition et continuité : fils aîné, nom de famille, culte des ancêtres}

Derrière la loi moderne, la tradition reste vivante. Trois faits à connaître :

\begin{itemize}
\item \textbf{Le fils aîné dans la Chine ancienne} : autrefois, le fils aîné (大儿子, \emph{dà érzi}) occupait une place privilégiée ; il recevait la part principale et surtout la \emph{responsabilité du culte des ancêtres} (\zh{祭祖} \emph{jìzǔ}). Aujourd'hui, la loi efface ce privilège juridique, mais l'idée d'une responsabilité particulière de l'aîné subsiste dans les mentalités, surtout à la campagne.
\item \textbf{Le nom de famille} (姓, \emph{xìng}) : les enfants portent traditionnellement le nom du \textbf{père} (\zh{跟父亲姓} \emph{gēn fùqīn xìng}). Évolution récente : la loi autorise désormais le nom de la mère (\zh{跟母亲姓} \emph{gēn mǔqīn xìng}), et quelques familles urbaines font ce choix. \zh{孩子一般跟父亲姓，不过现在也有孩子跟母亲姓。} \emph{Háizi yībān gēn fùqīn xìng, bùguò xiànzài yě yǒu háizi gēn mǔqīn xìng.} « Les enfants portent en général le nom du père, mais aujourd'hui certains portent celui de la mère. »
\item \textbf{Le respect des anciens et les rites} : après un décès, la famille observe le deuil (\zh{守孝} \emph{shǒuxiào}) ; chaque année, à la fête de Qingming (\zh{清明节} \emph{Qīngmíng jié}), on entretient les tombes (\zh{扫墓} \emph{sǎomù}). Hériter, c'est aussi \emph{se souvenir}.
\end{itemize}

\begin{savaistu}[清明节 : l'héritage, c'est aussi la mémoire]
À la fête de \cle{Qingming} (清明节, \emph{Qīngmíng jié}, début avril), des millions de Chinois se rendent sur les tombes de leurs ancêtres : ils les nettoient, déposent des fleurs et des offrandes, et saluent les défunts. Cette fête montre que, dans la culture chinoise, la transmission ne concerne pas seulement les biens matériels : elle inclut la \textbf{mémoire} et le \textbf{respect des ancêtres} (祖先, \emph{zǔxiān}). On peut comparer cette pratique, avec respect, aux rites funéraires et aux cérémonies de mémoire des familles camerounaises (veillées, deuil, hommage aux ancêtres).
\end{savaistu}

\coursec{Vocabulaire du thème : synthèse bilangue (Cameroun — Chine)}

Ce tableau regroupe le lexique essentiel pour le mini-exposé comparatif.

\begin{center}
\begin{tabularx}{\linewidth}{|X|X|X|X|}
\hline
\rowcolor{popblueL} Caractères & Pinyin & Nature & Traduction \\
\hline
继承 & jìchéng & v. & hériter, succession \\
\hline
遗产 & yíchǎn & n. & héritage, patrimoine \\
\hline
遗嘱 & yízhǔ & n. & testament \\
\hline
法律 & fǎlǜ & n. & loi \\
\hline
习俗 / 习惯法 & xísú / xíguànfǎ & n. & coutume / droit coutumier \\
\hline
统一 & tǒngyī & adj. & unifié, unique \\
\hline
配偶 & pèi'ǒu & n. & conjoint(e) \\
\hline
子女 & zǐnǚ & n. & enfants (fils et filles) \\
\hline
父母 & fùmǔ & n. & parents \\
\hline
兄弟姐妹 & xiōngdì jiěmèi & n. & frères et sœurs \\
\hline
长子 / 大儿子 & zhǎngzǐ / dà érzi & n. & fils aîné \\
\hline
继承权 & jìchéngquán & n. & droit d'héritage \\
\hline
平等 & píngděng & adj. & égalité, égal \\
\hline
一样 & yīyàng & adj. & pareil, identique \\
\hline
姓 & xìng & n./v. & nom de famille ; s'appeler \\
\hline
跟……姓 & gēn... xìng & expr. & porter le nom de... \\
\hline
父系 / 母系 & fùxì / mǔxì & adj. & patrilinéaire / matrilinéaire \\
\hline
家族 & jiāzú & n. & clan, famille élargie \\
\hline
族长 & zúzhǎng & n. & chef de clan/famille \\
\hline
守孝 & shǒuxiào & v. & observer le deuil \\
\hline
清明节 & Qīngmíng jié & n. propre & fête de Qingming \\
\hline
扫墓 & sǎomù & v. & nettoyer les tombes \\
\hline
祭祖 & jìzǔ & v. & rendre culte aux ancêtres \\
\hline
祖先 & zǔxiān & n. & ancêtres \\
\hline
纪念 & jìniàn & v. & commémorer \\
\hline
分配 & fēnpèi & v. & répartir, partager \\
\hline
规定 & guīdìng & v./n. & stipuler ; disposition \\
\hline
去世 & qùshì & v. & décéder (soutenu) \\
\hline
\end{tabularx}
\end{center}

\coursec{Comparer et donner son avis en chinois}

Pour réussir le mini-exposé, il faut maîtriser les structures de comparaison et de nuance.

\courssub{Structures de comparaison}

\begin{propriete}[Comparer avec 比 (bǐ) — supériorité / différence]
Structure : \textbf{A + 比 + B + Adjectif / Verbe d'état}. Pas de « de » après l'adjectif.
\zh{中国的法律比喀麦隆的习俗统一。}\\
\emph{Zhōngguó de fǎlǜ bǐ Kāmàilóng de xísú tǒngyī.}\\
« La loi chinoise est plus unifiée que les coutumes camerounaises. »\\
\zh{现在的继承法比古代平等。}\\
\emph{Xiànzài de jìchéngfǎ bǐ gǔdài píngděng.}\\
« La loi successorale actuelle est plus égalitaire qu'autrefois. »
\end{propriete}

\begin{propriete}[Égalité et similitude : 跟……一样 / 跟……不一样]
Structure : \textbf{A + 跟 + B + 一样 / 不一样 (+ Adjectif / Verbe)}.
\zh{儿子跟女儿一样，都有继承权。}\\
\emph{Érzi gēn nǚ'ér yīyàng, dōu yǒu jìchéngquán.}\\
« Le fils, comme la fille, a le droit d'héritage. »\\
\zh{中国的制度跟喀麦隆的不一样。}\\
\emph{Zhōngguó de zhìdù gēn Kāmàilóng de bù yīyàng.}\\
« Le système chinois n'est pas le même que le camerounais. »
\end{propriete}

\begin{propriete}[Contraste et concession : 虽然……，但是…… / 可是……]
Structure : \textbf{虽然 + Situation 1, 但是/可是 + Situation 2 (contraire)}.
\zh{虽然法律规定平等，但是传统观念还在。}\\
\emph{Suīrán fǎlǜ guīdìng píngděng, dànshì chuántǒng guānniàn hái zài.}\\
« Bien que la loi prévoie l'égalité, les mentalités traditionnelles persistent. »\\
\zh{喀麦隆有习惯法，可是中国只有统一法律。}\\
\emph{Kāmàilóng yǒu xíguànfǎ, kěshì Zhōngguó zhǐyǒu tǒngyī fǎlǜ.}\\
« Le Cameroun a le droit coutumier, mais la Chine n'a qu'une loi unifiée. »
\end{propriete}

\begin{propriete}[Exprimer son avis : 我认为 / 我觉得 / 依我看]
\zh{我认为中国的制度更清晰。} \emph{Wǒ rènwéi Zhōngguó de zhìdù gèng qīngxī.} « Je pense que le système chinois est plus clair. »\\
\zh{我觉得喀麦隆的习俗尊重家族。} \emph{Wǒ juéde Kāmàilóng de xísú zūnzhòng jiāzú.} « Je trouve que les coutumes camerounaises respectent le clan. »\\
\zh{依我看，两种制度各有优缺点。} \emph{Yī wǒ kàn, liǎng zhǒng zhìdù gè yǒu yōu quēdiǎn.} « À mon avis, les deux systèmes ont chacun leurs avantages et inconvénients. »
\end{propriete}

\begin{methode}[Réussir une phrase de comparaison structurée]
\begin{enumerate}
\item Identifier les deux termes comparés (A et B).
\item Choisir le critère (adjectif : \zh{统一}、\zh{平等}、\zh{复杂}、\zh{传统}).
\item Choisir la structure : \zh{比} (différence nette) OU \zh{跟...一样/不一样} (similitude/différence qualitative) OU \zh{虽然...但是...} (nuance).
\item Construire la phrase : \zh{A 比 B 复杂} / \zh{A 跟 B 不一样} / \zh{虽然 A..., 但是 B...}.
\item Vérifier l'ordre des mots (pas de « de » après l'adjectif avec \zh{比}).
\end{enumerate}
\end{methode}

\coursec{Questions de réflexion et préparation du mini-exposé}

\courssub{Questions de réflexion (écrit / oral)}

Répondez en chinois (phrases complètes) ou en français structuré. Ces questions guident votre exposé.

\begin{enumerate}
\item \textbf{Cadre légal} : \zh{中国和喀麦隆的继承法律有什么不同？} (Quelle différence entre les cadres juridiques chinois et camerounais ?) — Mots-clés : \zh{统一法律} vs \zh{习惯法 / 现代法律}。
\item \textbf{Égalité} : \zh{中国法律如何保障男女平等？喀麦隆呢？} (Comment la loi chinoise garantit-elle l'égalité ? Et le Cameroun ?) — Mots-clés : \zh{儿子和女儿一样}、\zh{现代法律保护}、\zh{习俗有时不平等}。
\item \textbf{Famille élargie vs nucléaire} : \zh{中国谁来分配遗产？喀麦隆呢？} (Qui partage l'héritage en Chine ? Au Cameroun ?) — Mots-clés : \zh{法院 / 遗嘱 / 继承人自己} vs \zh{族长 / 家族会议 / notable}。
\item \textbf{Tradition vivante} : \zh{中国人怎么纪念祖先？喀麦隆人呢？} (Comment les Chinois honorent-ils les ancêtres ? Les Camerounais ?) — Mots-clés : \zh{清明节、扫墓、祭祖} vs \zh{veillée、deuil、cérémonies coutumières}。
\item \textbf{Nom de famille} : \zh{中国孩子跟谁姓？喀麦隆呢？有变化吗？} (Nom de qui portent les enfants ? Changement ?) — Mots-clés : \zh{跟父亲姓}、\zh{现在也可以跟母亲姓}、\zh{喀麦隆一般跟父亲姓}。
\end{enumerate}

\begin{experience}[Atelier oral : « Débat croisé Cameroun — Chine »]
\begin{enumerate}
\item \textbf{Préparation (10 min)} : En groupes de 4. Chaque groupe tire un thème (Cadre légal / Égalité / Rôle de la famille / Rites ancestraux / Nom de famille).
\item \textbf{Production (15 min)} : Le groupe prépare un mini-exposé de 2 minutes \textbf{en chinois} (chacun parle 30 sec). Structure imposée :
      \begin{itemize}
      \item Introduction : \zh{我们组比较……} (\emph{Wǒmen zǔ bǐjiào...})
      \item Fait Chine : \zh{在中国，……} (\emph{Zài Zhōngguó...})
      \item Fait Cameroun : \zh{在喀麦隆，……} (\emph{Zài Kāmàilóng...})
      \item Comparaison : \zh{中国……，但是喀麦隆…… / 两者都……} (\emph{Zhōngguó..., dànshì Kāmàilóng... / Liǎngzhě dōu...})
      \item Avis personnel : \zh{我认为…… / 我觉得……} (\emph{Wǒ rènwéi... / Wǒ juéde...})
      \end{itemize}
\item \textbf{Présentation (10 min)} : Chaque groupe présente. La classe note 1 similitude et 1 différence entendues.
\item \textbf{Retour (5 min)} : Correction collective de la prononciation (tons sur \zh{继承}、\zh{习俗}、\zh{平等}、\zh{统一}) et de la structure \zh{比} / \zh{跟...一样}.
\end{enumerate}
\end{experience}

\begin{methode}[Rédiger le mini-exposé écrit (devoir ou évaluation)]
\begin{enumerate}
\item \textbf{Plan} : Rédigez 5 paragraphes courts (1 par question de réflexion).
\item \textbf{Vocabulaire} : Utilisez au moins 10 mots du tableau « Vocabulaire du thème ».
\item \textbf{Grammaire} : Incluez au moins : 1 phrase avec \zh{比}, 1 phrase avec \zh{跟...一样/不一样}, 1 phrase avec \zh{虽然...但是...}, 1 phrase avec \zh{把} (transmission).
\item \textbf{Caractères} : Écrivez en caractères (pas de pinyin seul). Vérifiez l'ordre des traits pour : \zh{继}、\zh{承}、\zh{遗}、\zh{嘱}、\zh{俗}、\zh{族}、\zh{祭}、\zh{祖}。
\item \textbf{Relecture} : Vérifiez les tons (dictionnaire), la ponctuation chinoise (， 。) et l'absence d'espaces entre caractères.
\end{enumerate}
\end{methode}

\coursec{Tableau récapitulatif comparatif (outil de révision)}

\begin{center}
\begin{tabularx}{\linewidth}{|X|X|X|}
\hline
\rowcolor{popblueL} Critère & Chine (中国) & Cameroun (喀麦隆) \\
\hline
Source du droit & \zh{统一法律 (民法典/继承法)} \emph{tǒngyī fǎlǜ} & \zh{现代法律 + 习惯法} \emph{xiàndài fǎlǜ + xíguànfǎ} \\
\hline
Héritiers prioritaires & \zh{配偶、子女、父母 (第一顺序)} \emph{pèi'ǒu, zǐnǚ, fùmǔ} & \zh{配偶、子女 (loi) ; famille élargie (coutume)} \emph{pèi'ǒu, zǐnǚ ; jiāzú} \\
\hline
Égalité filles / fils & \zh{法律规定完全平等} \emph{fǎlǜ guīdìng wánquán píngděng} & \zh{Loi moderne = égalité ; Coutume = souvent inégale} \emph{xiàndài fǎlǜ = píngděng; xísú = bù píngděng} \\
\hline
Rôle fils aîné & \zh{古代有特权，今无法律特权} \emph{gǔdài yǒu tequán, jīn wú fǎlǜ tequán} & \zh{Souvent privilégié (part principale, culte)} \emph{yībān tèquán (zhǔyào fēn, jìzǔ)} \\
\hline
Nom de famille & \zh{一般跟父亲姓，可跟母亲姓} \emph{yībān gēn fùqīn xìng, kě gēn mǔqīn xìng} & \zh{Presque toujours père (patrilinéaire)} \emph{jīhū dōu gēn fùqīn xìng (fùxì)} \\
\hline
Culte ancêtres & \zh{清明节扫墓、祭祖、守孝} \emph{Qīngmíng jié sǎomù, jìzǔ, shǒuxiào} & \zh{Veillées, deuil, rites coutumiers, libations} \emph{shǒuxiào, xísú jìsì, jièjiǔ} \\
\hline
Testament & \zh{De plus en plus fréquent (ville)} \emph{yuè lái yuè chángjiàn (chénglǐ)} & \zh{Existe (droit moderne), rare en coutume} \emph{cúnzài (xiàndài fǎlǜ), xísú shǎo} \\
\hline
\end{tabularx}
\end{center}

\begin{exoresolu}[Traduisez en chinois (niveau examen)]
Énoncé : Traduisez en chinois : « Bien que le Cameroun ait des coutumes variées, la loi moderne protège les droits des filles. »
\tcblower
Solution détaillée :\\
1. « Bien que » $\rightarrow$ \zh{虽然} (\emph{suīrán}) en tête de proposition.\\
2. « Le Cameroun a des coutumes variées » : \zh{喀麦隆有各种习俗} (\emph{Kāmàilóng yǒu gèzhǒng xísú}). \zh{各种} = variées, diverses.\\
3. « La loi moderne » : \zh{现代法律} (\emph{xiàndài fǎlǜ}).\\
4. « Protège les droits des filles » : \zh{保护女儿的权利} (\emph{bǎohù nǚ'ér de quánlì}) ou \zh{保障女儿的继承权} (\emph{bǎozhàng nǚ'ér de jìchéngquán}).\\
5. Structure : \zh{虽然喀麦隆有各种习俗，但是现代法律保护女儿的权利。}\\
Réponse : \zh{虽然喀麦隆有各种习俗，但是现代法律保护女儿的权利。}\\
\emph{Suīrán Kāmàilóng yǒu gèzhǒng xísú, dànshì xiàndài fǎlǜ bǎohù nǚ'ér de quánlì.}
\end{exoresolu}

\begin{attention}[Pièges fréquents des francophones — Révision finale]
\begin{itemize}
\item \textbf{« Les mêmes droits »} : Ne dites pas \zh{同样的权利} (structure nominale) comme prédicat principal. Utilisez \zh{有一样的权利} ou \zh{跟……一样的权利}. Ex: \zh{女儿跟儿子有一样的继承权。}
\item \textbf{« Décéder »} : \zh{去世} (\emph{qùshì}) est le terme respectueux standard. \zh{死} (\emph{sǐ}) est brutal, familier ou médical. \zh{去世了} (\emph{qùshì le}) = « est décédé(e) ».
\item \textbf{\zh{姓} (xìng)} : Verbe « porter le nom de ». \zh{他姓李。} (\emph{Tā xìng Lǐ.}) Pas de \zh{是} avant. \zh{跟父亲姓} = porter le nom du père.
\item \textbf{\zh{把} (bǎ)} : Objet \textbf{connu/défini} + action de disposition. \zh{把房子给了儿子} (OK). \zh{把一个房子给了儿子} (Maladroit, indéfini). \zh{给了儿子房子} (SVO standard, OK mais moins focalisé sur l'objet).
\item \textbf{Comparatif \zh{比}} : Pas de \zh{更} (\emph{gèng}) ni de \zh{很} (\emph{hěn}) après l'adjectif dans la structure de base. \zh{中国比喀麦隆统一} (OK). \zh{中国比喀麦隆更统一} (OK, \zh{更} avant l'adj). \zh{中国比喀麦隆统一很} (FAUX).
\item \textbf{Pluriel} : Pas de \zh{们} (\emph{men}) sur les noms inanimés ou collectifs (\zh{财产们} FAUX, \zh{习俗们} FAUX). \zh{子女} est déjà collectif.
\end{itemize}
\end{attention}

\begin{exercice}[Entraînement complet (type Bac)]
\begin{enumerate}
\item \textbf{Vocabulaire} : Associez le caractère à sa définition.
      \begin{itemize}
      \item A. 遗嘱 \quad B. 习俗 \quad C. 族长 \quad D. 父系 \quad E. 扫墓
      \item 1. Chef de clan/famille. \quad 2. Testament. \quad 3. Nettoyer les tombes. \quad 4. Coutume. \quad 5. Patrilinéaire.
      \end{itemize}
\item \textbf{Grammaire} : Transformez la phrase à la voix passive avec \zh{由}.
      \zh{第一顺序的继承人继承遗产。} $\rightarrow$ \zh{遗产由第一顺序的继承人继承。}
\item \textbf{Structure \zh{把}} : Réordonnez : \zh{了 / 房子 / 把 / 给 / 爷爷 / 儿子 / 大}.
      $\rightarrow$ \zh{爷爷把房子给了大儿子。}
\item \textbf{Comparaison} : Complétez avec \zh{比}、\zh{跟...一样}、\zh{虽然...但是} ou \zh{可是}.
      \begin{enumerate}
      \item 中国的法律 \_\_\_\_\_\_ 喀麦隆的习俗 \_\_\_\_\_\_ 统一。 (plus unifié)
      \item \_\_\_\_\_\_ 中国有统一法律，\_\_\_\_\_\_ 喀麦隆有习惯法。 (Bien que... mais...)
      \item 儿子 \_\_\_\_\_\_ 女儿 \_\_\_\_\_\_，都有继承权。 (comme... pareil)
      \end{enumerate}
\item \textbf{Expression écrite} : Rédigez 5 phrases en caractères + pinyin pour présenter \emph{une} différence majeure et \emph{une} similitude entre l'héritage en Chine et au Cameroun, en utilisant le vocabulaire de la leçon.
\end{enumerate}
\end{exercice}

\begin{corrige}[Exercice « Entraînement complet »]
1. A-2, B-4, C-1, D-5, E-3.\\
2. \zh{遗产由第一顺序的继承人继承。} \emph{Yíchǎn yóu dì-yī shùnxù de jìchéngrén jìchéng.}\\
3. \zh{爷爷把房子给了大儿子。} \emph{Yéye bǎ fángzi gěi le dà érzi.}\\
4. a) \zh{比} ... \zh{更} (ou \zh{比} ... \zh{统一}) — \emph{Zhōngguó de fǎlǜ bǐ Kāmàilóng de xísú gèng tǒngyī.} \\
   b) \zh{虽然} ... \zh{可是/但是} — \emph{Suīrán Zhōngguó yǒu tǒngyī fǎlǜ, kěshì/dànshì Kāmàilóng yǒu xíguànfǎ.} \\
   c) \zh{跟} ... \zh{一样} — \emph{Érzi gēn nǚ'ér yīyàng, dōu yǒu jìchéngquán.}\\
5. Exemple de réponse attendue :\\
   \zh{中国有统一的继承法，喀麦隆有现代法律和习惯法。} (\emph{Zhōngguó yǒu tǒngyī de jìchéngfǎ, Kāmàilóng yǒu xiàndài fǎlǜ hé xíguànfǎ.}) — Différence : cadre légal.\\
   \zh{中国和喀麦隆都重视纪念祖先。} (\emph{Zhōngguó hé Kāmàilóng dōu zhòngshì jìniàn zǔxiān.}) — Similitude : mémoire ancestrale.\\
   \zh{虽然中国法律规定男女平等，但是传统观念里长子责任大。} (\emph{Suīrán Zhōngguó fǎlǜ guīdìng nánnǚ píngděng, dànshì chuántǒng guānniàn lǐ zhǎngzǐ zérèn dà.}) — Nuance tradition/loi.\\
   \zh{喀麦隆的习俗有时不平等，女儿得不到遗产。} (\emph{Kāmàilóng de xísú yǒushí bù píngděng, nǚ'ér débu dào yíchǎn.}) — Point critique.\\
   \zh{我认为两个国家都在改变，法律越来越保护女儿。} (\emph{Wǒ rènwéi liǎng gè guójiā dōu zài gǎibiàn, fǎlǜ yuè lái yuè bǎohù nǚ'ér.}) — Avis personnel.
\end{corrige}

\coursec{Vocabulaire du thème}

Après avoir examiné les faits relatifs à l'hérédité en Chine, nous disposons maintenant du contexte nécessaire pour aborder le lexique spécialisé qui permet de décrire, comparer et analyser ces pratiques. Ce vocabulaire est la \cle{boîte à outils} indispensable pour réussir le mini-exposé final : il couvre les acteurs (héritiers, veuves), les objets (biens, titres, testaments), le cadre juridique (lois, droits, stipulations) et les concepts culturels (traditions, coutumes, comparaison).

\courssub{Lexique de base : acteurs, biens et cadre juridique}

Observons d'abord ces mots en situation à travers un court dialogue entre deux étudiants, l'un camerounais (Paul) et l'autre chinois (Li Ming), qui préparent leur exposé commun.

\begin{textebox}[Dialogue de préparation : le vocabulaire de l'hérédité]
保罗：李明，我们要准备关于继承的词汇。\\
Pǔluò: Lǐ Míng, wǒmen yào zhǔnbèi guānyú jìchéng de cíhuì.\\
Paul : Li Ming, nous devons préparer le vocabulaire sur l'héritage.

李明：好。首先，核心词是 \zh{继承} (jìchéng) 和 \zh{继承人} (jìchéngrén)。\\
Lǐ Míng: Hǎo. Shǒuxiān, héxīn cí shì jìchéng hé jìchéngrén.\\
Li Ming : Bien. D'abord, les mots-clés sont « hériter » et « héritier ».

保罗：那 \zh{遗产} (yíchǎn) 和 \zh{遗嘱} (yízhǔ) ?\\
Pǔluò: Nà yíchǎn hé yízhǔ?\\
Paul : Et « héritage/patrimoine » et « testament » ?

李明：\zh{遗产} 是留下的财产，\zh{遗嘱} 是立遗嘱的人写的文件。还有 \zh{财产} (cáichǎn) 和 \zh{头衔} (tóuxián)。\\
Lǐ Míng: Yíchǎn shì liúxià de cáichǎn, yízhǔ shì lì yízhǔ de rén xiě de wénjiàn. Hái yǒu cáichǎn hé tóuxián.\\
Li Ming : L'héritage ce sont les biens laissés, le testament est le document écrit par le testateur. Il y a aussi « biens/fortune » et « titre ».

保罗：在喀麦隆，\zh{长子} (zhǎngzǐ) 传统上继承一切，\zh{寡妇} (guǎfu) 处境困难。\\
Pǔluò: Zài Kāmàilóng, zhǎngzǐ chuántǒng shàng jìchéng yīqiè, guǎfu chǔjìng kùnnán.\\
Paul : Au Cameroun, le fils aîné hérite traditionnellement de tout, la veuve est en situation difficile.

李明：中国法律现在规定女儿也有 \zh{继承权} (jìchéngquán). \zh{法律} (fǎlǜ) 保障 \zh{权利} (quánlì)，\zh{规定} (guīdìng) 男女 \zh{一样} (yīyàng)。\\
Lǐ Míng: Zhōngguó fǎlǜ xiànzài guīdìng nǚ'ér yě yǒu jìchéngquán. Fǎlǜ bǎozhàng quánlì, guīdìng nánnǚ yīyàng.\\
Li Ming : La loi chinoise stipule maintenant que les filles ont aussi le droit d'hériter. La loi garantit les droits, stipule que hommes et femmes sont égaux.

保罗：但 \zh{传统} (chuántǒng) 和 \zh{习俗} (xísú) 变化慢。我们要 \zh{比较} (bǐjiào) 两国的 情况。\\
Pǔluò: Dàn chuántǒng hé xísú biànhuà màn. Wǒmen yào bǐjiào liǎng guó de qíngkuàng.\\
Paul : Mais les traditions et coutumes changent lentement. Nous devons comparer la situation des deux pays.
\end{textebox}

\begin{definition}[Glossaire du dialogue]
\begin{center}
\begin{tabularx}{\linewidth}{|X|X|X|X|}
\hline
\rowcolor{popblueL}
\textbf{Caractères} & \textbf{Pinyin} & \textbf{Nature} & \textbf{Traduction} \\
\hline
继承 & jìchéng & v. & hériter (de) \\
继承人 & jìchéngrén & n. & héritier \\
遗产 & yíchǎn & n. & héritage, patrimoine \\
遗嘱 & yízhǔ & n. & testament \\
财产 & cáichǎn & n. & biens, fortune, patrimoine \\
头衔 & tóuxián & n. & titre (honorifique, noble, fonction) \\
长子 & zhǎngzǐ & n. & fils aîné \\
寡妇 & guǎfu & n. & veuve \\
法律 & fǎlǜ & n. & loi \\
权利 & quánlì & n. & droit (juridique) \\
规定 & guīdìng & v./n. & stipuler, prescrire / stipulation, règle \\
传统 & chuántǒng & n./adj. & tradition / traditionnel \\
习俗 & xísú & n. & coutume, mœurs \\
一样 & yīyàng & adj. & pareil, identique, égal \\
比较 & bǐjiào & v./adv. & comparer / relativement, assez \\
\hline
\end{tabularx}
\end{center}
\end{definition}

\courssub{Analyse morphologique et familles de mots}

La méthode la plus efficace pour mémoriser ce lexique est de le regrouper par \textbf{familles morphologiques} (mots partageant un même caractère-clé). Cela révèle la logique interne du chinois et multiplie votre vocabulaire actif.

\begin{methode}[Mémoriser par familles de mots — 4 étapes]
\begin{enumerate}
\item \textbf{Identifiez le caractère-clé} (ex. : 继 jì « continuer, succéder »).
\item \textbf{Listez les mots composés} qui l'utilisent (继承, 继承人, 继承权, 继承法).
\item \textbf{Analysez la relation sémantique} : 继 + 承 (recevoir) = hériter ; 继 + 承 + 人 (personne) = héritier ; 继 + 承 + 权 (droit) = droit d'héritage.
\item \textbf{Produisez une phrase par mot} pour ancrer l'usage en contexte.
\end{enumerate}
\end{methode}

\begin{exemplebox}[Familles de mots du thème]
\textbf{Famille 继 (jì) — « succéder, continuer »}
\begin{itemize}
\item \zh{继承} jìchéng — hériter (litt. « succéder + recevoir »)
\item \zh{继承人} jìchéngrén — héritier
\item \zh{继承权} jìchéngquán — droit d'héritage
\item \zh{继承法} jìchéngfǎ — loi successorale
\end{itemize}

\textbf{Famille 遗 (yí) — « laisser derrière, legs »}
\begin{itemize}
\item \zh{遗产} yíchǎn — héritage, patrimoine (legs + biens)
\item \zh{遗嘱} yízhǔ — testament (legs + injonction/ordre écrit)
\item \zh{遗留} yíliú — laisser en héritage, legs (v.)
\item \zh{遗孀} yíshuāng — veuve (legs + femme seule) — terme formel/littéraire
\end{itemize}

\textbf{Famille 财 (cái) — « richesse, argent » (radical 贝)}
\begin{itemize}
\item \zh{财产} cáichǎn — biens, fortune
\item \zh{财富} cáifù — richesse
\item \zh{财务} cáiwù — finances, comptabilité
\item \zh{发财} fācái — s'enrichir, faire fortune
\end{itemize}

\textbf{Famille 权 (quán) — « droit, pouvoir, autorité »}
\begin{itemize}
\item \zh{权利} quánlì — droit (juridique)
\item \zh{权力} quánlì — pouvoir, autorité (politique, administratif) — \emph{attention : même pinyin, sens différent !}
\item \zh{人权} rénquán — droits de l'homme
\item \zh{选举权} xuǎnjǔquán — droit de vote
\end{itemize}

\textbf{Famille 定 (dìng) — « fixer, décider, stable »}
\begin{itemize}
\item \zh{规定} guīdìng — stipuler, règle
\item \zh{决定} juédìng — décider / décision
\item \zh{一定} yídìng — certain, sûr / certainement
\item \zh{稳定} wěndìng — stable, stabiliser
\end{itemize}
\end{exemplebox}

\begin{propriete}[Le radical 贝 (bèi) — « coquillage = monnaie ancienne »]
Le radical \textbf{贝} (bèi), qui représente une coquille de cauri (monnaie traditionnelle en Afrique et en Chine ancienne), apparaît dans de nombreux caractères liés à l'argent, la richesse, l'échange ou la valeur.
\begin{center}
\begin{tabularx}{\linewidth}{|X|X|X|X|}
\hline
\rowcolor{poporangeL}
\textbf{Caractère} & \textbf{Pinyin} & \textbf{Sens} & \textbf{Lien avec l'argent} \\
\hline
财 & cái & richesse, biens & argent + talent (才) \\
贵 & guì & cher, noble & argent + comparaison (比) — « qui vaut plus » \\
贫 & pín & pauvre & argent + diviser (分) — « partager l'argent → peu » \\
贸 & mào & commerce & argent + échanger \\
赚 & zhuàn & gagner (argent) & argent + avantage (兼) \\
赔 & péi & indemniser, perdre & argent + accompagner (陪) — « payer de sa poche » \\
债 & zhài & dette & personne (亻) + responsabilité (责) — radical 贝 en bas \\
\hline
\end{tabularx}
\end{center}
\emph{Astuce :} quand vous voyez 贝 dans un caractère, pensez « argent / valeur / échange ».
\end{propriete}

\courssub{Points de grammaire lexicale : 规定, 比较, 一样}

Trois mots du corpus fonctionnent comme des \textbf{outils grammaticaux} pour structurer la comparaison et l'argumentation.

\begin{propriete}[\zh{规定} guīdìng — v. « stipuler, prescrire » / n. « stipulation, règle, disposition »]
\begin{itemize}
\item \textbf{Comme verbe} : Sujet (loi, règlement, école) + 规定 + [que] + Proposition.
  \begin{itemize}
  \item \zh{法律规定女儿有继承权。} Fǎlǜ guīdìng nǚ'ér yǒu jìchéngquán. (La loi stipule que les filles ont le droit d'hériter.)
  \item \zh{学校规定学生必须穿校服。} Xuéxiào guīdìng xuéshēng bìxū chuān xiàofú. (L'école stipule que les élèves doivent porter l'uniforme.)
  \end{itemize}
\item \textbf{Comme nom} : souvent précédé de 的 (de) ou suivi de 内容 (nèiróng, contenu) / 精神 (jīngshén, esprit).
  \begin{itemize}
  \item \zh{根据有关规定……} Gēnjù yǒuguān guīdìng… (Selon les dispositions concernées…)
  \item \zh{这项规定很明确。} Zhè xiàng guīdìng hěn míngquè. (Cette règle est très claire.)
  \end{itemize}
\item \textbf{Ne pas confondre} avec \zh{决定} juédìng (décider) : 规定 = règle générale, impersonnelle ; 决定 = choix ponctuel, souvent personnel.
\end{itemize}
\end{propriete}

\begin{propriete}[\zh{比较} bǐjiào — v. « comparer » / adv. « relativement, assez, plutôt »]
\begin{itemize}
\item \textbf{Verbe transitif} : A + 比较 + B (+ 方面 / 情况 / 差异).
  \begin{itemize}
  \item \zh{我们比较两国的继承法。} Wǒmen bǐjiào liǎng guó de jìchéngfǎ. (Nous comparons les lois successorales des deux pays.)
  \end{itemize}
\item \textbf{Adverbe de degré} (très courant à l'oral et à l'écrit) : placé devant un adjectif ou un verbe d'état mental, signifie « relativement, assez ».
  \begin{itemize}
  \item \zh{这个问题比较复杂。} Zhège wèntí bǐjiào fùzá. (Cette question est assez complexe.)
  \item \zh{中国的法律比较完善。} Zhōngguó de fǎlǜ bǐjiào wánshàn. (La loi chinoise est relativement complète/bien développée.)
  \item \zh{我比较喜欢喀麦隆的传统音乐。} Wǒ bǐjiào xǐhuan Kāmàilóng de chuántǒng yīnyuè. (J'aime assez la musique traditionnelle camerounaise.)
  \end{itemize}
\item \textbf{Niveau de langue} : plus neutre/poli que \zh{很} hěn (très) ; exprime une nuance, une réserve.
\end{itemize}
\end{propriete}

\begin{propriete}[\zh{一样} yīyàng — adj. « pareil, identique, égal »]
\begin{itemize}
\item \textbf{Structure affirmative} : A + 跟/和 + B + 一样 (+ adj./verbe d'état).
  \begin{itemize}
  \item \zh{男女的继承权一样。} Nánnǚ de jìchéngquán yīyàng. (Les droits d'héritage des hommes et des femmes sont identiques.)
  \item \zh{喀麦隆的习俗和中国的不一样。} Kāmàilóng de xísú hé Zhōngguó de bù yīyàng. (Les coutumes du Cameroun ne sont pas pareilles à celles de la Chine.)
  \end{itemize}
\item \textbf{Structure « de la même façon » (adverbialisation)} : 一样 + 地 (de) + Verbe.
  \begin{itemize}
  \item \zh{儿子女儿一样地继承财产。} Érzi nǚ'ér yīyàng de jìchéng cáichǎn. (Fils et filles héritent des biens de la même façon.)
  \end{itemize}
\item \textbf{Attention} : \zh{一样} n'est \textbf{pas} un verbe transitif (*他们一样 quelque chose). En prédicat, on dit \zh{他们一样} (ils sont pareils) ou \zh{他们相同} (tāmen xiāngtóng — ils sont identiques, plus formel).
\end{itemize}
\end{propriete}

\courssub{Radicaux clés et ordre des traits : écrire pour mémoriser}

L'écriture manuelle des caractères ancre la mémoire visuelle et motrice. Voici les trois caractères « pivots » du thème avec leur décomposition radicale et l'ordre des traits officiel.

\begin{definition}[Radicaux-clés du thème]
\begin{center}
\begin{tabularx}{\linewidth}{|X|X|X|X|}
\hline
\rowcolor{popgreenL}
\textbf{Radical} & \textbf{Nom} & \textbf{Sens originel} & \textbf{Exemples du thème} \\
\hline
纟 & sī (soie) & fil, continuité & \zh{继} (jì) — continuer la lignée \\
宀 & mián (toit) & maison, foyer, protection & \zh{家} (jiā), \zh{寡} (guǎ dans 寡妇) \\
贝 & bèi (coquillage) & monnaie, richesse & \zh{财} (cái), \zh{贵} (guì), \zh{贫} (pín) \\
土 & tǔ (terre) & sol, terre, tombeau & \zh{坟} (fén), \zh{墓} (mù) — tombe \\
\hline
\end{tabularx}
\end{center}
\end{definition}

\begin{exemplebox}[Ordre des traits — Démonstration guidée]
\textbf{1. 继 (jì) — 10 traits} — Radical 纟 (3 traits) + partie droite (7 traits)
\begin{enumerate}
\item 纟 : ① piě (撇) ② zhé (折) ③ diǎn (点)
\item Partie droite : ④ piě (撇) ⑤ héng (横) ⑥ piě (撇) ⑦ nà (捺) ⑧ héng (横) ⑨ shù (竖) ⑩ héng (横 — trait final long qui « scelle » la succession)
\end{enumerate}
\emph{Mnémotechnique :} « Sous le toit de la soie (纟), on tend la main (partie droite) pour saisir la succession. »

\textbf{2. 遗 (yí) — 12 traits} — Radical 辶 (chuò, marche, 3 traits) + 贵 (guì, 9 traits)
\begin{enumerate}
\item 贵 : ① héng ② shù ③ héng ④ héng ⑤ héng ⑥ shù ⑦ piě ⑧ diǎn ⑨ diǎn
\item 辶 : ⑩ diǎn (11) piě (12) nà (le trait final « marche » part du bas du caractère)
\end{enumerate}
\emph{Mnémotechnique :} « Ce qui est \zh{贵} (précieux) s'en va (辶) — c'est le legs. »

\textbf{3. 权 (quán) — 6 traits} — Caractère unique (pas de radical séparé en écriture simplifiée ; historiquement 木 + 券)
\begin{enumerate}
\item ① héng (横) ② shù (竖) ③ héng (横) ④ piě (撇) ⑤ diǎn (点) ⑥ héng (横 — trait final long qui « équilibre » le tout)
\end{enumerate}
\emph{Mnémotechnique :} « Un arbre (木) dont on a taillé les branches pour ne garder que le tronc droit = l'autorité, le droit. »
\end{exemplebox}

\begin{attention}[Piège de prononciation : \zh{长} zhǎng vs cháng]
Le caractère \zh{长} a deux lectures selon le sens :
\begin{itemize}
\item \textbf{zhǎng} (3e ton) = « aîné, chef, grandir (verbe), diriger » — \textbf{noms de personnes/fonctions}.
  \begin{itemize}
  \item \zh{长子} zhǎngzǐ — fils aîné
  \item \zh{长女} zhǎngnǚ — fille aînée
  \item \zh{班长} bānzhǎng — délégué de classe
  \item \zh{部长} bùzhǎng — ministre
  \item \zh{长大} zhǎngdà — grandir (v.)
  \end{itemize}
\item \textbf{cháng} (2e ton) = « long, longueur, durée » — \textbf{adjectif/nom de mesure}.
  \begin{itemize}
  \item \zh{长短} chángduǎn — longueur
  \item \zh{长度} chángdù — longueur (physique)
  \item \zh{长江} Cháng Jiāng — le Yangtsé (le « Long Fleuve »)
  \item \zh{长久} chángjiǔ — durable, long temps
  \end{itemize}
\end{itemize}
\emph{Erreur fréquente :} dire *chángzǐ pour « fils aîné ». \textbf{Règle mnémotechnique} : « Celui qui \textbf{dirige} (zhǎng) la fratrie est l'aîné ; ce qui est \textbf{long} (cháng) se mesure. »
\end{attention}

\courssub{Exemples traduits et mise en contexte}

\begin{exemplebox}[Phrases modèles pour l'exposé]
\begin{enumerate}
\item \zh{女儿也有继承权。} Nǚ'ér yě yǒu jìchéngquán. — Les filles aussi ont le droit d'hériter.
\item \zh{每个地方的传统不一样。} Měi gè dìfang de chuántǒng bù yīyàng. — Les traditions diffèrent selon les endroits.
\item \zh{中国法律规定：男女享有同等的继承权。} Zhōngguó fǎlǜ guīdìng: nánnǚ xiǎngyǒu tóngděng de jìchéngquán. — La loi chinoise stipule : hommes et femmes jouissent de droits d'héritage égaux.
\item \zh{喀麦隆某些地区的习俗是：长子继承父亲的头衔和财产。} Kāmàilóng mǒuxiē dìqū de xísú shì: zhǎngzǐ jìchéng fùqīn de tóuxián hé cáichǎn. — Dans certaines régions du Cameroun, la coutume veut que le fils aîné hérite du titre et des biens du père.
\item \zh{寡妇的权利在法律上受保护，但在习俗中比较困难。} Guǎfu de quánlì zài fǎlǜ shàng shòu bǎohù, dàn zài xísú zhōng bǐjiào kùnnán. — Les droits de la veuve sont protégés par la loi, mais dans la coutume c'est relativement difficile.
\item \zh{立遗嘱可以避免家庭纠纷。} Lì yízhǔ kěyǐ bìmiǎn jiātíng jiūfēn. — Faire un testament peut éviter les conflits familiaux.
\item \zh{我们要比较两国的异同，不能只看表面。} Wǒmen yào bǐjiào liǎng guó de yìtóng, bùnéng zhǐ kàn biǎomiàn. — Nous devons comparer les similitudes et différences des deux pays, ne pas nous fier aux apparences.
\end{enumerate}
\end{exemplebox}

\courssub{Texte support original : « Deux regards sur l'héritage »}

Le texte suivant (environ 250 caractères chinois) met en œuvre l'essentiel du vocabulaire de la leçon. Il sert de base à la compréhension écrite, à la traduction et à la préparation de l'exposé.

\begin{textebox}[Texte original : 两国继承观念的对比 — Comparaison des conceptions de l'héritage dans les deux pays]
喀麦隆和中国都重视家庭传承，但具体做法有很大不同。\\
Kāmàilóng hé Zhōngguó dōu zhòngshì jiātíng chuánchéng, dàn jùtǐ zuòfǎ yǒu hěn dà bùtóng.\\
Le Cameroun et la Chine accordent tous deux de l'importance à la transmission familiale, mais les pratiques concrètes sont très différentes.

在喀麦隆许多民族中，习俗规定长子继承父亲的头衔、财产和主要遗产。\\
Zài Kāmàilóng xǔduō mínzú zhōng, xísú guīdìng zhǎngzǐ jìchéng fùqīn de tóuxián, cáichǎn hé zhǔyào yíchǎn.\\
Chez de nombreuses ethnies du Cameroun, la coutume stipule que le fils aîné hérite du titre, des biens et du patrimoine principal du père.

寡妇和女儿往往只能获得少量生活费，继承权受限。\\
Guǎfu hé nǚ'ér wǎngwǎng zhǐ néng huòdé shǎoliàng shēnghuófèi, jìchéngquán shòuxiàn.\\
La veuve et les filles ne reçoivent souvent qu'une petite pension de subsistance, leur droit d'héritage est limité.

中国的继承法则完全不同。法律规定男女继承权一样，第一顺序继承人包括配偶、子女和父母。\\
Zhōngguó de jìchéngfǎ zé wánquán bùtóng. Fǎlǜ guīdìng nánnǚ jìchéngquán yīyàng, dì yī shùnxù jìchéngrén bāokuò pèiǒu, zǐnǚ hé fùmǔ.\\
La loi successorale chinoise est complètement différente. La loi stipule que les droits d'héritage des hommes et des femmes sont identiques ; les héritiers du premier ordre incluent le conjoint, les enfants et les parents.

现在城市家庭常立遗嘱，自由分配遗产，比较灵活。\\
Xiànzài chéngshì jiātíng cháng lì yízhǔ, zìyóu fēnpèi yíchǎn, bǐjiào línghuó.\\
Aujourd'hui, les familles urbaines font souvent un testament, répartissent librement l'héritage, c'est relativement flexible.

但是，农村地区的传统观念依然强烈，有些家庭还是希望长子多继承一点。\\
Dànshì, nóngcūn dìqū de chuántǒng guānniàn yīrán qiángliè, yǒuxiē jiātíng hái shì xīwàng zhǎngzǐ duō jìchéng yīdiǎn.\\
Cependant, dans les zones rurales, les conceptions traditionnelles restent fortes ; certaines familles espèrent encore que le fils aîné hérite un peu plus.

比较两国，我们发现：法律可以改变，但习俗改变得比较慢。\\
Bǐjiào liǎng guó, wǒmen fāxiàn: fǎlǜ kěyǐ gǎibiàn, dàn xísú gǎibiàn de bǐjiào màn.\\
En comparant les deux pays, nous découvrons : la loi peut changer, mais les coutumes changent relativement lentement.

理解这些异同，有助于我们尊重文化多样性。\\
Lǐjiě zhèxiē yìtóng, yǒuzhù yú wǒmen zūnzhòng wénhuà duōyàngxìng.\\
Comprendre ces similitudes et différences nous aide à respecter la diversité culturelle.
\end{textebox}

\begin{exoresolu}[Exploitation du texte — Questions de compréhension et traduction]
\textbf{Énoncé :}
\begin{enumerate}
\item Relevez dans le texte tous les mots du vocabulaire de la leçon (liste de la définition) et soulignez-les.
\item Traduisez en français les deux phrases suivantes :
  \begin{enumerate}
  \item \zh{寡妇和女儿往往只能获得少量生活费，继承权受限。}
  \item \zh{现在城市家庭常立遗嘱，自由分配遗产，比较灵活。}
  \end{enumerate}
\item Répondez en chinois (phrase complète) :
  \begin{enumerate}
  \item \zh{根据法律，中国的第一顺序继承人是谁？}
  \item \zh{为什么说“习俗改变得比较慢”？}
  \end{enumerate}
\end{enumerate}

\tcblower
\textbf{Solution détaillée :}
\begin{enumerate}
\item \textbf{Mots repérés (18 occurrences) :} 传承 (chuánchéng — transmission, variante de 传统), 习俗 (xísú), 规定 (guīdìng), 长子 (zhǎngzǐ), 头衔 (tóuxián), 财产 (cáichǎn), 遗产 (yíchǎn), 寡妇 (guǎfu), 女儿 (nǚ'ér), 继承权 (jìchéngquán), 受限 (shòuxiàn), 继承法 (jìchéngfǎ), 法律 (fǎlǜ), 规定 (guīdìng), 男女 (nánnǚ), 一样 (yīyàng), 继承人 (jìchéngrén), 配偶 (pèiǒu), 子女 (zǐnǚ), 父母 (fùmǔ), 遗嘱 (yízhǔ), 分配 (fēnpèi), 灵活 (línghuó), 传统 (chuántǒng), 观念 (guānniàn), 比较 (bǐjiào), 异同 (yìtóng).
\item \textbf{Traduction :}
  \begin{enumerate}
  \item « La veuve et les filles ne peuvent souvent obtenir qu'une petite somme pour les frais de vie, leur droit d'héritage est restreint. »
  \item « Aujourd'hui, les familles citadines font souvent un testament, répartissent librement l'héritage, c'est relativement flexible. »
  \end{enumerate}
\item \textbf{Réponses en chinois :}
  \begin{enumerate}
  \item \zh{根据法律，中国的第一顺序继承人包括配偶、子女和父母。} (Gēnjù fǎlǜ, Zhōngguó de dì yī shùnxù jìchéngrén bāokuò pèiǒu, zǐnǚ hé fùmǔ.)
  \item \zh{因为农村地区的传统观念依然强烈，有些家庭还是希望长子多继承一点。} (Yīnwèi nóngcūn dìqū de chuántǒng guānniàn yīrán qiángliè, yǒuxiē jiātíng hái shì xīwàng zhǎngzǐ duō jìchéng yīdiǎn.)
  \end{enumerate}
\end{enumerate}
\end{exoresolu}

\courssub{Exercice d'application : familles de mots et structures}

\begin{exercice}[Entraînement vocabulaire + grammaire]
\begin{enumerate}
\item \textbf{Complétez} avec le bon mot de la famille indiquée (forme correcte) :
  \begin{enumerate}
  \item Famille \zh{继} : 根据《\_\_\_\_\_\_法》，女儿和儿子享有同等的\_\_\_\_\_\_。 (loi successorale / droit d'héritage)
  \item Famille \zh{遗} : 老人生前立了\_\_\_\_\_\_，把房子留给\_\_\_\_\_\_。 (testament / veuve — terme formel)
  \item Famille \zh{财} : 这家公司\_\_\_\_\_\_雄厚，每年\_\_\_\_\_\_都在增长。 (richesse / finances)
  \item Famille \zh{权} : 公民享有受教育\_\_\_\_\_\_和\_\_\_\_\_\_。 (droit / droit de vote)
  \end{enumerate}
\item \textbf{Transformez} les phrases en utilisant \zh{比较} (adverbe) + adjectif :
  \begin{enumerate}
  \item 这个问题很复杂。 → \_\_\_\_\_\_
  \item 中国的法律很完善。 → \_\_\_\_\_\_
  \item 我喜欢喀麦隆的菜。 → \_\_\_\_\_\_
  \end{enumerate}
\item \textbf{Traduisez} en chinois (caractères + pinyin) :
  \begin{enumerate}
  \item Au Cameroun, la coutume veut que le fils aîné hérite du titre.
  \item La loi garantit les droits de la veuve.
  \item Les traditions des deux pays ne sont pas pareilles.
  \end{enumerate}
\end{enumerate}
\end{exercice}

\begin{corrige}[Exercice Vocabulaire du thème]
\begin{enumerate}
\item \textbf{Complétion :}
  \begin{enumerate}
  \item \zh{继承} (jìchéng) / \zh{继承权} (jìchéngquán)
  \item \zh{遗嘱} (yízhǔ) / \zh{遗孀} (yíshuāng) — ou \zh{寡妇} (guǎfu) à l'oral
  \item \zh{财力} (cáilì) ou \zh{财富} (cáifù) / \zh{财务} (cáiwù) ou \zh{财产} (cáichǎn)
  \item \zh{权利} (quánlì) / \zh{选举权} (xuǎnjǔquán)
  \end{enumerate}
\item \textbf{Transformation avec \zh{比较} :}
  \begin{enumerate}
  \item \zh{这个问题比较复杂。} (Zhège wèntí bǐjiào fùzá.)
  \item \zh{中国的法律比较完善。} (Zhōngguó de fǎlǜ bǐjiào wánshàn.)
  \item \zh{我比较喜欢喀麦隆的菜。} (Wǒ bǐjiào xǐhuan Kāmàilóng de cài.)
  \end{enumerate}
\item \textbf{Traduction :}
  \begin{enumerate}
  \item \zh{喀麦隆的习俗规定：长子继承头衔。} (Kāmàilóng de xísú guīdìng: zhǎngzǐ jìchéng tóuxián.)
  \item \zh{法律保障寡妇的权利。} (Fǎlǜ bǎozhàng guǎfu de quánlì.)
  \item \zh{两国的传统不一样。} (Liǎng guó de chuántǒng bù yīyàng.)
  \end{enumerate}
\end{enumerate}
\end{corrige}

\begin{savaistu}[Le cauri (贝 bèi) : monnaie commune Afrique-Chine]
Le radical \zh{贝} (bèi) témoigne d'une histoire économique fascinante. La coquille de \textbf{cauri} (Cypraea moneta), petit coquillage blanc brillant, a servi de monnaie pendant des siècles \textbf{à la fois en Afrique de l'Ouest (dont le Cameroun) et en Chine ancienne}.
\begin{itemize}
\item En Chine, les cauris (贝贝 bèibèi) étaient utilisés dès la dynastie Shang (vers 1600-1046 av. J.-C.) ; le caractère \zh{贝} en est la représentation pictographique.
\item Au Cameroun, le cauri (appelé \emph{njangi} ou \emph{cowrie} en pidgin/anglais) a circulé comme monnaie d'échange, dot, paiement de tributs, bien into le XXe siècle.
\item C'est pourquoi \textbf{les mêmes radicaux « argent » structurent le vocabulaire financier des deux cultures} : \zh{财} (richesse), \zh{贸} (commerce), \zh{贷} (prêt), \zh{账} (compte)… Une clé de lecture commune pour nos élèves !
\end{itemize}
\end{savaistu}

\begin{aretenir}[Synthèse lexicale — Ce qu'il faut maîtriser pour l'exposé]
\begin{itemize}
\item \textbf{15 mots-clés} (caractères, pinyin, nature, traduction) — voir tableau p. 1.
\item \textbf{4 familles morphologiques} : 继\_, 遗\_, 财\_, 权\_ — savoir en produire 3-4 mots chacune.
\item \textbf{3 outils grammaticaux} : \zh{规定} (v./n.), \zh{比较} (v./adv.), \zh{一样} (adj. + structure 跟/和…一样 / 一样地+V).
\item \textbf{2 pièges} : \zh{长} (zhǎng vs cháng) ; \zh{权利} vs \zh{权力} (même pinyin, sens différents).
\item \textbf{4 radicaux} : 纟, 宀, 贝, 土 — reconnaître et écrire 继 (10 traits), 遗 (12 traits), 权 (6 traits).
\end{itemize}
\end{aretenir}

\begin{experience}[Atelier oral : « Mon mot, ma phrase, mon image »]
\textbf{Objectif :} Réactiver le vocabulaire par la production orale guidée.
\textbf{Déroulement (15 min) :}
\begin{enumerate}
\item Chaque élève tire au sort une carte « mot-clé » (parmi les 15 du tableau).
\item En 1 minute, il prépare : (a) la lecture correcte (tons !), (b) une phrase personnelle en chinois, (c) un geste/mime qui illustre le sens.
\item Passage à l'oral : l'élève dit le mot → fait le geste → dit sa phrase. La classe devine le sens en français.
\item Exemple : \zh{遗嘱} → mime « écrire un papier solennel » → \zh{我爷爷立了遗嘱。} (Wǒ yéye lìle yízhǔ.)
\end{enumerate}
\textbf{Variante débat :} « \zh{长子应该多继承一点，对吗？} » — deux groupes (pour/contre), 3 arguments chacun en chinois simple.
\end{experience}

\coursec{Comparer et donner son avis en chinois}

Dans la section précédente, nous avons acquis le vocabulaire de base sur l'hérédité. Pour réussir l'exposé comparatif demandé au programme, il faut désormais savoir \textbf{structurer une comparaison} et \textbf{argumenter un point de vue} en chinois correct. Cette section vous donne les trois outils syntaxiques indispensables : \cle{比} (différence), \cle{和……一样} (similitude) et \cle{我觉得}/\cle{我认为} (opinion), ainsi que les connecteurs logiques pour les relier.

\begin{definition}[Rappel de vocabulaire clé (Leçon 8)]
\begin{center}
\begin{tabularx}{\linewidth}{|X|X|X|X|}
\hline
\rowcolor{popblueL} Caractères & Pinyin & Nature & Traduction \\
\hline
继承 & jìchéng & v & hériter \\
财产 & cáichǎn & n & biens, patrimoine \\
传统 & chuántǒng & n & tradition \\
法律 & fǎlǜ & n & loi \\
习惯法 & xíguànfǎ & n & droit coutumier \\
成文法 & chéngwénfǎ & n & droit écrit, loi codifiée \\
统一 & tǒngyī & adj & unifié, uniforme \\
公平 & gōngpíng & adj & juste, équitable \\
照顾 & zhàogù & v & s'occuper de, prendre soin de \\
保护 & bǎohù & v & protéger \\
丰富 & fēngfù & adj & riche (contenu) \\
清楚 & qīngchu & adj & clair \\
受尊重 & shòu zūnzhòng & v (passif) & être respecté \\
一员 & yí yuán & n & membre \\
比 & bǐ & prep & (marque de comparatif) plus... que \\
和 & hé & conj & et, avec \\
一样 & yíyàng & adj/adv & pareil, de la même façon \\
觉得 & juéde & v & trouver, penser (avis personnel) \\
认为 & rènwéi & v & considérer, estimer (avis argumenté) \\
因为 & yīnwèi & conj & parce que \\
所以 & suǒyǐ & adv & donc, par conséquent \\
但是 & dànshì & conj & mais \\
不过 & bùguò & conj & cependant, toutefois \\
而且 & érqiě & conj & de plus, en plus \\
\hline
\end{tabularx}
\end{center}
\end{definition}

\courssub{Observer d'abord : trois phrases-clés}

\begin{textebox}[Trois opinions d'élèves de Terminale]
Lisez ces trois phrases prononcées lors d'un débat en classe à Yaoundé. Repérez les structures en gras qui permettent de comparer ou d'opiner.\\[4pt]
\textbf{中国的继承法}\textbf{比}\textbf{喀麦隆的习惯法更统一。}\\
\emph{Zhōngguó de jìchéngfǎ bǐ Kāmàilóng de xíguànfǎ gèng tǒngyī.}\\
« La loi chinoise sur l'héritage est plus unifiée que le droit coutumier camerounais. »\\[6pt]
\textbf{中国}\textbf{和}\textbf{喀麦隆一样}，家庭很重要。\\
\emph{Zhōngguó hé Kāmàilóng yíyàng, jiātíng hěn zhòngyào.}\\
« En Chine comme au Cameroun, la famille est très importante. »\\[6pt]
\textbf{我觉得}\textbf{女儿也应该继承财产。}\\
\emph{Wǒ juéde nǚ'ér yě yīnggāi jìchéng cáichǎn.}\\
« Je pense que les filles devraient aussi hériter des biens. »
\end{textebox}

Trois phrases, trois outils grammaticaux distincts. Étudions-les en détail.

\courssub{Structure 1 : la comparaison de différence avec 比}

\begin{propriete}[La comparaison : A 比 B + adjectif]
Pour exprimer qu'A possède une qualité à un degré supérieur à B, le chinois utilise une structure \textbf{rigide et immuable} :\\
\centering \textbf{A + 比 + B + adjectif (éventuellement précédé de 更 / 得多)}\\[4pt]
\raggedright
\begin{itemize}
\item \cle{比} (bǐ) signifie « comparé à ». Il se place \textbf{strictement entre} les deux termes comparés. L'ordre est crucial : le terme « supérieur » (A) est le sujet, placé \textbf{avant} 比.
\item \textbf{Interdiction absolue} : on n'écrit \textbf{jamais} \zh{很} (hěn) après 比. La comparaison porte déjà l'intensité. \zh{中国比喀麦隆很大} est \textbf{faux}.
\item Renforcement : \cle{更} (gèng, « encore plus ») ou \cle{得多} (dé duō, « beaucoup plus ») se placent \textbf{avant} l'adjectif. Ex. : \zh{A 比 B 更大}, \zh{A 比 B 大得多}.
\item Différence chiffrée : le complément de mesure se place \textbf{après} l'adjectif. Ex. : \zh{大三岁} (trois ans de plus), \zh{贵五块钱} (cinq yuans plus cher).
\item Négation : \textbf{A 没有 B (+ 这么/那么) + adjectif} = « A n'est pas aussi... que B ». L'ajout de \zh{这么} (zhème) ou \zh{那么} (nàme) « à ce point » est fréquent et plus naturel. Ex. : \zh{习惯法没有成文法(那么)统一。}
\end{itemize}
\end{propriete}

\begin{popfigure}
\begin{tikzpicture}[node distance=0.6cm, every node/.style={font=\small}]
\node[draw=popblue, fill=popblueL, rounded corners, minimum width=2.4cm, minimum height=1cm, align=center] (A) {\textbf{A (Sujet)}\\ 中国\\ \emph{Zhōngguó}};
\node[draw=poporange, fill=poporangeL, rounded corners, minimum width=1.6cm, minimum height=1cm, align=center, right=of A] (bi) {\textbf{比}\\ \emph{bǐ}\\ \textit{comparé à}};
\node[draw=popblue, fill=popblueL, rounded corners, minimum width=2.4cm, minimum height=1cm, align=center, right=of bi] (B) {\textbf{B (Objet)}\\ 喀麦隆\\ \emph{Kāmàilóng}};
\node[draw=popgreen, fill=popgreenL, rounded corners, minimum width=2.8cm, minimum height=1cm, align=center, right=of B] (adj) {\textbf{Adjectif (+ modif.)}\\ 更统一 / 大得多\\ \emph{gèng tǒngyī / dà de duō}};
\draw[-{Stealth}, thick, popdark] (A) -- (bi);
\draw[-{Stealth}, thick, popdark] (bi) -- (B);
\draw[-{Stealth}, thick, popdark] (B) -- (adj);
\node[below=0.8cm of bi, align=center, text=popmuted, font=\small] (trad) {中国比喀麦隆大得多。\ \emph{Zhōngguó bǐ Kāmàilóng dà de duō.}\\\textit{La Chine est bien plus grande que le Cameroun.}};
\draw[-{Stealth}, dashed, popmuted] (bi) -- (trad);
\end{tikzpicture}
\legende{Schéma de la structure comparative A 比 B Adj. L'ordre A > B est encodé par la position du sujet avant 比.}
\end{popfigure}

\begin{exemplebox}[La comparaison en action : héritage et géographie]
\begin{itemize}
\item \zh{中国比喀麦隆大得多。} \emph{Zhōngguó bǐ Kāmàilóng dà de duō.} « La Chine est beaucoup plus grande que le Cameroun. » (Différence chiffrée implicite : \zh{得多}).
\item \zh{中国的成文法比喀麦隆的习惯法更统一。} \emph{Zhōngguó de chéngwénfǎ bǐ Kāmàilóng de xíguànfǎ gèng tǒngyī.} « Le droit écrit chinois est plus unifié que le droit coutumier camerounais. » (Renforcement par \zh{更}).
\item \zh{在城市，女儿继承财产的机会比以前多了。} \emph{Zài chéngshì, nǚ'ér jìchéng cáichǎn de jīhuì bǐ yǐqián duō le.} « En ville, les occasions pour les filles d'hériter sont plus nombreuses qu'avant. » (Comparaison temporelle + \zh{了} marquant le changement).
\item \zh{弟弟比妹妹小三岁。} \emph{Dìdi bǐ mèimei xiǎo sān suì.} « Le petit frère est plus jeune que la petite sœur de trois ans. » (Différence chiffrée explicite).
\item \zh{喀麦隆的习惯法没有中国的成文法(那么)统一。} \emph{Kāmàilóng de xíguànfǎ méiyǒu Zhōngguó de chéngwénfǎ (nàme) tǒngyī.} « Le droit coutumier camerounais n'est pas aussi unifié que le droit écrit chinois. » (Négation standard).
\end{itemize}
\end{exemplebox}

\begin{attention}[Piège n°1 : le « que » fantôme et le 很 intrus]
En français : « plus grand \textbf{que} le Cameroun ». En chinois : \zh{中国比喀麦隆大} — \textbf{aucun mot ne traduit « que »}. \cle{比} fait tout le travail.
\\[4pt]
\ding{55} \zh{中国比喀麦隆很大} \quad (Faute : \zh{很} interdit après 比)
\\[2pt]
\ding{55} \zh{中国比喀麦隆是大} \quad (Faute : \zh{比} n'est pas un verbe, pas de \zh{是})
\\[2pt]
\ding{51} \zh{中国比喀麦隆大} \quad (Correct : Sujet + 比 + Objet + Adj nu)
\end{attention}

\courssub{Structure 2 : la similitude avec 和……一样}

\begin{propriete}[L'égalité / similitude : A 和 B 一样 (+ adjectif)]
Pour dire que deux éléments partagent une caractéristique au même degré :\\
\centering \textbf{A + 和 + B + 一样 (+ adjectif)}\\[4pt]
\raggedright
\begin{itemize}
\item \cle{和} (hé) = « et / avec » (préposition/connecteur). \cle{一样} (yíyàng) = « identique, pareil ». Les deux forment un bloc \textbf{inséparable} : \zh{和} seul ne suffit pas pour exprimer la similitude.
\item L'adjectif est \textbf{optionnel} : \zh{A 和 B 一样} « A et B sont pareils » ; \zh{A 和 B 一样重要} « A est aussi important que B ».
\item Variante courante : \cle{跟} (gēn) ou \cle{同} (tóng) peuvent remplacer \zh{和} (même structure). \zh{跟} est plus oral, \zh{同} plus formel/écrit.
\item Négation : \textbf{A 和 B 不一样} (+ adjectif possible). Ex. : \zh{中国的法律和喀麦隆的法律不一样。} « Les lois chinoise et camerounaise ne sont pas identiques. »
\item Avec \zh{都} (dōu, « tous/deux ») : pour insister sur le partage de la propriété par les deux termes. \zh{喀麦隆和中国一样，老人都很受尊重。} « Au Cameroun comme en Chine, les aînés sont tous respectés. »
\end{itemize}
\end{propriete}

\begin{exemplebox}[La similitude en action]
\begin{itemize}
\item \zh{中国和喀麦隆一样，家庭很重要。} \emph{Zhōngguó hé Kāmàilóng yíyàng, jiātíng hěn zhòngyào.} « En Chine comme au Cameroun, la famille est très importante. » (Sujet composé en thème).
\item \zh{儿子和女儿一样，都可以继承财产。} \emph{Érzi hé nǚ'ér yíyàng, dōu kěyǐ jìchéng cáichǎn.} « Fils et filles, de la même façon, peuvent tous hériter. » (\zh{都} distribue la possibilité).
\item \zh{我的看法和你的看法一样。} \emph{Wǒ de kànfǎ hé nǐ de kànfǎ yíyàng.} « Mon avis est le même que le tien. » (Comparaison de noms).
\item \zh{喀麦隆跟中国一样，都重视孝道。} \emph{Kāmàilóng gēn Zhōngguó yíyàng, dōu zhòngshì xiàodào.} « Le Cameroun comme la Chine valorisent tous deux la piété filiale. » (Variante orale avec \zh{跟}).
\end{itemize}
\end{exemplebox}

\begin{attention}[Piège n°2 : le mélange interdit 比 + 一样]
\zh{中国和喀麦隆一样重要} \ding{51} (Égalité).
\zh{中国比喀麦隆重要} \ding{51} (Différence : la Chine > Cameroun).
\zh{中国比喀麦隆一样重要} \ding{55} \textbf{INCORRECT}. \cle{比} marque l'inégalité, \cle{一样} marque l'égalité. \textbf{Choisissez un seul outil selon le sens visé.}
\end{attention}

\courssub{Structure 3 : exprimer son opinion avec 我觉得 / 我认为}

\begin{definition}[Deux registres pour opiner]
\begin{itemize}
\item \cle{我觉得} (wǒ juéde) = « je trouve que / j'ai l'impression que / je pense que ». \textbf{Registre courant, oral, subjectif}. Idéal pour une conversation ou un débat informel.
\item \cle{我认为} (wǒ rènwéi) = « je considère que / j'estime que / mon avis est que ». \textbf{Registre plus soutenu, argumenté, écrit}. \textbf{Recommandé pour l'exposé oral noté et la rédaction}.
\end{itemize}
\textbf{Syntaxe commune} : Sujet + 觉得/认为 + \textbf{Proposition complète (Sujet + Verbe + Complément)}. \textit{Aucun mot de liaison (que, de) n'est nécessaire.}
\end{definition}

\begin{exemplebox}[Des avis traduits mot à mot (analyse structurelle)]
\begin{itemize}
\item \zh{我觉得女儿也应该继承财产。} \\
\emph{Wǒ juéde nǚ'ér yě yīnggāi jìchéng cáichǎn.} \\
Mot à mot : moi / trouver / fille / aussi / devoir / hériter / biens.
\item \zh{我认为传统很重要，但是法律也重要。} \\
\emph{Wǒ rènwéi chuántǒng hěn zhòngyào, dànshì fǎlǜ yě zhòngyào.} \\
« J'estime que la tradition est importante, mais la loi l'est aussi. » (Coordination par \zh{但是}).
\item \zh{我认为每个人都应该尊重老人。} \\
\emph{Wǒ rènwéi měi ge rén dōu yīnggāi zūnzhòng lǎorén.} \\
« J'estime que chacun devrait respecter les aînés. » (\zh{每个人都} = « tout le monde »).
\end{itemize}
\end{exemplebox}

\courssub{Les connecteurs d'argumentation : construire le raisonnement}

Pour passer de l'opinion brute à l'exposé structuré, quatre connecteurs sont indispensables. Maîtrisez leur place dans la phrase.

\begin{center}
\begin{tabularx}{\linewidth}{|X|X|X|}
\hline
\rowcolor{popblueL} Connecteur & Exemple en caractères + pinyin & Traduction \\
\hline
因为……所以…… (yīnwèi… suǒyǐ…) & 因为家庭是核心，所以遗产问题很敏感。 \emph{Yīnwèi jiātíng shì héxīn, suǒyǐ yíchǎn wèntí hěn mǐngǎn.} & Parce que la famille est le noyau, la question de l'héritage est sensible. \\
\hline
但是 (dànshì) & 传统必须尊重，但是法律必须公平。 \emph{Chuántǒng bìxū zūnzhòng, dànshì fǎlǜ bìxū gōngpíng.} & La tradition doit être respectée, mais la loi doit être juste. \\
\hline
不过 (bùguò) & 法律改了，不过观念变得慢。 \emph{Fǎlǜ gǎi le, bùguò guānniàn biàn de màn.} & La loi a changé, cependant les mentalités évoluent lentement. \\
\hline
而且 (érqiě) & 女儿能继承，而且她们管理得好。 \emph{Nǚ'ér néng jìchéng, érqiě tāmen guǎnlǐ de hǎo.} & Les filles peuvent hériter, et de plus elles gèrent bien. \\
\hline
\end{tabularx}
\end{center}

\begin{aretenir}[Règles d'or des connecteurs pour francophones]
\begin{enumerate}
\item \textbf{因为……所以……} : Contrairement au français (*« Parce que..., donc... » est lourd), le chinois \textbf{exige} souvent les deux marqueurs pour marquer clairement la relation cause-effet. Gardez les deux !
\item \textbf{但是 vs 不过} : \zh{但是} = « mais » (opposition franche, écrit/oral). \zh{不过} = « toutefois / cependant » (opposition plus douce, nuance, très oral). En exposé, \zh{不过} montre une maturité argumentative.
\item \textbf{而且} : Ne se place \textbf{jamais} en début de phrase (contrairement à « De plus, » en français). Il relie deux propositions : \zh{Prop 1, 而且 Prop 2}.
\end{enumerate}
\end{aretenir}

\courssub{Méthode : Structurer un avis complet en 4 étapes (Modèle d'exposé)}

\begin{methode}[Produire un mini-exposé argumenté]
Pour répondre à une question de réflexion (ex. : « Faut-il que les filles héritent comme les fils ? »), suivez ce canevas en 4 temps. Chaque étape correspond à une phrase ou un groupe de phrases.

\begin{enumerate}
\item \textbf{OPINION} : Posez votre thèse avec \zh{我认为} (préférable à l'écrit/exposé) ou \zh{我觉得}.
\zh{我认为女儿和儿子一样，都应该继承财产。} \emph{(Je considère que filles et fils, pareillement, doivent hériter.)}
\item \textbf{RAISON} : Justifiez avec \zh{因为} (la cause). La conséquence (\zh{所以}) viendra à l'étape 4.
\zh{因为女儿也是家庭的一员，也照顾父母。} \emph{(Car la fille est aussi membre de la famille, elle soigne aussi les parents.)}
\item \textbf{EXEMPLE / FAIT} : Ancrez dans le réel (Cameroun ou Chine). Utilisez \zh{在……} (Dans...), \zh{现在……} (Maintenant...), \zh{比如……} (Par exemple...).
\zh{在很多喀麦隆家庭，女儿帮助父母很多年。} \emph{(Dans beaucoup de familles camerounaises, les filles aident les parents pendant des années.)}
\item \textbf{CONCLUSION} : Tirez la conséquence logique avec \zh{所以} et/ou réaffirmez l'égalité avec \zh{和……一样}.
\zh{所以，传统和法律一样重要，都要保护女儿的权利。} \emph{(Donc, tradition et loi sont aussi importantes, toutes deux doivent protéger les droits des filles.)}
\end{enumerate}

\textbf{Résultat final (texte continu pour l'oral) :}
\zh{我认为女儿和儿子一样，都应该继承财产，因为女儿也是家庭的一员，也照顾父母。在很多喀麦隆家庭，女儿帮助父母很多年，所以传统和法律一样重要，都要保护女儿的权利。}
\end{methode}

\begin{exoresolu}[Construire un avis complet : thème « Tradition vs Loi »]
Énoncé : Traduisez en chinois structuré (méthode 4 étapes) : « J'estime que la tradition camerounaise est plus riche qu'on ne le croit, car elle protège la famille ; cependant, la loi doit aussi protéger les filles. Donc tradition et loi sont toutes deux importantes. »
\tcblower
\textbf{Solution détaillée :}
\begin{enumerate}
\item \textbf{Opinion} : \zh{我认为喀麦隆的传统很丰富。} \emph{Wǒ rènwéi Kāmàilóng de chuántǒng hěn fēngfù.}
\item \textbf{Raison (Cause)} : \zh{因为它保护家庭。} \emph{Yīnwèi tā bǎohù jiātíng.}
\item \textbf{Nuance (Concession)} : \zh{不过，法律也应该保护女儿。} \emph{Bùguò, fǎlǜ yě yīnggāi bǎohù nǚ'ér.} (\zh{不过} adoucit l'opposition).
\item \textbf{Conclusion (Conséquence + Égalité)} : \zh{所以，传统和法律一样重要。} \emph{Suǒyǐ, chuántǒng hé fǎlǜ yíyàng zhòngyào.}
\end{enumerate}
\textbf{Texte final cohérent :}
\zh{我认为喀麦隆的传统很丰富，因为它保护家庭。不过，法律也应该保护女儿，所以传统和法律一样重要。}
\\[4pt]
\textit{Vérification : pas de \zh{很} après \zh{比} (pas de 比 ici), \zh{一样} bien présent après \zh{和}, pas de conjugaison, ordre SOV respecté dans les subordonnées.}
\end{exoresolu}

\begin{exercice}[À vous de jouer : comparer et argumenter]
\begin{enumerate}
\item Traduisez en chinois : « Le droit écrit est plus clair que le droit coutumier. » (Voc. : \zh{成文法}, \zh{习惯法}, \zh{清楚}).
\item Traduisez : « Au Cameroun comme en Chine, les aînés sont respectés. » (Utilisez \zh{老人} et \zh{受尊重} ou \zh{很重要} ; structure \zh{和……一样} + \zh{都}).
\item Corrigez la phrase fautive (erreur de mélange) : \zh{中国比喀麦隆一样大。} Donnez les deux versions correctes possibles (différence OU égalité).
\item \textbf{Production écrite (Mini-exposé)} : Rédigez un avis complet en 4 étapes (Opinion, Raison, Exemple, Conclusion) sur le sujet : \textbf{« Faut-il que les filles héritent comme les fils ? »} Utilisez \zh{我认为}, \zh{因为}, \zh{比如} / \zh{在……}, \zh{所以} et \zh{和……一样}. Écrivez en caractères, pinyin et français.
\end{enumerate}
\end{exercice}

\begin{corrige}[Exercice : Comparer et donner son avis]
\begin{enumerate}
\item \zh{成文法比习惯法更清楚。} \emph{Chéngwénfǎ bǐ xíguànfǎ gèng qīngchu.} — \textit{Note : pas de \zh{很} après \zh{比} ; \zh{更} renforce le comparatif.}
\item \zh{喀麦隆和中国一样，老人都很受尊重。} \emph{Kāmàilóng hé Zhōngguó yíyàng, lǎorén dōu hěn shòu zūnzhòng.} — \textit{Structure : A 和 B 一样, Sujet 都 Adj. \zh{都} insiste sur le fait que les deux groupes respectent les aînés.}
\item \textbf{Phrase fautive} : \zh{中国比喀麦隆一样大。} (Mélange \zh{比} + \zh{一样}).
\\[2pt]
\textbf{Correction A (Différence)} : \zh{中国比喀麦隆大。} \emph{Zhōngguó bǐ Kāmàilóng dà.} (La Chine est plus grande).
\\[2pt]
\textbf{Correction B (Égalité)} : \zh{中国和喀麦隆一样大。} \emph{Zhōngguó hé Kāmàilóng yíyàng dà.} (La Chine est aussi grande que le Cameroun — fait géographiquement faux, mais grammaticalement correct).
\item \textbf{Modèle de production (un exemple parmi d'autres) :}
\\[4pt]
\textbf{Opinion} : \zh{我认为女儿应该和儿子一样继承财产。} \emph{Wǒ rènwéi nǚ'ér yīnggāi hé érzi yíyàng jìchéng cáichǎn.}
\\[4pt]
\textbf{Raison} : \zh{因为女儿也是家庭的一员，而且女儿往往照顾父母更多。} \emph{Yīnwèi nǚ'ér yě shì jiātíng de yì yuán, érqiě nǚ'ér wǎngwǎng zhàogù fùmǔ gèng duō.}
\\[4pt]
\textbf{Exemple} : \zh{比如在我家，姐姐照顾父母十几年了。} \emph{Bǐrú zài wǒ jiā, jiějie zhàogù fùmǔ shí jǐ nián le.}
\\[4pt]
\textbf{Conclusion} : \zh{所以，法律应该保障女儿和儿子一样的权利。} \emph{Suǒyǐ, fǎlǜ yīnggāi bǎozhàng nǚ'ér hé érzi yíyàng de quánlì.}
\\[4pt]
\textbf{Texte continu pour l'oral (environ 45--50 secondes) :}
\zh{我认为女儿应该和儿子一样继承财产，因为女儿也是家庭的一员，而且女儿往往照顾父母更多。比如在我家，姐姐照顾父母十几年了，所以法律应该保障女儿和儿子一样的权利。}
\end{enumerate}
\end{corrige}

\begin{aretenir}[L'essentiel pour l'exposé comparatif (Leçon 8 -- Section 5)]
\begin{itemize}
\item \textbf{Comparer (Différence)} : \textbf{A 比 B Adj} — \ding{55} \zh{很}, \ding{55} mot pour « que ». Renfort : \zh{更}, \zh{得多}. Négation : \zh{A 没有 B (那么) Adj}.
\item \textbf{Comparer (Égalité)} : \textbf{A 和 B 一样 (Adj)} — \ding{55} oublier \zh{一样}. \ding{55} mélanger avec \zh{比}. Variante orale : \zh{跟}.
\item \textbf{Opiner} : \zh{我觉得} (oral/subjectif) vs \zh{我认为} (écrit/argumenté/exposé). Suivi direct d'une phrase complète.
\item \textbf{Argumenter} : \zh{因为……所以……} (paire inséparable), \zh{但是} (mais fort), \zh{不过} (cependant, nuance), \zh{而且} (de plus, en milieu de phrase).
\item \textbf{Méthode 4 étapes} : \zh{我认为...} $\rightarrow$ \zh{因为...} $\rightarrow$ \zh{比如/在...} $\rightarrow$ \zh{所以...和...一样...}.
\end{itemize}
Vous êtes maintenant prêt(e) pour la dernière section : \textbf{Questions de réflexion et préparation du mini-exposé noté}.
\end{aretenir}

\coursec{Questions de réflexion et préparation du mini-exposé}

Dans la section précédente, nous avons appris à comparer les systèmes d'héritage du Cameroun et de la Chine avec les structures \zh{和……一样}, \zh{比} et \zh{跟……不一样}, et à donner un avis simple avec \zh{我觉得} ou \zh{我认为}. Il est temps maintenant de mettre tout cela en pratique : d'abord en réfléchissant ensemble à trois grandes questions, puis en préparant un \cle{mini-exposé} oral de huit phrases, exactement comme on vous le demandera au baccalauréat.

\courssub{Trois questions de réflexion}

Lisez chaque question à voix haute, repérez les mots connus, puis essayez d'y répondre en une ou deux phrases simples. L'objectif n'est pas de trouver « la » bonne réponse unique, mais de construire une phrase chinoise correcte et claire.

\textbf{Question 1 : 喀麦隆和中国的继承方法有什么一样的地方？}

\emph{Kāmàilóng hé Zhōngguó de jìchéng fāngfǎ yǒu shénme yíyàng de dìfang ?}

« Quels points communs y a-t-il entre les méthodes d'héritage du Cameroun et de la Chine ? »

Pour répondre, on réutilise la structure de comparaison vue dans la section 5 : \textbf{A 和 B 一样 + adjectif / verbe}. Exemple de réponse possible :

\zh{两个国家都有继承的法律，也都尊重家庭的传统。}

\emph{Liǎng ge guójiā dōu yǒu jìchéng de fǎlǜ, yě dōu zūnzhòng jiātíng de chuántǒng.}

« Les deux pays ont des lois sur l'héritage et respectent aussi les traditions familiales. »

\textbf{Question 2 : 女儿应该继承财产吗？为什么？}

\emph{Nǚ'ér yīnggāi jìchéng cáichǎn ma ? Wèishénme ?}

« Les filles devraient-elles hériter des biens ? Pourquoi ? »

C'est une question d'opinion. On répond avec \zh{我觉得} + phrase, puis on justifie avec \zh{因为}. Exemple :

\zh{我觉得女儿应该继承财产，因为儿子和女儿都是父母的孩子。}

\emph{Wǒ juéde nǚ'ér yīnggāi jìchéng cáichǎn, yīnwèi érzi hé nǚ'ér dōu shì fùmǔ de háizi.}

« Je pense que les filles devraient hériter, parce que les fils et les filles sont tous les enfants de leurs parents. »

\textbf{Question 3 : 传统和法律，哪个更重要？}

\emph{Chuántǒng hé fǎlǜ, nǎ ge gèng zhòngyào ?}

« La tradition et la loi, laquelle est plus importante ? »

On répond avec \zh{我觉得 + A 比 B 更 + adjectif} ou avec une réponse nuancée (\zh{两个都重要}, « les deux sont importants »). Exemple :

\zh{我觉得法律更重要，因为法律保护每一个人。}

\emph{Wǒ juéde fǎlǜ gèng zhòngyào, yīnwèi fǎlǜ bǎohù měi yí ge rén.}

« Je pense que la loi est plus importante, parce que la loi protège chaque personne. »

\begin{center}
\begin{tabularx}{\linewidth}{|X|X|X|X|}
\hline
\rowcolor{popblueL} Caractères & Pinyin & Nature & Traduction \\
\hline
继承 & jìchéng & verbe & hériter \\
\hline
方法 & fāngfǎ & nom & méthode, manière \\
\hline
一样 & yíyàng & adjectif & pareil, identique \\
\hline
地方 & dìfang & nom & point, aspect, endroit \\
\hline
应该 & yīnggāi & verbe modal & devoir (il faudrait) \\
\hline
财产 & cáichǎn & nom & biens, patrimoine \\
\hline
传统 & chuántǒng & nom & tradition \\
\hline
法律 & fǎlǜ & nom & loi, législation \\
\hline
重要 & zhòngyào & adjectif & important \\
\hline
保护 & bǎohù & verbe & protéger \\
\hline
看法 & kànfǎ & nom & point de vue, avis \\
\hline
尊重 & zūnzhòng & verbe & respecter \\
\hline
\end{tabularx}
\end{center}

\begin{attention}[看法 se prononce kànfǎ, pas kānfǎ]
Le caractère \zh{看} a deux prononciations. Dans \zh{看法} (\emph{kànfǎ}, « point de vue »), \zh{看见} (\emph{kànjiàn}, « voir ») et \zh{看书} (\emph{kàn shū}, « lire »), il se prononce \textbf{kàn} (4\up{e} ton). La prononciation \textbf{kān} (1\ier{} ton) existe dans \zh{看家} (\emph{kān jiā}, « garder la maison »), mot rare à votre niveau. À l'oral du baccalauréat, dire \emph{kānfǎ} pour « mon avis » est une erreur de ton très fréquente chez les francophones : retenez \zh{我的看法} = \emph{wǒ de kànfǎ}.
\end{attention}

\courssub{La grille de préparation de l'exposé}

Un bon exposé n'est pas une improvisation : c'est une construction. Voici la grille officielle en cinq étapes, pour un total de huit phrases. C'est exactement le nombre attendu pour un mini-exposé de deux minutes.

\begin{enumerate}
\item \textbf{Introduction} (1 phrase) : annoncer le sujet avec \zh{今天我说一说……} (« Aujourd'hui, je vais parler de… »).
\item \textbf{Le Cameroun} (2 phrases) : présenter la situation camerounaise.
\item \textbf{La Chine} (2 phrases) : présenter la situation chinoise.
\item \textbf{La comparaison} (2 phrases) : un point commun avec \zh{一样}, une différence avec \zh{比} ou \zh{不一样}.
\item \textbf{Avis personnel} (1 phrase) : \zh{我的看法是……} ou \zh{我觉得……}.
\end{enumerate}

\begin{popfigure}
\begin{tikzpicture}[
  node distance=0.5cm,
  etape/.style={rectangle, rounded corners=3pt, draw=popblue, fill=popblueL, text width=8.5cm, align=center, inner sep=5pt, font=\small},
  etapeLast/.style={rectangle, rounded corners=3pt, draw=poporange, fill=poporangeL, text width=8.5cm, align=center, inner sep=5pt, font=\small},
  fleche/.style={-{Stealth[length=2.5mm]}, thick, poporange}
]
\node[etape] (e1) {\textbf{1. Introduction} — 今天我说一说继承的问题。};
\node[etape, below=of e1] (e2) {\textbf{2. Cameroun} — 在喀麦隆，传统很重要。很多家庭有习惯法。};
\node[etape, below=of e2] (e3) {\textbf{3. Chine} — 在中国，继承法保护儿子和女儿。};
\node[etape, below=of e3] (e4) {\textbf{4. Comparaison} — 两个国家都有法律，但是传统不一样。};
\node[etapeLast, below=of e4] (e5) {\textbf{5. Avis personnel} — 我的看法是：法律和传统都重要。};
\draw[fleche] (e1) -- (e2);
\draw[fleche] (e2) -- (e3);
\draw[fleche] (e3) -- (e4);
\draw[fleche] (e4) -- (e5);
\end{tikzpicture}
\legende{Organigramme de la structure de l'exposé en cinq étapes : chaque étape s'appuie sur la précédente.}
\end{popfigure}

\courssub{Ordonner son exposé : 首先 / 然后 / 最后}

Pour que votre exposé soit clair, il faut guider l'auditeur avec des \cle{mots de liaison} qui annoncent l'ordre des idées, comme « d'abord, ensuite, enfin » en français.

\begin{propriete}[Les marqueurs d'ordre]
\begin{itemize}
\item \zh{首先} (\emph{shǒuxiān}) = « d'abord, en premier lieu » : ouvre l'exposé.
\item \zh{然后} (\emph{ránhòu}) = « ensuite, puis » : enchaîne les parties (on peut le répéter).
\item \zh{最后} (\emph{zuìhòu}) = « enfin, pour finir » : introduit la conclusion ou l'avis personnel.
\end{itemize}
Structure : \textbf{首先，……。然后，……。最后，……。} — la virgule chinoise \zh{，} suit toujours le marqueur.
\end{propriete}

\begin{exemplebox}[Un squelette d'exposé complet]
\zh{首先，我说一说喀麦隆的情况。然后，我介绍中国的继承法。最后，我说说我的看法。}

\emph{Shǒuxiān, wǒ shuō yi shuō Kāmàilóng de qíngkuàng. Ránhòu, wǒ jièshào Zhōngguó de jìchéngfǎ. Zuìhòu, wǒ shuōshuo wǒ de kànfǎ.}

« D'abord, je parle de la situation du Cameroun. Ensuite, je présente la loi chinoise sur l'héritage. Enfin, je donne mon avis. »
\end{exemplebox}

Remarquez la forme \zh{说一说} et \zh{说说} : la répétition du verbe adoucit le ton (« parler un peu de »), c'est très naturel à l'oral. Notez aussi \zh{继承法} (\emph{jìchéngfǎ}) : « loi sur l'héritage », mot composé de \zh{继承} + \zh{法} (loi).

\begin{methode}[Réussir son mini-exposé en cinq conseils]
\begin{enumerate}
\item \textbf{Écrivez d'abord vos huit phrases} sur papier, en suivant la grille, puis vérifiez chaque ton.
\item \textbf{Parlez lentement} : en chinois, la clarté des tons compte plus que la vitesse. Une phrase lente et juste vaut mieux que trois phrases rapides et fausses.
\item \textbf{Ordonnez avec 首先 / 然后 / 最后} : l'examinateur entend immédiatement que votre exposé est structuré.
\item \textbf{Regardez votre auditoire}, pas votre feuille : apprenez votre plan par cœur, gardez seulement les mots-clés sous les yeux.
\item \textbf{Terminez par votre avis} avec \zh{我的看法是……} : c'est la phrase qui montre que vous savez « exprimer un point de vue simple en chinois », compétence exigée par le programme.
\end{enumerate}
\end{methode}

\courssub{Entraînement par binômes}

\begin{experience}[Présenter et interroger]
Formez des binômes. L'élève A prépare son exposé de huit phrases avec la grille et le présente en deux minutes, debout, sans lire. L'élève B écoute, puis pose \textbf{une question} en chinois, choisie parmi les trois questions de réflexion ou inventée : \zh{女儿应该继承财产吗？} / \zh{传统和法律，哪个更重要？} L'élève A répond en une phrase avec \zh{我觉得……因为……}. Puis on échange les rôles. Chronométrez-vous !
\end{experience}

\courssub{Critères d'évaluation}

Votre exposé sera noté selon quatre critères, comme à l'examen. Gardez-les en tête pendant la préparation.

\begin{center}
\begin{tabularx}{\linewidth}{|X|X|X|}
\hline
\rowcolor{popblueL} Critère & Ce qu'on attend & Exemple de réussite \\
\hline
Exactitude du vocabulaire & employer les mots du thème : 继承， 财产， 法律， 传统 & 在中国，继承法保护女儿。 \\
\hline
Structures de comparaison & au moins une structure 一样 ou 比 & 两个国家都有继承的法律。 \\
\hline
Respect des tons & tons justes, surtout kànfǎ, nǚ'ér, fǎlǜ & 我的看法 (wǒ de kànfǎ) \\
\hline
Clarté et ordre & parler lentement, utiliser 首先 / 然后 / 最后 & 首先……然后……最后…… \\
\hline
\end{tabularx}
\end{center}

\begin{exoresolu}[Construire la phrase de comparaison d'un exposé]
Énoncé : vous êtes à l'étape 4 de votre exposé. Dites en une phrase que le Cameroun et la Chine ont tous les deux des lois sur l'héritage, puis en une phrase que les traditions sont différentes.
\tcblower
Solution détaillée : pour le point commun, on utilise \zh{A 和 B 一样} ou plus simplement \zh{两个国家都有……} : \zh{喀麦隆和中国都有继承的法律。} (\emph{Kāmàilóng hé Zhōngguó dōu yǒu jìchéng de fǎlǜ.}) « Le Cameroun et la Chine ont tous les deux des lois sur l'héritage. » Pour la différence, on utilise \zh{但是} (« mais ») ou \zh{不一样} : \zh{但是，两个国家的传统不一样。} (\emph{Dànshì, liǎng ge guójiā de chuántǒng bù yíyàng.}) « Mais les traditions des deux pays sont différentes. » Piège évité : ne dites pas \zh{两个国家都有有} (répétition) et n'oubliez pas le classificateur \zh{个} dans \zh{两个国家}.
\end{exoresolu}

\begin{attention}[Erreurs fréquentes des francophones à l'oral]
\begin{itemize}
\item Traduire « je pense que » par \zh{我想} : non, \zh{我想} signifie « je veux / je pense à ». Pour donner un avis, dites \zh{我觉得} ou \zh{我认为}.
\item Oublier \zh{的} dans \zh{我的看法} : on ne dit pas \zh{我看法}.
\item Placer \zh{因为} comme en français en fin de phrase sans sujet répété : en chinois, on peut dire \zh{因为……，所以……} (« parce que…, donc… »), structure impossible en français mais correcte et élégante en chinois : \zh{因为法律保护每个人，所以法律很重要。}
\end{itemize}
\end{attention}

\begin{savaistu}[Cette grille est votre passeport pour le Bac]
Au baccalauréat, les questions socioculturelles demandent très souvent de comparer le Cameroun et la Chine en trois à cinq phrases : la famille, l'école, les fêtes, la nourriture, l'héritage… La grille que vous venez d'apprendre (introduction → Cameroun → Chine → comparaison → avis) est directement réutilisable pour n'importe quel sujet. Apprenez-la une fois, elle vous servira toute l'année — et les marqueurs \zh{首先 / 然后 / 最后} impressionnent toujours un jury.
\end{savaistu}

\begin{exercice}[À vous de jouer]
\begin{enumerate}
\item Répondez par écrit, en une phrase chinoise complète, à la question : \zh{女儿应该继承财产吗？为什么？}
\item Rédigez votre mini-exposé complet de huit phrases en suivant la grille (introduction, Cameroun, Chine, comparaison, avis). Soulignez les marqueurs \zh{首先}, \zh{然后}, \zh{最后}.
\item Préparez deux questions à poser à votre binôme sur le thème de l'héritage, avec \zh{吗} ou \zh{哪个}.
\end{enumerate}
\end{exercice}

\begin{corrige}[Exercice « À vous de jouer » — éléments de réponse]
\begin{enumerate}
\item Réponse possible : \zh{我觉得女儿应该继承财产，因为女儿也是父母的孩子。} (\emph{Wǒ juéde nǚ'ér yīnggāi jìchéng cáichǎn, yīnwèi nǚ'ér yě shì fùmǔ de háizi.}) Toute phrase avec \zh{我觉得 / 我认为} + \zh{因为} est acceptable si les tons et l'ordre des mots sont corrects.
\item Modèle : \zh{首先，今天我说一说继承的问题。在喀麦隆，很多家庭尊重传统。习惯法很重要。然后，在中国，继承法保护儿子和女儿。城市和农村都使用这个法律。两个国家都有继承的法律，但是传统不一样。最后，我的看法是：法律和传统都重要。}
\item Exemples : \zh{传统和法律，哪个更重要？} / \zh{你的国家有继承法吗？}
\end{enumerate}
\end{corrige}

\begin{aretenir}[L'essentiel de la section]
\begin{itemize}
\item Trois questions types : points communs (\zh{有什么一样的地方？}), opinion (\zh{应该……吗？为什么？}), choix (\zh{哪个更重要？}).
\item Grille de l'exposé : introduction (1) → Cameroun (2) → Chine (2) → comparaison (2) → avis (1) = huit phrases.
\item Ordonner avec \zh{首先 / 然后 / 最后} ; donner son avis avec \zh{我的看法是……} (\emph{kànfǎ}, 4\up{e} ton !).
\item Évaluation : vocabulaire exact, structures de comparaison, tons respectés, exposé clair et lent.
\end{itemize}
\end{aretenir}

\newpage

\coursec{Activité d'intégration}

\courssub{Objectifs de l'activité}
\noindent Cette activité d'intégration vise à mobiliser l'ensemble des acquis de la leçon (vocabulaire socioculturel, structures de comparaison, expression de l'opinion) dans une tâche finale authentique : la réalisation d'un \textbf{mini-exposé comparatif oral} (et sa variante écrite). Les élèves devront :
\begin{itemize}
    \item Restituer de façon structurée des faits culturels sur l'hérédité au Cameroun et en Chine.
    \item Employer correctement les structures de comparaison (\zh{比}, \zh{跟\ldots{}一样}, \zh{没有\ldots{}那么}, \zh{相比之下}) et les marqueurs de relation logique (\zh{首先}, \zh{但是}, \zh{因此}).
    \item Exprimer un point de vue personnel nuancé et respectueux (\zh{我觉得}, \zh{我认为}, \zh{应该}).
    \item Travailler la cohésion de groupe (répartition des rôles, écoute active, évaluation par les pairs).
\end{itemize}

\courssub{Situation}
\noindent \textbf{Contexte :} Dans le cadre d'un jumelage numérique entre le \textbf{Lycée Général Leclerc de Yaoundé} et le \textbf{Lycée n°4 de Pékin}, une visioconférence thématique est organisée sur le thème \og \textit{Transmission du patrimoine familial : regards croisés Cameroun-Chine} \fg{}. Chaque établissement doit présenter un exposé collaboratif de 6 minutes (3 élèves $\times$ 2 minutes) suivi d'un échange de 10 minutes. Vous faites partie de la délégation camerounaise. Votre groupe de trois doit préparer la partie \og \textit{L'héritage chez moi et en Chine} \fg{}.

\noindent Vous recevez un message de votre correspondant chinois, Li Wei, qui partage les grandes lignes de l'exposé chinois pour que vous puissiez préparer la comparaison.

\begin{textebox}[Message de Li Wei (WeChat) -- Contexte de l'échange]
李伟：你好！我们下周三视频会议的主题是“家族遗产传承”。我们中国组准备介绍：中国传统上遵循“长子继承制”（传给长子），但现在法律规定男女平等继承（继承法）。农村还是倾向于把房子、土地留给儿子，城市则更平均。还有“宗祠”、家谱文化。请你们准备喀麦隆的情况，我们到时对比。期待交流！
\\ \textit{Lǐ Wěi: Nǐ hǎo! Wǒmen xià zhōusān shìpín huìyì de zhǔtí shì "Jiāzú yíchán chuánchéng". Wǒmen Zhōngguó zǔ zhǔnbèi jièshào: Zhōngguó chuántǒng shàng zūnxún "zhǎngzǐ jìchéngzhì" (chuán gěi zhǎngzǐ), dàn xiànzài fǎlǜ guīdìng nánnǚ píngděng jìchéng (Jìchéngfǎ). Nóngcūn háishì qīngxiàngyú bǎ fángzi, tǔdì liú gěi érzi, chéngshì zé gèng píngjūn. Hái yǒu "zōngcí", jiāpǔ wénhuà. Qǐng nǐmen zhǔnbèi Kāmàilóng de qíngkuàng, wǒmen dào shí duìbǐ. Qídài jiāoliú!}
\\ \textit{Traduction : Li Wei : Salut ! Le thème de notre visioconférence mercredi prochain est "Transmission du patrimoine familial". Notre groupe chinois va présenter : traditionnellement la Chine suit le "système d'héritage par l'aîné" (transmis au fils aîné), mais la loi actuelle stipule l'égalité homme-femme dans l'héritage (Loi sur les successions). En zone rurale, on tend encore à laisser maison et terre aux fils, en ville c'est plus égalitaire. Il y a aussi la culture des "temples ancestraux" et des registres généalogiques. Préparez la situation du Cameroun, nous comparerons alors. Impatients d'échanger !}
\end{textebox}

\courssub{Consignes}
\noindent \textbf{Travail en groupe de trois (rôles tournants). Durée totale : 1h30 (préparation 45 min, passage 15 min, écoute/vote 30 min).}

\begin{enumerate}
    \item \textbf{Analyse du document source (10 min) :} Lisez le message de Li Wei. Identifiez les \cle{mots-clés} chinois (vocabulaire HSK 3-4) : \zh{长子继承制}, \zh{继承法}, \zh{男女平等}, \zh{倾向于}, \zh{宗祠}, \zh{家谱}. Notez leur sens en français. Relevez deux faits sur la Chine (un traditionnel, un moderne) et une différence ville/campagne.
    \item \textbf{Collecte \& structuration des faits Cameroun (15 min) :} À partir de la leçon 4 et de votre connaissance, listez 3 à 4 faits sur l'hérédité au Cameroun (ex: droit coutumier vs droit écrit, rôle de l'oncle maternel chez certains peuples, partage des terres, place des filles). Rédigez-les en chinois simple (caractères + pinyin) sur une fiche brouillon. Utilisez le vocabulaire du tableau de la leçon (\zh{习俗}, \zh{分配}, \zh{长子}, \zh{女儿}, \zh{土地}, \zh{法律}).
    \item \textbf{Rédaction du script collaboratif (20 min) :} Répartissez les 5 étapes de l'exposé (voir grille ci-dessous). Chaque élève rédige sa partie (environ 40 secondes / 5-6 phrases). Vérifiez ensemble : tons, ordre des mots, connecteurs logiques, respect des structures de comparaison. \textbf{Contrainte :} Chaque partie doit contenir au moins \textbf{une structure de comparaison} et \textbf{un connecteur logique}.
    \item \textbf{Répétition \& enregistrement (15 min) :} Répétez à voix haute en groupe. Chronométrez (2 min par élève = 6 min total). Enregistrez-vous (dictaphone/téléphone) pour auto-évaluation de la prononciation (tons, \zh{ü}, \zh{zh/ch/sh/r}).
    \item \textbf{Passage devant la classe \& Écoute active (20 min) :} Chaque groupe présente. Les élèves auditeurs remplissent la \textbf{Fiche d'écoute} (voir ci-dessous).
    \item \textbf{Vote \& Retour méta-cognitif (10 min) :} Vote à bulletin secret pour l'exposé \og Le plus clair \fg{} et \og Le plus respectueux des deux cultures \fg{}. Debriefing oral : \og Quelle structure vous a le plus aidé ? Quel mot de vocabulaire manquait ? \fg{}.
\end{enumerate}

\noindent \textbf{Grille de l'exposé (5 étapes -- 1 élève par étape) :}
\begin{center}
\begin{tabularx}{\linewidth}{|X|X|X|}\hline
\rowcolor{popblueL} \textbf{Étape} & \textbf{Contenu attendu} & \textbf{Structures cibles obligatoires} \\ \hline
1. Introduction & Présenter le thème, le contexte (visioconférence), annoncer le plan. & \zh{今天我们要介绍\ldots{}；首先\ldots{}，其次\ldots{}，最后\ldots{}} \\ \hline
2. Héritage au Cameroun & 2-3 faits clés (coutume, droit, famille élargie, genre). & \zh{在喀麦隆\ldots{}；根据习俗\ldots{}；但是根据法律\ldots{}} \\ \hline
3. Héritage en Chine (reprise Li Wei) & Résumer les points de Li Wei (tradition, loi, ville/campagne, culture). & \zh{在中国\ldots{}；传统上\ldots{}；现在\ldots{}；农村\ldots{}，城市\ldots{}} \\ \hline
4. Comparaison explicite & 1 similitude + 1 différence majeure + 1 nuance (évolution). & \zh{跟中国一样\ldots{}；跟中国不同\ldots{}；相比之下\ldots{}；没有\ldots{}那么\ldots{}} \\ \hline
5. Avis personnel + Conclusion & Opinion nuancée (valeurs, modernité), remerciements, ouverture. & \zh{我觉得\ldots{}；我认为\ldots{}应该\ldots{}；谢谢大家！} \\ \hline
\end{tabularx}
\end{center}

\noindent \textbf{Fiche d'écoute (pour les élèves auditeurs) -- À remplir pendant chaque exposé :}
\begin{center}
\begin{tabularx}{\linewidth}{|X|X|}\hline
\rowcolor{popgreenL} \textbf{Élément à repérer} & \textbf{Votre note (Mots-clés en chinois / français)} \\ \hline
Un point commun Cameroun -- Chine & \\ \hline
Une différence marquante & \\ \hline
Une structure de comparaison entendue & \\ \hline
Une question à poser aux présentateurs & \\ \hline
Note globale : Clarté /10 -- Respect culturel /10 -- Chinois /10 & \\ \hline
\end{tabularx}
\end{center}

\courssub{Ressources à mobiliser}
\noindent \textbf{1. Vocabulaire thématique (révision + extension)}

\begin{center}
\begin{tabularx}{\linewidth}{|X|X|X|X|}\hline
\rowcolor{popblueL} \textbf{Caractères} & \textbf{Pinyin} & \textbf{Nature} & \textbf{Traduction} \\ \hline
遗产 & yíchán & n. & héritage, patrimoine \\ \hline
继承 & jìchéng & v. & hériter, succéder \\ \hline
继承法 & jìchéngfǎ & n. & loi sur les successions / code successoral \\ \hline
长子继承制 & zhǎngzǐ jìchéngzhì & n. & système d'héritage par l'aîné (primogéniture) \\ \hline
男女平等 & nánnǚ píngděng & adj./n. & égalité homme-femme \\ \hline
习俗 & xísú & n. & coutume, mœurs \\ \hline
法律 & fǎlǜ & n. & loi, législation \\ \hline
分配 & fēnpèi & v. & répartir, distribuer, partager \\ \hline
倾向于 & qīngxiàngyú & v. & tendre à, pencher vers \\ \hline
宗祠 & zōngcí & n. & temple ancestral, salle des ancêtres \\ \hline
家谱 & jiāpǔ & n. & registre généalogique, arbre généalogique \\ \hline
长子 / 次子 & zhǎngzǐ / cìzǐ & n. & fils aîné / cadet \\ \hline
女儿 / 儿子 & nǚ'ér / érzi & n. & fille / fils \\ \hline
土地 & tǔdì & n. & terre, terrain \\ \hline
房子 & fángzi & n. & maison \\ \hline
相比之下 & xiāngbǐ zhīxià & adv. & en comparaison, par contraste \\ \hline
既\ldots{}又\ldots{} & jì\ldots{}yòu\ldots{} & conj. & à la fois\ldots{}et\ldots{} \\ \hline
\end{tabularx}
\end{center}

\noindent \textbf{2. Structures grammaticales clés (fiches-mémo)}

\begin{propriete}[Comparaison : \texorpdfstring{\zh{比}}{bi} (supériorité) vs \texorpdfstring{\zh{没有\ldots{}那么/这么}}{meiyou name/zheme} (infériorité/égalité négative)]
\noindent \textbf{Structure A (Supériorité) :} \zh{Sujet A + 比 + Sujet B + Adjectif / Verbe d'état} \\
\zh{喀麦隆的习俗比中国的习俗复杂。} (Kāmàilóng de xísú bǐ Zhōngguó de xísú fùzá.) -- Les coutumes camerounaises sont plus complexes que les chinoises. \\
\textbf{Structure B (Infériorité/Négation de l'égalité) :} \zh{Sujet A + 没有 + Sujet B + 那么/这么 + Adjectif} \\
\zh{中国的农村传统没有以前那么强。} (Zhōngguó de nóngcūn chuántǒng méiyǒu yǐqián nàme qiáng.) -- La tradition rurale chinoise n'est plus aussi forte qu'avant. \\
\textbf{Attention :} Ne pas dire \zh{比\ldots{}不} pour l'infériorité. \zh{不比} existe mais change le sens (n'est pas comparable / ne fait pas le poids).
\end{propriete}

\begin{propriete}[Comparaison : \texorpdfstring{\zh{跟\ldots{}一样/不一样}}{gen yiyang/buyiyang} (similitude/différence) et \texorpdfstring{\zh{相比之下}}{xiangbi zhixia} (transition comparative)]
\noindent \textbf{Similitude :} \zh{Sujet A + 跟 + Sujet B + 一样 + Adjectif / Verbe} \\
\zh{喀麦隆的法律跟中国的法律一样规定男女平等。} (Kāmàilóng de fǎlǜ gēn Zhōngguó de fǎlǜ yīyáng guīdìng nánnǚ píngděng.) -- La loi camerounaise, comme la loi chinoise, stipule l'égalité homme-femme. \\
\textbf{Différence :} \zh{Sujet A + 跟 + Sujet B + 不一样} (Sujet A diffère de Sujet B) \\
\zh{喀麦隆的分配方式跟中国不一样。} (Kāmàilóng de fēnpèi fāngshì gēn Zhōngguó bù yīyáng.) -- Le mode de partage au Cameroun diffère de la Chine. \\
\textbf{Transition :} \zh{相比之下} (Par comparaison / En revanche) s'utilise en tête de phrase pour opposer deux constats. \\
\zh{中国重视家谱。相比之下，喀麦隆更重视口头传承。} (Zhōngguó zhòngshì jiāpǔ. Xiāngbǐ zhīxià, Kāmàilóng gèng zhòngshì kǒutóu chuánchéng.)
\end{propriete}

\begin{propriete}[Expression de l'opinion et de la recommandation : \texorpdfstring{\zh{觉得 / 认为 / 应该}}{jue de / ren wei / ying gai}]
\noindent \zh{我觉得 (juéde)} : ressenti subjectif, plus oral. \zh{我认为 (rènwéi)} : opinion réfléchie, plus formel (exposé). \\
\zh{应该 (yīnggāi)} : devoir moral, recommandation (se place \textbf{avant} le verbe). \\
\zh{我觉得保护传统很重要。} (Wǒ juéde bǎohù chuántǒng hěn zhòngyào.) -- Je trouve que protéger la tradition est important. \\
\zh{我认为我们应该尊重法律，也尊重习俗。} (Wǒ rènwéi wǒmen yīnggāi zūnzhòng fǎlǜ, yě zūnzhòng xísú.) -- Je pense qu'on doit respecter la loi, mais aussi les coutumes.
\end{propriete}

\begin{propriete}[Marqueurs de progression (Discours structuré)]
\noindent \zh{首先 (shǒuxiān)} -- Premièrement ; \zh{其次 (qícì)} -- Deuxièmement ; \zh{然后 (ránhòu)} -- Ensuite ; \zh{最后 (zuìhòu)} -- Enfin. \\
\zh{但是 (dànshì) / 不过 (búguò)} -- Mais (néanmoins). \\
\zh{因为\ldots{}所以\ldots{} (yīnwèi\ldots{}suǒyǐ\ldots{})} -- Parce que\ldots{}donc\ldots{} (cause/conséquence). \\
\zh{不但\ldots{}而且\ldots{} (búdàn\ldots{}érqiě\ldots{})} -- Non seulement\ldots{}mais encore\ldots{}.
\end{propriete}

\courssub{Proposition de production et corrigé commenté}
\noindent Voici une proposition de script complet pour le groupe (Élève A, B, C) et sa version écrite condensée (variante cahier).

\begin{exoresolu}[Mini-exposé modèle : « L'héritage chez moi et en Chine » -- Groupe 3 (Élèves A, B, C)]
\textbf{ÉLÈVE A (Introduction + Cameroun) :}
\\ \zh{大家好！今天我们要介绍“喀麦隆和中国的遗产传承”。首先，介绍喀麦隆的情况。}
\\ \textit{Dàjiā hǎo! Jīntiān wǒmen yào jièshào "Kāmàilóng hé Zhōngguó de yíchán chuánchéng". Shǒuxiān, jièshào Kāmàilóng de qíngkuàng.}
\\ \zh{在喀麦隆，习俗和法律有时不一样。根据传统习俗，在很多民族里，土地和房子传给长子，或者传给舅舅的儿子（姐妹的儿子）。女儿通常不分土地。}
\\ \textit{Zài Kāmàilóng, xísú hé fǎlǜ yǒushí bù yīyáng. Gēnjù chuántǒng xísú, zài hěnduō mínzú lǐ, tǔdì hé fángzi chuán gěi zhǎngzǐ, huòzhě chuán gěi jiùjiu de érzi (jiěmèi de érzi). Nǚ'ér tōngcháng bù fēn tǔdì.}
\\ \zh{但是，根据现代法律（民法典），男女平等继承，女儿也有权分财产。所以，现在城里家庭常常平均分配。}
\\ \textit{Dànshì, gēnjù xiàndài fǎlǜ (mínfǎdiǎn), nánnǚ píngděng jìchéng, nǚ'ér yě yǒu quán fēn cáichǎn. Suǒyǐ, xiànzài chénglǐ jiātíng chángcháng píngjūn fēnpèi.}

\textbf{ÉLÈVE B (Chine d'après Li Wei + Transition) :}
\\ \zh{其次，介绍中国的情况，根据李伟的信息。}
\\ \textit{Qícì, jièshào Zhōngguó de qíngkuàng, gēnjù Lǐ Wěi de xìnxī.}
\\ \zh{中国传统上有“长子继承制”，家产传给长子。现在《继承法》规定男女平等，儿子女儿都继承。}
\\ \textit{Zhōngguó chuántǒng shàng yǒu "zhǎngzǐ jìchéngzhì", jiāchǎn chuán gěi zhǎngzǐ. Xiànzài "Jìchéngfǎ" guīdìng nánnǚ píngděng, érzi nǚ'ér dōu jìchéng.}
\\ \zh{农村倾向于把房子留给儿子，城市比较平均。中国还有宗祠和家谱文化，很重视家族历史。}
\\ \textit{Nóngcūn qīngxiàngyú bǎ fángzi liú gěi érzi, chéngshì bǐjiào píngjūn. Zhōngguó hái yǒu zōngcí hé jiāpǔ wénhuà, hěn zhòngshì jiāzú lìshǐ.}

\textbf{ÉLÈVE C (Comparaison + Avis personnel + Conclusion) :}
\\ \zh{然后，我们比较一下。}
\\ \textit{Ránhòu, wǒmen bǐjiào yíxià.}
\\ \zh{跟中国一样，喀麦隆现在法律也规定男女平等，这是一个共同点。}
\\ \textit{Gēn Zhōngguó yīyáng, Kāmàilóng xiànzài fǎlǜ yě guīdìng nánnǚ píngděng, zhè shì yí gè gòngtóngdiǎn.}
\\ \zh{跟中国不同，喀麦隆习俗里“舅舅继承”（姐妹的儿子继承）很特别，中国没有这个习俗。}
\\ \textit{Gēn Zhōngguó bù yīyáng, Kāmàilóng xísú lǐ "jiùjiu jìchéng" (jiěmèi de érzi jìchéng) hěn tèbié, Zhōngguó méiyǒu zhège xísú.}
\\ \zh{相比之下，中国更重视家谱和宗祠（书面记录），喀麦隆更重视口头传承（长者讲故事）。}
\\ \textit{Xiāngbǐ zhīxià, Zhōngguó gèng zhòngshì jiāpǔ hé zōngcí (shūmiàn jìlù), Kāmàilóng gèng zhòngshì kǒutóu chuánchéng (zhǎngzhě jiǎng gùshi).}
\\ \zh{最后，我的看法。我认为传统习俗有文化价值，但法律必须保护平等权利。我觉得我们应该尊重法律，也尊重长者的智慧，让传承更和谐。谢谢大家！}
\\ \textit{Zuìhòu, wǒ de kànfǎ. Wǒ rènwéi chuántǒng xísú yǒu wénhuà jiàzhí, dàn fǎlǜ bìxū bǎohù píngděng quánlì. Wǒ juéde wǒmen yīnggāi zūnzhòng fǎlǜ, yě zūnzhòng zhǎngzhě de zhìhuì, ràng chuánchéng gèng héxié. Xièxiè dàjiā!}

\tcblower
\textbf{Commentaire des choix de langue (Corrigé commenté) :}
\begin{itemize}
    \item \textbf{Cohésion :} Utilisation systématique des connecteurs \zh{首先 / 其次 / 然后 / 最后} pour baliser l'exposé (critère "Clarté").
    \item \textbf{Vocabulaire précis :} \zh{舅舅} (oncle maternel) pour la spécificité camerounaise (système matrilinéaire partiel chez les Bassa, Bamileke, etc.) vs \zh{宗祠 / 家谱} pour la spécificité chinoise. Le terme \zh{民法典} (Code civil) ancre le propos dans le droit moderne camerounais.
    \item \textbf{Grammaire ciblée :}
    \begin{itemize}
        \item \zh{跟中国一样\ldots{}共同点} (Similitude légale).
        \item \zh{跟中国不同\ldots{}特别} (Différence coutumière forte).
        \item \zh{相比之下} (Transition vers la comparaison culturelle profonde : écrit vs oral).
        \item \zh{虽然\ldots{}但是} (impliqué dans la partie A : \zh{习俗和法律有时不一样}).
        \item \zh{应该} + verbe (recommandation finale).
    \end{itemize}
    \item \textbf{Registre :} Langage soutenu mais accessible (HSK 3-4), phrases courtes, sujet souvent topicalisé (\zh{在喀麦隆\ldots{}, 根据\ldots{}}). Pas de structures trop complexes (pas de \zh{把} / \zh{被} inutiles).
    \item \textbf{Respect culturel :} Aucune hiérarchisation de valeur ("meilleur/pire"), termes neutres : \zh{特别} (particulier/spécifique), \zh{重视} (accorder de l'importance), \zh{和谐} (harmonieux).
\end{itemize}

\textbf{Variante écrite (Cahier -- 7 phrases) :}
\\ \zh{今天我们对比喀麦隆和中国的遗产传承。首先，喀麦隆习俗常传给长子或舅舅的儿子，但法律规定男女平等。其次，中国传统有长子继承制，现在法律也平等，农村仍倾向留房给儿子。跟中国一样，两国法律都保障平等；跟中国不同，喀麦隆有“舅舅继承”习俗。相比之下，中国重家谱书面记录，喀麦隆重口头传承。我认为应该尊重法律平等，也珍惜传统文化智慧。这样传承才和谐。}
\\ \textit{Jīntiān wǒmen duìbǐ Kāmàilóng hé Zhōngguó de yíchán chuánchéng. Shǒuxiān, Kāmàilóng xísú cháng chuán gěi zhǎngzǐ huò jiùjiu de érzi, dàn fǎlǜ guīdìng nánnǚ píngděng. Qícì, Zhōngguó chuántǒng yǒu zhǎngzǐ jìchéngzhì, xiànzài fǎlǜ yě píngděng, nóngcūn réng qīngxiàng liú fáng gěi érzi. Gēn Zhōngguó yīyáng, liǎngguó fǎlǜ dōu bǎozhàng píngděng; gēn Zhōngguó bù yīyáng, Kāmàilóng yǒu "jiùjiu jìchéng" xísú. Xiāngbǐ zhīxià, Zhōngguó zhòng jiāpǔ shūmiàn jìlù, Kāmàilóng zhòng kǒutóu chuánchéng. Wǒ rènwéi yīnggāi zūnzhòng fǎlǜ píngděng, yě zhēnxī chuántǒng wénhuà zhìhuì. Zhèyàng chuánchéng cái héxié.}
\end{exoresolu}

\courssub{Bilan de l'activité}
\noindent À l'issue de cette activité, l'enseignant guide une courte méta-réflexion collective (5 min) à l'aide des questions suivantes, affichées au tableau :

\begin{itemize}
    \item \textbf{Sur le fond (Socio-culturel) :} \og Avez-vous découvert une similitude surprenante ? Une différence qui vous a interpellé ? Pourquoi le \zh{舅舅} (oncle maternel) apparaît-il au Cameroun et pas en Chine ? (Réponse attendue : systèmes de parenté matrilinéaire vs patrilinéaire). \fg{}
    \item \textbf{Sur la forme (Langue) :} \og Quelle structure de comparaison avez-vous trouvée la plus naturelle à l'oral ? (\zh{跟\ldots{}一样} vs \zh{比} vs \zh{相比之下}). Avez-vous réussi à placer \zh{应该} avant le verbe ? Avez-vous respecté l'ordre : \zh{时间 / 地点 / 方式 + Verbe} ? \fg{}
    \item \textbf{Sur la compétence orale :} \og Le chronomètre a-t-il été respecté ? La prononciation des tons (notamment \zh{倾向于 qīngxiàngyú}, \zh{相比之下 xiāngbǐ zhīxià}, \zh{继承 jìchéng}) était-elle stable ? \fg{}
    \item \textbf{Sur la citoyenneté / Vivre ensemble :} \og Comment le fait de \textbf{comparer sans juger} (vocabulaire : \zh{特别}, \zh{不同}, \zh{重视} au lieu de \zh{奇怪}, \zh{错误}, \zh{落后}) change-t-il la qualité du dialogue interculturel ? \fg{}
\end{itemize}

\noindent \textbf{Trace écrite obligatoire dans le cahier :}
\begin{enumerate}
    \item Copier le tableau de vocabulaire (section Ressources).
    \item Coller/Recopier la version écrite du mini-exposé (7 phrases) en caractères, pinyin et traduction.
    \item Noter 2 structures de comparaison "modèles" à réutiliser (ex: \zh{跟中国不同，喀麦隆\ldots{} ; 相比之下，中国\ldots{}}).
    \item Rédiger une phrase personnelle avec \zh{我认为\ldots{}应该\ldots{}} sur un autre sujet de société (ex: mariage, éducation, fête).
\end{enumerate}

\noindent \textbf{Évaluation formative (Grille professeur) :}
\begin{center}
\begin{tabularx}{\linewidth}{|X|X|X|X|X|}\hline
\rowcolor{poporangeL} \textbf{Critère} & \textbf{4 -- Maîtrise} & \textbf{3 -- Acquis} & \textbf{2 -- Fragile} & \textbf{1 -- Non acquis} \\ \hline
Vocabulaire thème (10 mots+) & Richesse, précision, tons justes & Correct, 7-9 mots, tons OK & Limité ($<$7), erreurs tons & Très pauvre, tons faux \\ \hline
Structures comparaison & 3 structures variées, justes & 2 structures justes & 1 structure / maladresses & Aucune / erreurs bloquantes \\ \hline
Cohérence / Connecteurs & Fluide, \zh{首先\ldots{}最后}, logiques & Correct, connecteurs présents & Saccadé, connecteurs manquants & Désorganisé \\ \hline
Expression opinion & \zh{我认为/觉得/应该} justes, nuancés & Opinion claire, structure OK & Opinion simple, structure fautive & Pas d'opinion / HSK 1 only \\ \hline
Prononciation / Tons & Naturels, tons 1-4 stables, \zh{ü} correct & Compréhensible, quelques hésitations & Difficile à suivre, tons instables & Incompréhensible \\ \hline
Respect culturel / Contenu & Factuel, nuancé, 0 jugement de valeur & Factuel, 1 jugement léger & Stéréotypes, jugements & Offensant / HSK 0 \\ \hline
\end{tabularx}
\end{center}

\noindent \textbf{Prolongement possible (Devoirs / Projet) :} Rédiger un email de réponse à Li Wei (80-100 caractères) pour le remercier et lui poser 2 questions précises sur la \zh{宗祠} ou la \zh{家谱} aujourd'hui en ville. Utiliser \zh{请问}, \zh{想了解}, \zh{谢谢}.

\newpage

\coursec{Méthodes et exercices résolus}

\courssub{Rappel ciblé et mise en route}

\begin{aretenir}[Rappel de la Leçon 4 : Vocabulaire et concepts de base]
La leçon 4 a introduit les notions fondamentales de la famille et de l'héritage. Voici les points clés à réactiver :
\begin{itemize}
    \item \textbf{Structure familiale :} \zh{大家庭} (dàjiātíng, famille élargie) vs \zh{小家庭} (xiǎojiātíng, famille nucléaire) ; \zh{族长} (zúzhǎng, chef de clan) ; \zh{酋长} (qiúzhǎng, chef traditionnel/village).
    \item \textbf{Principes coutumiers :} \zh{重男轻女} (zhòng nán qīng nǔ, préférer les garçons aux filles) ; \zh{长子继承制} (zhǎngzǐ jìchéngzhì, primogéniture) ; \zh{习惯法} (xíguànfǎ, droit coutumier).
    \item \textbf{Base légale moderne :} \zh{法律} (fǎlǜ, loi) ; \zh{民法典} (Mínfǎdiǎn, Code civil) ; \zh{继承法} (Jìchéngfǎ, Loi sur les successions -- Chine, 1985/2021) ; \zh{男女平等} (nánnǚ píngděng, égalité hommes-femmes).
    \item \textbf{Actes juridiques :} \zh{立遗嘱} (lì yízhǔ, faire un testament) ; \zh{公证} (gōngzhèng, notarier/notarisation) ; \zh{分割遗产} (fēngē yíchǎn, partager la succession).
\end{itemize}
\textbf{Mise en route :} En binôme, citez \textbf{en chinois} 3 différences entre l'héritage « selon la coutume » et l'héritage « selon la loi » que vous avez retenus. Utilisez \zh{但是} (dànshì) ou \zh{而} (ér).
\end{aretenir}

\courssub{L'hérédité au Cameroun : faits}

\begin{textebox}[Faits socioculturels : Le Cameroun -- Dualisme juridique et coutumier]
\zh{喀麦隆的遗产继承呈现“双轨制”特点：现代成文法与传统习惯法并存。} \\
\zh{1. 习惯法地位高：宪法承认习惯法，农村地区尤为盛行。} \\
\zh{2. 土地与身份：祖传土地（祖传土地）多按男系传承，长子（长子）负责祭祀与管理，女儿（女儿）嫁出后常失去继承权，仅保留使用权（使用权）。} \\
\zh{3. 寡妇处境：传统上寡妇（寡妇）无继承权，甚至可能被“继承”给亡夫兄弟（ levirate / 继娶），但法院判例越来越多保护寡妇居住权。} \\
\zh{4. 现代法典进步：2016年《民法典》（民法典）明确规定男女平等继承权（男女平等继承权），法定继承（法定继承）顺序不分性别。} \\
\zh{5. 司法实践：城市法院常依据《民法典》判决女儿分得份额（份额），但执行在乡村常受阻于酋长调解（酋长调解）。}

\textbf{Pinyin :} Kāmàilóng de yíchǎn jìchéng chéngxiàn "shuāngguīzhì" tèdìng: xiàndài chéngwénfǎ yǔ chuántǒng xíguànfǎ bìngcún. 1. Xíguànfǎ dìwèi gāo: xiànfǎ chéngrèn xíguànfǎ, nóngcūn dìqū yóuwéi shèngxíng. 2. Tǔdì yǔ shēnfèn: zǔchuán tǔdì duō àn nánxì chuánchéng, zhǎngzǐ fùzé jìsì yǔ guǎnlǐ, nǚ'ér jiàchū hòu cháng shīqù jìchéngquán, jǐn bǎoliú shǐyòngquán. 3. Guǎfù chǔjìng: chuántǒng shàng guǎfù wú jìchéngquán, shènzhì kěnéng bèi "jìchéng" gěi wángfū xiōngdì (levirate / jìqǔ), dàn fǎyuàn pànlì yuèláiyuèduō bǎohù guǎfù jūzhùquán. 4. Xiàndài fǎdiǎn jìnbù: 2016 nián "Mínfǎdiǎn" míngquè guīdìng nánnǚ píngděng jìchéngquán, fǎdìng jìchéng shùnxù bù fēn xìngbié. 5. Sīfǎ shíjiàn: chéngshì fǎyuàn cháng yījù "Mínfǎdiǎn" pànjué nǚ'ér fēndé fèn'é, dàn zhíxíng zài xiāngcūn cháng zǔzú yú qiúzhǎng tiáojiě.

\textbf{Traduction :} L'héritage au Cameroun présente une « double voie » : coexistence de la loi écrite moderne et du droit coutumier traditionnel. 1. Statut fort du droit coutumier : la Constitution le reconnaît, prédominant en zone rurale. 2. Terre et identité : les terres ancestrales se transmettent par la lignée masculine, l'aîné gère les rites et le domaine, les filles mariées perdent souvent leurs droits, ne gardant qu'un droit d'usage. 3. Situation des veuves : traditionnellement sans droits, parfois « léguées » au frère du défunt (lévirat), mais la jurisprudence protège de plus en plus leur droit au logement. 4. Progrès du Code civil (2016) : stipule l'égalité successorale hommes-femmes, ordre successoral sans distinction de sexe. 5. Pratique judiciaire : les tribunaux urbains jugent selon le Code civil pour attribuer des parts aux filles, mais l'exécution en zone rurale bute souvent sur la médiation du chef.
\end{textebox}

\courssub{L'hérédité en Chine : faits}

\begin{textebox}[Faits socioculturels : La Chine -- Unification légale et réalités urbaines/rurales]
\zh{中国遗产继承经历了从封建礼教到法治统一的转型。} \\
\zh{1. 立法统一：1985年《继承法》（继承法）废除封建长子继承制（封建长子继承制），确立法定继承男女平等（法定继承男女平等）。2021年《民法典》（民法典）继承编延续此原则，新增遗嘱形式（打印、录像、口头等）和宽容遗嘱制度（宽容遗嘱制度）。} \\
\zh{2. 城市实践：房产（房产）、存款为主，普遍通过公证遗嘱（公证遗嘱）或自书遗嘱分配，男女份额均等（份额均等），纠纷进法院（法院诉讼）按证据判决。} \\
\zh{3. 农村特殊性：宅基地（宅基地）所有权归集体，仅使用权可继承。“分家析产”（分家析产）习俗依存，长子/儿子常分得主宅（主宅），出嫁女（出嫁女）常被排除或仅得现金补偿。} \\
\zh{4. 养老与继承挂钩：赡养义务（赡养义务）影响继承份额，「尽了主要赡养义务的，可以多分」（尽了主要赡养义务的，可以多分）。}

\textbf{Pinyin :} Zhōngguó yíchǎn jìchéng jīnglì le cóng fēngjiàn lǐjiào dào fǎzhì tǒngyī de zhuǎnxíng. 1. Lìfǎ tǒngyī: 1985 nián "Jìchéngfǎ" fèichú fēngjiàn zhǎngzǐ jìchéngzhì, quèdìng fǎdìng jìchéng nánnǚ píngděng. 2021 nián "Mínfǎdiǎn" jìchéng biān yánxù cǐ yuánzé, xīnzēng yízhǔ xíngshì (dǎyìn, lùxiàng, kǒutóu děng) hé kuāngróng yízhǔ zhìdù. 2. Chéngshì shíjiàn: fángchǎn, cúnkuǎn wéi zhǔ, pǔbiàn tōngguò gōngzhèng yízhǔ huò zìshū yízhǔ fēnpèi, nánnǚ fèn'é jūnděng, jiūfēn jìn fǎyuàn àn zhèngjù pànjué. 3. Nóngcūn tèshūxìng: zhájīdì suǒyǒuquán guī jítǐ, jǐn shǐyòngquán kě jìchéng. "Fēnjiā xīchǎn" xísú yícún, zhǎngzǐ/érzi cháng fēndé zhǔzhái, jiàchū nǚ cháng bèi páichú huò jǐn dé xiànjīn bǔcháng. 4. Yǎnglǎo yǔ jìchéng guàgōu: shànyǎng yìwù yǐngxiǎng jìchéng fèn'é, "jìn le zhǔyào shànyǎng yìwù de, kěyǐ duōfēn".

\textbf{Traduction :} La Chine a opéré une transition de l'éthique féodale vers l'unification par l'État de droit. 1. Unification législative : Loi sur les successions 1985 abolit la primogéniture féodale, instaure l'égalité successorale légale. Code civil 2021 (partie successions) poursuit ce principe, ajoute formes de testaments (imprimé, vidéo, oral) et régime du testament tolérant. 2. Pratique urbaine : immobilier/épargne, partage via testament notarié/olographe, parts égales, contentieux au tribunal sur preuves. 3. Spécificité rurale : terre collective (droit d'usage seulement transmissible), coutume « division familiale » persiste, aîné/fils gardent la maison principale, filles mariées exclues ou compensées en numéraire. 4. Lien entretien/héritage : l'obligation d'entretien influence la part successorale (« celui qui a assumé la charge principale d'entretien peut recevoir davantage »).
\end{textebox}

\courssub{Vocabulaire du thème : l'hérédité (juridique et coutumier)}

\begin{definition}[Lexique spécialisé -- Tableau récapitulatif]
\begin{center}
\begin{tabularx}{\linewidth}{|X|X|X|X|}
\hline
\rowcolor{popblueL} \textbf{Caractères} & \textbf{Pinyin} & \textbf{Nature} & \textbf{Traduction} \\
\hline
\zh{习惯法} & xíguànfǎ & n. & droit coutumier \\
\zh{法定继承} & fǎdìng jìchéng & n. & succession légale / dévolution légale \\
\zh{长子继承制} & zhǎngzǐ jìchéngzhì & n. & primogéniture (système de l'aîné) \\
\zh{男女平等} & nánnǚ píngděng & adj./n. & égalité hommes-femmes \\
\zh{遗产} & yíchǎn & n. & succession, héritage, patrimoine \\
\zh{遗嘱} & yízhǔ & n. & testament \\
\zh{公证} & gōngzhèng & v./n. & notarier / notarisation, acte notarié \\
\zh{继承权} & jìchéngquán & n. & droit de succession \\
\zh{份额} & fèn'é & n. & part, quote-part \\
\zh{纠纷} & jiūfēn & n. & litige, conflit, dispute \\
\zh{分割遗产} & fēngē yíchǎn & v. & partager la succession, partir l'héritage \\
\zh{立遗嘱} & lì yízhǔ & v. & faire un testament \\
\zh{放弃继承权} & fàngqì jìchéngquán & v. & renoncer à ses droits successoraux \\
\zh{祖传土地} & zǔchuán tǔdì & n. & terre ancestrale, terre de lignée \\
\zh{宅基地} & zhájīdì & n. & terrain d'assiette (rural, collectif) \\
\zh{分家析产} & fēnjiā xīchǎn & v./n. & division familiale, partage du patrimoine familial \\
\zh{寡妇} & guǎfù & n. & veuve \\
\zh{妇女权益} & fùnǚ quányì & n. & droits et intérêts des femmes \\
\zh{赡养义务} & shànyǎng yìwù & n. & obligation d'entretien (parents âgés) \\
\zh{酋长调解} & qiúzhǎng tiáojiě & n. & médiation du chef traditionnel \\
\zh{双轨制} & shuāngguīzhì & n. & double système / dualisme (légal + coutumier) \\
\hline
\end{tabularx}
\end{center}
\end{definition}

\begin{attention}[Caractères proches et pièges d'écriture]
\begin{itemize}
    \item \zh{遗} (yí, legs, radical \zh{辶} « marche ») \textbf{vs} \zh{疑} (yí, douter, radical \zh{疒} « maladie »). \textit{Astuce :} ce qu'on laisse derrière soi en partant (\zh{辶}) = héritage.
    \item \zh{承} (chéng, porter, recevoir, radical \zh{手} « main ») \textbf{vs} \zh{成} (chéng, devenir, réussir, radical \zh{戈} « hallebarde »). \textit{Astuce :} hériter = porter (\zh{手}) la charge du défunt.
    \item \zh{继} (jì, continuer, hériter, radical \zh{纟} « fil ») \textbf{vs} \zh{纪} (jì, discipline, radical \zh{纟}). \zh{继承} (jìchéng) \neq \zh{纪成}.
    \item \zh{法} (fǎ, loi, radical \zh{氵} « eau ») \textbf{vs} \zh{灋} (fǜ, variante ancienne/bouddhique). \textbf{Utilisez toujours \zh{法} en chinois moderne} (\zh{法律}, \zh{法典}, \zh{法院}, \zh{普法}).
\end{itemize}
\end{attention}

\courssub{Comparer et donner son avis en chinois}

\begin{methode}[Comparer deux pratiques socioculturelles : structures \cle{比}, \cle{跟……一样/不一样}, \cle{而/但是}, \cle{相比之下}]
\textbf{Objectif :} Produire des phrases correctes pour opposer ou rapprocher les coutumes d'hérédité au Cameroun et en Chine.
\begin{enumerate}
    \item \textbf{Identifier les deux termes de la comparaison} (A et B) : ex. \zh{喀麦隆} (Cameroun) vs \zh{中国} (Chine) ; \zh{传统习俗} (coutumes traditionnelles) vs \zh{现代法律} (loi moderne) ; \zh{城市} (ville) vs \zh{农村} (campagne).
    \item \textbf{Choisir la structure adaptée} :
        \begin{itemize}
            \item \textbf{Supériorité / inégalité (quantifiable) :} \textbf{A + 比 + B + Adj (+ 程度补语)}. Ex: \zh{喀麦隆的习惯法比现代法律复杂得多。}
            \item \textbf{Différence marquée (nature) :} \textbf{A + 跟/和 + B + 不一样} ; \textbf{A + 而/但是 + B + ...} (A..., alors que B...). Ex: \zh{中国靠立法统一，而喀麦隆靠司法纠偏。}
            \item \textbf{Similitude :} \textbf{A + 跟/和 + B + 一样} ; \textbf{A + 和 + B + 都...}. Ex: \zh{两国都在走向男女平等。}
            \item \textbf{Comparaison globale (niveau discours) :} \zh{相比之下} (xiāngbǐ zhīxià, en comparaison / par contraste) en tête de phrase.
        \end{itemize}
    \item \textbf{Placer le complément de degré} (\zh{得多}, \zh{多了}, \zh{一点儿}) \textbf{après l'adjectif} dans la structure \zh{比}. \zh{比中国复杂得多} (bǐ Zhōngguó fùzá de duō).
    \item \textbf{Ordre des mots impératif :} Sujet A + 比/跟 + Objet B + Adverbe + Adjectif. \textbf{Piège francophone :} Ne pas calquer « La Chine est plus... que le Cameroun » $\rightarrow$ \zh{中国比喀麦隆...} (Correct si la Chine est le sujet A). Si le sujet est le Cameroun : \zh{喀麦隆比中国...}.
\end{enumerate}
\end{methode}

\begin{exoresolu}[Comparaison écrite : droits successoraux -- Application des structures]
\textbf{Énoncé :} Traduisez en chinois les phrases suivantes en utilisant les structures de comparaison vues en classe.
\begin{enumerate}
    \item Au Cameroun, la coutume traditionnelle est plus complexe que la loi moderne.
    \item En Chine, la loi actuelle accorde les mêmes droits aux fils et aux filles, alors qu'autrefois le fils aîné héritait de tout.
    \item Les deux pays ont maintenant des lois écrites pour protéger l'héritage des femmes.
    \item Par contraste, l'unification légale en Chine est plus complète qu'au Cameroun.
\end{enumerate}
\tcblower
\textbf{Solution détaillée :}
\begin{enumerate}
    \item \textbf{Analyse :} Comparaison de complexité. Sujet A = \zh{喀麦隆的传统习俗} (Kāmàilóng de chuántǒng xísú), Sujet B = \zh{现代法律} (xiàndài fǎlǜ). Adj = \zh{复杂} (fùzá). Structure \cle{比} + degré \zh{得多}.
    \textbf{Réponse :} \zh{喀麦隆的传统习俗比现代法律复杂得多。} (Kāmàilóng de chuántǒng xísú bǐ xiàndài fǎlǜ fùzá de duō.)
    \item \textbf{Analyse :} Deux propositions contrastées. P1 : Égalité actuelle (\zh{跟……一样} / \zh{享有同样的...}). P2 : Inégalité passée (\zh{而} formel / \zh{但是} oral). Vocab : \zh{儿子} (érzi), \zh{女儿} (nǚ'ér), \zh{长子} (zhǎngzǐ), \zh{继承一切} (jìchéng yīqiè), \zh{规定} (guīdìng), \zh{享有} (xiǎngyǒu), \zh{同样的} (tóngyàng de).
    \textbf{Réponse :} \zh{中国现在的法律规定儿子跟女儿享有同样的继承权，而过去长子继承一切。} (Zhōngguó xiànzài de fǎlǜ guīdìng érzi gēn nǚ'ér xiǎngyǒu tóngyàng de jìchéngquán, ér guòqù zhǎngzǐ jìchéng yīqiè.)
    \item \textbf{Analyse :} Similitude (\zh{都} + verbe d'action \zh{制定} / \zh{有}). Sujet composé : \zh{喀麦隆和中国} (Kāmàilóng hé Zhōngguó). Objet : \zh{保护妇女继承权的法律} (bǎohù fùnǚ jìchéngquán de fǎlǜ).
    \textbf{Réponse :} \zh{喀麦隆和中国都制定了保护妇女继承权的法律。} (Kāmàilóng hé Zhōngguó dōu zhìdìng le bǎohù fùnǚ jìchéngquán de fǎlǜ.)
    \item \textbf{Analyse :} Comparaison globale \zh{相比之下} + \cle{比}. Sujet A = \zh{中国的法律统一} (Zhōngguó de fǎlǜ tǒngyī), B = \zh{喀麦隆} (Kāmàilóng). Adj = \zh{完全} (wánquán) / \zh{彻底} (chèdǐ).
    \textbf{Réponse :} \zh{相比之下，中国的法律统一比喀麦隆更彻底。} (Xiāngbǐ zhīxià, Zhōngguó de fǎlǜ tǒngyī bǐ Kāmàilóng gèng chèdǐ.)
\end{enumerate}
\textbf{Conclusion :} Respecter l'ordre Sujet--Comparateur--Objet--Adj. Utiliser \zh{而} pour l'opposition formelle (écrit), \zh{但是} à l'oral. \zh{相比之下} structure le paragraphe.
\end{exoresolu}

\begin{methode}[Mobiliser le lexique socioculturel spécialisé : collocations et classification]
\textbf{Objectif :} Employer avec justesse le vocabulaire de l'hérédité (coutumier / légal) en caractères, pinyin et sens.
\begin{enumerate}
    \item \textbf{Classer par domaine sémantique} (voir Tableau en \textbf{Definition} ci-dessus).
    \item \textbf{Mémoriser les collocations verbe + objet (VP) indispensables} :
        \begin{center}
        \begin{tabularx}{\linewidth}{|X|X|X|}
        \hline
        \rowcolor{popgreenL} \textbf{Verbe} & \textbf{Objet} & \textbf{Collocation (VP)} \\
        \hline
        \zh{办理} (bànlǐ) & \zh{公证} (gōngzhèng) & \zh{办理公证} (faire la notarisation) \\
        \zh{发生} (fāshēng) & \zh{纠纷} (jiūfēn) & \zh{发生纠纷} (survenir un litige) \\
        \zh{解决} (jiějué) / \zh{调解} (tiáojiě) & \zh{纠纷} (jiūfēn) & \zh{解决/调解纠纷} \\
        \zh{放弃} (fàngqì) & \zh{继承权} (jìchéngquán) & \zh{放弃继承权} \\
        \zh{享有} (xiǎngyǒu) & \zh{份额/权利} (fèn'é/quánlì) & \zh{享有份额/权利} \\
        \zh{保护} (bǎohù) & \zh{妇女权益} (fùnǚ quányì) & \zh{保护妇女权益} \\
        \zh{废除} (fèichú) & \zh{长子继承制} (zhǎngzǐ jìchéngzhì) & \zh{废除长子继承制} \\
        \hline
        \end{tabularx}
        \end{center}
    \item \textbf{Attention aux tons :} \zh{继承} (jì-chéng, 4-2), \zh{土地} (tǔ-dì, 3-4), \zh{公证} (gōng-zhèng, 1-4), \zh{份额} (fèn-é, 4-2).
\end{enumerate}
\end{methode}

\begin{exoresolu}[Vocabulaire en contexte : complétion et traduction]
\textbf{Énoncé :} Complétez les phrases avec le mot juste (caractères + pinyin) puis traduisez.
\textit{Banque :} \zh{公证, 习惯法, 份额, 立遗嘱, 纠纷, 长子继承制, 法定继承}
\begin{enumerate}
    \item 为了避免家庭 \_\_\_\_\_\_ ，老人决定提前 \_\_\_\_\_\_ 并去 \_\_\_\_\_\_ 。
    \item 许多非洲国家保留了 \_\_\_\_\_\_ ，但在城市里 \_\_\_\_\_\_ 逐渐被 \_\_\_\_\_\_ 取代。
    \item 每个继承人依法享有平等的 \_\_\_\_\_\_ 。
\end{enumerate}
\tcblower
\textbf{Solution détaillée :}
\begin{enumerate}
    \item \textbf{Contexte :} Prévention de conflit $\rightarrow$ \zh{纠纷} (jiūfēn). Action préventive $\rightarrow$ \zh{立遗嘱} (lì yízhǔ). Formalisation $\rightarrow$ \zh{公证} (gōngzhèng).
    \textbf{Phrase :} \zh{为了避免家庭纠纷，老人决定提前立遗嘱并去公证。} (Wèile bìmiǎn jiātíng jiūfēn, lǎorén juédìng tíqián lì yízhǔ bìng qù gōngzhèng.)
    \textbf{Traduction :} Pour éviter les litiges familiaux, la personne âgée a décidé de faire un testament à l'avance et d'aller chez le notaire.
    \item \textbf{Contexte :} Système traditionnel général $\rightarrow$ \zh{习惯法} (xíguànfǎ). Pratique spécifique remplacée $\rightarrow$ \zh{长子继承制} (zhǎngzǐ jìchéngzhì). Remplaçant moderne $\rightarrow$ \zh{法定继承} (fǎdìng jìchéng).
    \textbf{Phrase :} \zh{许多非洲国家保留了习惯法，但在城市里长子继承制逐渐被法定继承取代。} (Xǔduō Fēizhōu guójiā bǎoliú le xíguànfǎ, dàn zài chéngshì lǐ zhǎngzǐ jìchéngzhì zhújiàn bèi fǎdìng jìchéng tìdài.)
    \textbf{Traduction :} De nombreux pays africains ont conservé le droit coutumier, mais en ville le système de primogéniture est progressivement remplacé par la succession légale.
    \item \textbf{Contexte :} Droit individuel dans le partage $\rightarrow$ \zh{份额} (fèn'é).
    \textbf{Phrase :} \zh{每个继承人依法享有平等的份额。} (Měige jìchéngrén yīfǎ xiǎngyǒu píngděng de fèn'é.)
    \textbf{Traduction :} Chaque héritier jouit légalement d'une quote-part égale.
\end{enumerate}
\end{exoresolu}

\begin{methode}[Exprimer un point de vue personnel argumenté : \cle{我认为}, \cle{应该}, \cle{因为……所以……}, \cle{一方面……另一方面……}]
\textbf{Objectif :} Rédiger un paragraphe cohérent (80--120 caractères) donnant son avis sur l'évolution des pratiques successorales.
\begin{enumerate}
    \item \textbf{Poser la thèse} : \zh{我认为……} (Wǒ rènwéi...) / \zh{在我看来……} (Zài wǒ kànlái...).
    \item \textbf{Structurer deux arguments} (Dialectique) :
        \begin{itemize}
            \item \textbf{Argument 1 (Modernité/Droit) :} \zh{一方面，现代法律规定男女平等，保护了女儿的继承权和妇女权益……} (Yīfāngmiàn...).
            \item \textbf{Argument 2 (Tradition/Realité sociale) :} \zh{另一方面，尊重传统习俗有利于家族团结，但不能违背法律原则……} (Lìngyī fāngmiàn...).
        \end{itemize}
    \item \textbf{Conclure par une recommandation nuancée} : \zh{所以，我们应该……} (Suǒyǐ, wǒmen yīnggāi...) / \zh{建议……} (Jiànyì...). Utiliser \zh{应该} (devoir moral/légal), \zh{必须} (obligation forte), \zh{最好} (conseil).
    \item \textbf{Connecteurs logiques (cohérence) :}
        \begin{itemize}
            \item Cause $\rightarrow$ Conséquence : \zh{因为……所以……} (Yīnwèi... suǒyǐ...).
            \item Concession : \zh{虽然……但是/可是……} (Suīrán... dànshì/kěshì...).
            \item Progressif : \zh{此外} (cǐwài, de plus), \zh{而且} (érqiě, de plus/même).
        \end{itemize}
    \item \textbf{Ponctuation chinoise :} ， (clauses), 。 (fin), 、 (énumération noms), ； (phrases indépendantes liées).
\end{enumerate}
\end{methode}

\begin{exoresolu}[Production écrite : « Faut-il garder la primogéniture ? »]
\textbf{Énoncé :} Rédigez un court paragraphe (environ 100 caractères) en chinois pour exprimer votre avis sur le maintien ou l'abandon du système de primogéniture (\zh{长子继承制}). Utilisez au moins 6 mots du vocabulaire de la leçon (surlignés ci-dessous).
\tcblower
\textbf{Solution détaillée :}
\textbf{Plan :} Thèse (abandon nécessaire) -- Arg 1 (égalité/droit) -- Arg 2 (conflits/paix) -- Conclusion (loi + dialogue/conciliation).
\textbf{Proposition de rédaction :}
\zh{我认为应该废除长子继承制。一方面，现代法律规定男女平等，女儿也有继承权，保护妇女权益是法治的体现。另一方面，只让长子继承容易引起家庭纠纷，破坏亲情团结。所以，我们应该依法分割遗产，同时提倡家庭协商，既尊重传统孝道，又实现公平正义。}
(Pinyin : Wǒ rènwéi yīnggāi fèichú zhǎngzǐ jìchéngzhì. Yīfāngmiàn, xiàndài fǎlǜ guīdìng nánnǚ píngděng, nǚ'ér yě yǒu jìchéngquán, bǎohù fùnǚ quányì shì fǎzhì de tǐxiàn. Lìngyī fāngmiàn, zhǐ ràng zhǎngzǐ jìchéng róngyì yǐnqǐ jiātíng jiūfēn, pòhuài qīnqíng tuánjié. Suǒyǐ, wǒmen yīnggāi yīfǎ fēngē yíchǎn, tóngshí tíchàng jiātíng xiéshāng, jì zūnzhòng chuántǒng xiàodào, yòu shíxiàn gōngpíng zhèngyì.)
\textbf{Traduction :} Je pense qu'il faut abolir la primogéniture. D'un côté, la loi moderne stipule l'égalité hommes-femmes, les filles ont aussi des droits successoraux, protéger les droits des femmes est la manifestation de l'État de droit. De l'autre, ne laisser hériter que l'aîné provoque facilement des litiges familiaux, détruisant l'union affective. Donc, nous devons partager la succession selon la loi, et en même temps promouvoir la concertation familiale, respectant à la fois la piété filiale traditionnelle et réalisant l'équité et la justice.
\textbf{Vocabulaire utilisé (7) :} \zh{长子继承制, 男女平等, 继承权, 妇女权益, 家庭纠纷, 分割遗产, 协商} (+ \zh{法治}, \zh{孝道}, \zh{公平正义}).
\end{exoresolu}

\courssub{Questions de réflexion et préparation du mini-exposé}

\begin{methode}[Analyser un texte comparatif authentique : stratégie de lecture en 3 temps]
\textbf{Objectif :} Comprendre un texte chinois comparant l'hérédité au Cameroun et en Chine pour en extraire les faits et le point de vue de l'auteur.
\begin{enumerate}
    \item \textbf{Repérage global (1ère lecture) :} Identifier le \textbf{thème} (héritage), les \textbf{deux pays} (\zh{喀麦隆}, \zh{中国}), la \textbf{structure} (paragraphes : 1. Intro / 2. Cameroun / 3. Chine / 4. Comparaison / 5. Avis). Souligner les \textbf{connecteurs logiques} : \zh{然而} (cependant), \zh{此外} (de plus), \zh{相比之下} (en comparaison), \zh{总的来说} (en résumé), \zh{但是} (mais).
    \item \textbf{Lecture de détail (2ème lecture) :} Chercher les \textbf{faits chiffrés/datés} (ex. \zh{2016年}, \zh{1985年}, \zh{《民法典》}), les \textbf{termes techniques} (\zh{习惯法}, \zh{法定继承}, \zh{公证处}, \zh{酋长}), les \textbf{verbes modaux} (\zh{必须}, \zh{可以}, \zh{倾向于}, \zh{往往}). Noter les \textbf{différences} (coutume vs loi, oral vs écrit, dualisme vs unification) et \textbf{similitudes} (évolution vers l'égalité, rôle de la médiation/famille).
    \item \textbf{Synthèse (3ème lecture) :} Remplir un \textbf{tableau comparatif} (voir ci-dessous). Rédiger une \textbf{phrase résumé} en français.
\end{enumerate}
\end{methode}

\begin{exoresolu}[Compréhension de texte : « Deux chemins vers l'égalité »]
\textbf{Énoncé :} Lisez le texte ci-dessous et répondez aux questions en français.

\begin{textebox}[Texte support : 两条通往平等的路]
\zh{喀麦隆和中国在遗产继承问题上有相似的历史：传统上都偏向男性。} \\
\zh{在喀麦隆，习惯法依然强势，许多家庭按“长子继承制”分割祖传土地，女儿往往没有份额。} \\
\zh{但是，现代法典《民法典》已规定男女平等继承权，法院也常判决保护寡妇和女儿。} \\
\zh{在中国，1985年《继承法》确立了法定继承男女平等，废除了封建长子继承制。} \\
\zh{现在，城市家庭多通过公证遗嘱分配房产，农村地区仍有“分家析产”的习俗。} \\
\zh{相比之下，中国的法律统一性更强，喀麦隆则呈现“双轨制”并存。} \\
\zh{总的来说，两国都在从“重男轻女”走向“男女平等”，但路径不同：中国靠立法统一，喀麦隆靠司法逐步纠偏。}
\end{textebox}

\begin{enumerate}
    \item Quel est le point commun historique entre les deux pays ? (Citez le passage chinois).
    \item Quelle est la situation actuelle au Cameroun concernant le droit coutumier et le code civil ? (Deux faits distincts).
    \item Quelle date clé marque le changement légal en Chine ? Quelle pratique a été abolie ?
    \item Comment l'auteur caractérise-t-il la différence de « route » (路径) vers l'égalité ?
    \item Traduisez la dernière phrase en français.
\end{enumerate}
\tcblower
\textbf{Solution détaillée :}
\begin{enumerate}
    \item \textbf{Repérage :} 1ère phrase. \zh{传统上都偏向男性} (chuántǒng shàng dōu piānxiàng nánxìng).
    \textbf{Réponse :} Traditionnellement, les deux pays favorisaient les hommes (héritage patrilinéaire / primogéniture).
    \item \textbf{Repérage :} Paragraphe 2 (Cameroun). Fait 1 : \zh{习惯法依然强势} (xíguànfǎ yīrán qiángshì - le droit coutumier reste fort/prédominant). Fait 2 : \zh{现代法典《民法典》已规定男女平等继承权} (xiàndài fǎdiǎn "Mínfǎdiǎn" yǐ guīdìng nánnǚ píngděng jìchéngquán - le code civil moderne stipule l'égalité successorale).
    \textbf{Réponse :} Coexistence d'un droit coutumier fort (primogéniture, exclusion des filles) et d'un code civil moderne garantissant l'égalité, appliqué par les tribunaux.
    \item \textbf{Repérage :} Paragraphe 3 (Chine). Date : \zh{1985年}. Loi : \zh{《继承法》}. Pratique abolie : \zh{封建长子继承制} (fēngjiàn zhǎngzǐ jìchéngzhì - système féodal de primogéniture).
    \textbf{Réponse :} 1985, promulgation de la Loi sur les successions ; abolition de la primogéniture féodale.
    \item \textbf{Repérage :} Dernier paragraphe. \zh{中国靠立法统一，喀麦隆靠司法逐步纠偏} (Zhōngguó kào lìfǎ tǒngyī, Kāmàilóng kào sīfǎ zhúbù jiūpiān).
    \textbf{Réponse :} La Chine a unifié le système par la législation (loi unique nationale) ; le Cameroun corrige progressivement les inégalités par la jurisprudence/jugements (double système légal/coutumier).
    \item \textbf{Traduction :} « En résumé, les deux pays sont en train de passer du "préférer les garçons aux filles" à "l'égalité hommes-femmes", mais les chemins sont différents : la Chine compte sur la législation pour unifier, le Cameroun compte sur la justice pour corriger progressivement le tir. »
\end{enumerate}
\end{exoresolu}

\begin{experience}[Débat guidé : « Tradition et Droit : comment concilier ? »]
\textbf{Consigne :} En groupes de 4. Chaque groupe tire un rôle : 1. \zh{酋长/族长} (Chef coutumier), 2. \zh{法官/律师} (Juge/Avocat), 3. \zh{女儿/寡妇} (Fille/Veuve), 4. \zh{长子/继承人} (Aîné/Héritier principal).
\textbf{Sujet :} \zh{“祖传土地该不该分给女儿？”} (La terre ancestrale doit-elle être partagée aux filles ?)
\textbf{Déroulement (15 min) :}
\begin{enumerate}
    \item Préparation arguments (3 min) : Chaque rôle note 2 arguments en mots-clés chinois.
    \item Débat (8 min) : Tour de table, puis discussion libre. Contrainte : utiliser au moins 3 structures de comparaison (\zh{比}, \zh{而}, \zh{跟...不一样}) et 5 mots du vocabulaire.
    \item Synthèse (4 min) : Un rapporteur par groupe résume en 30s \textbf{en chinois} : \zh{我们组认为……因为……所以建议……}.
\end{enumerate}
\textbf{Évaluation :} Prononciation (tons), vocabulaire précis, structures de comparaison, cohérence argumentative.
\end{experience}

\begin{savaistu}[Le saviez-vous ? : Évolution récente du Code civil chinois (2021) et Droit camerounais]
\begin{itemize}
    \item \textbf{Chine -- Code civil 2021 (民法典) :} L'« Édition des successions » (继承编) innove : 1. \textbf{Testament tolérant (宽容遗嘱)} : un testament non parfait sur la forme (ex. témoin unique, écriture approximative) peut être valide si la volonté est claire. 2. \textbf{Pardon testamentaire (遗赦扶养协议)} : contrat par lequel un héritier s'engage à soigner le testateur contre legs. 3. \textbf{Partage inégalitaire justifié} : « Celui qui a assumé la charge principale d'entretien (\zh{尽主要赡养义务}) ou a vécu avec le défunt peut recevoir davantage » (Art. 1130).
    \item \textbf{Cameroun -- Loi n°2016/007 (Code civil) :} L'article 745 pose le principe : « La succession s'ouvre par la mort... la loi ne fait pas de distinction de sexe ». Cependant, l'article 757 prévoit que « les coutumes non contraires à l'ordre public s'appliquent ». C'est la base juridique du \zh{双轨制}.
    \item \textbf{Comparatif :} La Chine a \textbf{aboli} la primogéniture par la loi unique (1985). Le Cameroun \textbf{tolère} la coutume sous réserve d'ordre public (jurisprudence \textit{Essomba c. Essomba}, Cour Suprême, 2000+ : protection veuve/filles).
\end{itemize}
\end{savaistu}

\begin{methode}[Préparer et présenter un mini-exposé oral (1--2 min) : structure, notes, prononciation]
\textbf{Objectif :} Structurer une prise de parole continue en chinois sur le thème comparatif, sans lire un texte écrit.
\begin{enumerate}
    \item \textbf{Choisir un angle précis} (pas tout le cours) : ex. « Le rôle du notaire / du chef de village », « L'héritage des terres vs l'héritage des logements », « Le sort de la veuve », « Ville vs Campagne ».
    \item \textbf{Construire le canevas (mots-clés seulement, pas de phrases complètes)} :
        \begin{center}
        \begin{tikzpicture}[node distance=1.1cm, every node/.style={align=center, font=\small, inner sep=3pt}]
        \node[draw, rounded corners, fill=popblueL, text width=4.5cm] (titre) {Titre : \zh{中喀遗产继承对比} \\ (Comparaison Sino-Camerounaise)};
        \node[draw, rounded corners, fill=popgreenL, text width=3.5cm, below=of titre] (intro) {Intro : \zh{主题、两国背景、相似起点}};
        \node[draw, rounded corners, fill=poporangeL, text width=3.5cm, below left=of intro] (cameroun) {Cameroun : \zh{习惯法、长子、土地、酋长调解、双轨制}};
        \node[draw, rounded corners, fill=poporangeL, text width=3.5cm, below right=of intro] (chine) {Chine : \zh{法定继承、公证、房产、分家析产、立法统一}};
        \node[draw, rounded corners, fill=poppinkL, text width=4cm, below=of cameroun] (comp) {Comparer : \zh{相同点：向平等}\\\zh{不同点：路径、执行}};
        \node[draw, rounded corners, fill=popgoldL, text width=4cm, below=of comp] (concl) {Conclusion : \zh{趋势、建议、我观点}};
        \draw[->, thick, popblue] (titre) -- (intro);
        \draw[->, thick, popblue] (intro) -- (cameroun);
        \draw[->, thick, popblue] (intro) -- (chine);
        \draw[->, thick, popblue] (cameroun) -- (comp);
        \draw[->, thick, popblue] (chine) -- (comp);
        \draw[->, thick, popblue] (comp) -- (concl);
        \end{tikzpicture}
        \end{center}
    \item \textbf{Rédiger les "amorces" orales} (phrases toutes faites pour démarrer/transitionner) :
        \zh{大家好，今天我谈谈……} (Bonjour, aujourd'hui je parle de...)
        \zh{首先，在喀麦隆……} (D'abord, au Cameroun...)
        \zh{相比之下，在中国……} (En comparaison, en Chine...)
        \zh{共同点是……不同点是……} (Point commun... Différence...)
        \zh{最后，我认为……谢谢大家。} (Enfin, je pense que... Merci à tous.)
    \item \textbf{Entraînement tonal ciblé} : Identifier les mots à risque et les lire 3 fois.
        \begin{itemize}
            \item \zh{继承} (jì-chéng, 4-2) vs \zh{结成} (jié-chéng).
            \item \zh{土地} (tǔ-dì, 3-4) -- 3e ton + 4e ton.
            \item \zh{公证} (gōng-zhèng, 1-4).
            \item \zh{习惯法} (xí-guàn-fǎ, 2-4-3).
            \item \zh{双轨制} (shuāng-guǐ-zhì, 1-3-4).
        \end{itemize}
    \item \textbf{Gestion du temps (90s visés) :} Intro 15s -- Corps (Cameroun 25s + Chine 25s + Comparaison 15s) -- Conclusion 10s. Chronométrer à l'entraînement.
\end{enumerate}
\end{methode}

\begin{exoresolu}[Simulation d'exposé : notes et script oral]
\textbf{Énoncé :} À partir du canevas de la méthode, rédigez les \textbf{notes de l'orateur} (mots-clés + connecteurs en caractères) et le \textbf{script oral complet} (phrases en caractères + pinyin) pour un exposé de 90 secondes intitulé : \zh{“城市与农村：遗产分配的两张脸”} (Ville et campagne : deux visages du partage successoral).
\tcblower
\textbf{Solution détaillée :}

\textbf{1. Notes de l'orateur (fiche bristol -- mots-clés) :}
\begin{itemize}
    \item \zh{题目：城市 vs 农村 -- 遗产分配两张脸}
    \item \zh{开场：中喀同题：传统重男 $\rightarrow$ 现代重法}
    \item \zh{中国城市：公证遗嘱、房产、男女平等、份额均等}
    \item \zh{中国农村：宅基地、分家析产、儿子/长子优先、主宅}
    \item \zh{喀麦隆城市：民法典、法院判例、保护寡妇女儿}
    \item \zh{喀麦隆农村：酋长调解、习惯法、祖传土地归长子、女儿仅用权}
    \item \zh{核心对比：法律统一 (中国) vs 双轨并行 (喀麦隆)；书面遗嘱 vs 口头调解}
    \item \zh{结语：趋势向平等，需普法教育、法律援助下乡。}
\end{itemize}

\textbf{2. Script oral complet (≈ 230 caractères = ~90s débit normal) :}
\zh{大家好。今天我的话题是“城市与农村：遗产分配的两张脸”。} \\
\zh{在中国和喀麦隆，城乡差异都非常明显。} \\
\zh{在中国城市，家庭通常用公证遗嘱分割房产，法律保障男女平等，女儿和儿子拿一样的份额。} \\
\zh{但在中国农村，宅基地继承常按“分家析产”办，儿子尤其是长子往往得到主宅，女儿嫁出去就少分或不分。} \\
\zh{喀麦隆也是这样。城市里法院依据民法典判决，保护寡妇和女儿的权利。} \\
\zh{而在村里，酋长按习惯法调解，祖传土地通常归长子，女儿只有使用权没有所有权。} \\
\zh{相比之下，中国靠一部统一法律管全国，喀麦隆是“双轨制”并存。} \\
\zh{但是，两国趋势都是走向男女平等。} \\
\zh{建议加强普法宣传，让农村妇女知道自己的权利。谢谢大家。}

\textbf{Pinyin (clé de lecture) :}
Dàjiā hǎo. Jīntiān wǒ de huàtí shì "Chéngshì yǔ nóngcūn: yíchǎn fēnpèi de liǎng zhāng liǎn".
Zài Zhōngguó hé Kāmàilóng, chéngxiāng chāyì dōu fēi cháng míngxiǎn.
Zài Zhōngguó chéngshì, jiātíng tōngcháng yòng gōngzhèng yízhǔ fēngē fángchǎn, fǎlǜ bǎozhàng nánnǚ píngděng, nǚ'ér hé érzi ná yīyàng de fèn'é.
Dàn zài Zhōngguó nóngcūn, zhájīdì jìchéng cháng àn "fēnjiā xīchǎn" bàn, érzi yóuqí shì zhǎngzǐ wǎngwǎng dédào zhǔzhái, nǚ'ér jiàchūqù jiù shǎofēn huò bùfēn.
Kāmàilóng yě shì zhèyàng. Chéngshì lǐ fǎyuàn yījù Mínfǎdiǎn pànjué, bǎohù guǎfù hé nǚ'ér de quánlì.
Ér zài cūn lǐ, qiúzhǎng àn xíguànfǎ tiáojiě, zǔchuán tǔdì tōngcháng guī zhǎngzǐ, nǚ'ér zhǐ yǒu shǐyòngquán méiyǒu suǒyǒuquán.
Xiāngbǐ zhīxià, Zhōngguó kào yī bù tǒngyī fǎlǜ guǎn quánguó, Kāmàilóng shì "shuāngguīzhì" bìngcún.
Dànshì, liǎng guó qūshì dōu shì zǒuxiàng nánnǚ píngděng.
Jiànyì jiāqiáng pǔfǎ xuānchuán, ràng nóngcūn fùnǚ zhīdào zìjǐ de quánlì. Xièxie dàjiā.
\end{exoresolu}

\begin{attention}[Les 5 erreurs qui coûtent des points au Bac -- Focus Héritage]
\begin{enumerate}
    \item \textbf{Tons : confusion \cle{jì} (4e) / \cle{jǐ} (3e) / \cle{qí} (2e).} \zh{继承} (jìchéng, hériter) $\neq$ \zh{纪成} (jì chéng) $\neq$ \zh{奇成}. \zh{土地} (tǔdì, 3-4) $\neq$ \zh{兔弟} (tùdì). \textbf{Réflexe :} dictionnaire audio + répétition chorale sur \zh{继承权、土地、份额}.
    \item \textbf{Caractères proches : \cle{遗} (yí, legs, rad. \zh{辶}) vs \cle{疑} (yí, doute, rad. \zh{疒}) ; \cle{承} (chéng, porter, rad. \zh{手}) vs \cle{成} (chéng, réussir, rad. \zh{戈}).} \zh{遗产} (yíchǎn) $\neq$ \zh{疑产}. \zh{继承} (jìchéng) $\neq$ \zh{继成}. \textbf{Astuce :} analyser le radical (\zh{辶} = laisser derrière ; \zh{手} = porter la charge).
    \item \textbf{Ordre des mots : Place du complément de lieu/temps (Règle : Sujet + Temps + Lieu + Verbe).} \textbf{Erreur :} \zh{我继承了土地在喀麦隆。} (Calque français). \textbf{Correct :} \zh{我在喀麦隆继承了祖传土地。} (Wǒ zài Kāmàilóng jìchéng le zǔchuán tǔdì). Le lieu \zh{在喀麦隆} précède le verbe.
    \item \textbf{Mots-outils : \cle{了} (aspect accompli/réalisé) vs \cle{过} (expérience vécue).} \zh{他立了遗嘱} (Tā lì le yízhǔ -- Il a fait son testament [action terminée, focus sur le fait accompli]). \zh{他立过遗嘱} (Tā lì guo yízhǔ -- Il a déjà fait un testament [expérience de vie, peut-être révoqué]). \textbf{Piège Bac :} Narration = \zh{了} ; Question « As-tu déjà... ? » / « Il a déjà fait... » = \zh{过}.
    \item \textbf{Classificateurs : \cle{份} (fèn, parts, documents, héritage) vs \cle{项} (xiàng, projets, clauses, droits) vs \cle{笔} (bǐ, sommes d'argent) vs \cle{位} (wèi, personnes honorifiques).} \zh{一份遗嘱} (yī fèn yízhǔ), \zh{一份遗产} (yī fèn yíchǎn), \zh{一笔巨款} (yī bǐ jùkuǎn), \zh{一项权利} (yī xiàng quánlì), \zh{一位继承人} (yī wèi jìchéngrén). \textbf{Erreur fréquente :} \zh{一个遗产} (yī gè yíchǎn). \textbf{Méthode :} apprendre le couple \textbf{Nom + Classificateur} par cœur (fiche vocab).
\end{enumerate}
\end{attention}

\courssub{Exercices d'entraînement (Devoir / Classe)}

\begin{exercice}[Expression écrite : Comparaison argumentée]
\textbf{Consigne :} Rédigez un paragraphe de 100 à 120 caractères en chinois sur le thème : \zh{“如何看待城乡继承差异”} (Comment voir la différence successorale ville/campagne).
\textbf{Contraintes :}
\begin{enumerate}
    \item Utiliser \textbf{au moins 3 structures de comparaison} (\zh{比}, \zh{而}, \zh{跟...一样/不一样}, \zh{相比之下}).
    \item Intégrer \textbf{au moins 6 mots du tableau de vocabulaire} (surlignez-les dans votre copie).
    \item Structure : Thèse -- Argument 1 (Ville) -- Argument 2 (Campagne) -- Conclusion/Proposition.
    \item Soigner la ponctuation chinoise (， 。 、 ；).
\end{enumerate}
\end{exercice}

\begin{exercice}[Traduction : Thème (Français $\rightarrow$ Chinois)]
\textbf{Consigne :} Traduisez les phrases suivantes en chinois. Attention aux structures de comparaison, aux classificateurs et aux mots-outils (\zh{了}/\zh{过}).
\begin{enumerate}
    \item Bien que le Cameroun ait une loi moderne, la coutume reste forte dans les villages.
    \item La Chine a aboli la primogéniture féodale il y a longtemps (用 \zh{早已}), mais la campagne conserve l'habitude de « division familiale ».
    \item Cette veuve n'a pas hérité de la terre ancestrale, car le chef a jugé selon la coutume.
    \item Avez-vous déjà fait un testament notarié ? (Utiliser \zh{过}).
    \item Selon la loi, les filles ont le même droit de succession que les fils.
\end{enumerate}
\end{exercice}

\begin{exercice}[Compréhension orale / Expression orale : Préparation Bac]
\textbf{Consigne :} Préparez un exposé de 1 minute 30 (environ 250 caractères) sur l'un des deux sujets au choix. Vous remettrez vos \textbf{notes (canevas)} et enregistrerez votre \textbf{script oral} (ou le présenterez en classe).
\begin{enumerate}
    \item \zh{题目一：寡妇的继承权：中喀法律保护的对比} (Droits successoraux des veuves : comparaison de la protection légale Chine/Cameroun).
    \item \zh{题目二：遗嘱制度的现代化：从手写到数字化} (Modernisation du testament : de l'olographe au numérique -- Chine \zh{民法典} 2021 vs Cameroun).
\end{enumerate}
\textbf{Critères d'évaluation (Grille Bac) :}
\begin{itemize}
    \item Contenu (Faits exacts, comparaison claire) : 5 pts.
    \item Vocabulaire (Précision, richesse, 8+ mots du thème) : 4 pts.
    \item Grammaire (Structures comparaison, ordre mots, \zh{了}/\zh{过}) : 4 pts.
    \item Prononciation / Tons / Fluidité : 4 pts.
    \item Cohérence / Connecteurs logiques : 3 pts.
\end{itemize}
\end{exercice}

\begin{corrige}[Exercice : Traduction (Thème)]
\begin{enumerate}
    \item \zh{虽然喀麦隆有现代法律，但是在村里习惯法依然很强势。} (Suīrán Kāmàilóng yǒu xiàndài fǎlǜ, dànshì zài cūn lǐ xíguànfǎ yīrán hěn qiángshì.) -- \textit{Structure concession \zh{虽然...但是...} ; lieu avant verbe.}
    \item \zh{中国早已废除了封建长子继承制，但农村仍保留“分家析产”的习惯。} (Zhōngguó zǎiyǐ fèichú le fēngjiàn zhǎngzǐ jìchéngzhì, dàn nóngcūn réng bǎoliú "fēnjiā xīchǎn" de xíguàn.) -- \textit{\zh{早已} + \zh{了} (accompli passé) ; \zh{但} transition.}
    \item \zh{这位寡妇没有继承祖传土地，因为酋长按习惯法判决/调解。} (Zhè wèi guǎfù méiyǒu jìchéng zǔchuán tǔdì, yīnwèi qiúzhǎng àn xíguànfǎ pànjué/tiáojiě.) -- \textit{\zh{这位} (classif. honorifique) ; \zh{没有} + V (négation passé) ; \zh{按} (selon/conformément à).}
    \item \zh{你立过公证遗嘱吗？} (Nǐ lì guo gōngzhèng yízhǔ ma?) -- \textit{\zh{过} pour l'expérience de vie.}
    \item \zh{根据法律，女儿跟儿子享有同样的继承权。} (Gēnjù fǎlǜ, nǚ'ér gēn érzi xiǎngyǒu tóngyàng de jìchéngquán.) -- \textit{\zh{根据} (selon/sur la base de) ; \zh{跟...一样} structure égalité.}
\end{enumerate}
\end{corrige}

\begin{corrige}[Exercice : Expression écrite -- Proposition de correction type]
\textbf{Sujet :} \zh{如何看待城乡继承差异}
\textbf{Proposition :}
\zh{我认为城乡继承差异反映了法治普及的不平衡。一方面，城市依靠公证遗嘱和法定继承，男女份额均等，保障了妇女权益。另一方面，农村受传统习惯法影响，宅基地常由长子继承，出嫁女往往没有份额，这跟法律规定不一样。相比之下，中国立法统一性强，喀麦隆双轨制并存，但两国趋势都向男女平等。建议加强普法宣传和法律援助下乡，让法律真正走进村庄，实现公平正义。}
(≈ 115 caractères)
\textbf{Vocabulaire ciblé :} \zh{法治, 公证遗嘱, 法定继承, 份额均等, 妇女权益, 习惯法, 宅基地, 长子, 出嫁女, 双轨制, 普法宣传, 法律援助, 公平正义}.
\end{corrige}

\newpage

\coursec{Exercices}

\courssub{Je vérifie mes connaissances}

\begin{exercice}[QCM : vocabulaire et grammaire du thème]
Chaque item ne comporte qu'une seule réponse correcte. Choisis la bonne proposition.
\begin{enumerate}
    \item Quel mot signifie « testament » ?
    \begin{itemize}
        \item[A)] 遗产 (yíchǎn)
        \item[B)] 遗嘱 (yízhǔ)
        \item[C)] 继承 (jìchéng)
        \item[D)] 继承人 (jìchéngrén)
    \end{itemize}
    \item Pour dire « Au Cameroun, l'héritage suit la tradition », on utilise la structure :
    \begin{itemize}
        \item[A)] 喀麦隆的遗产 \textbf{跟} 传统 一样。
        \item[B)] 喀麦隆的遗产 \textbf{按} 传统 分配。
        \item[C)] 喀麦隆的遗产 \textbf{比} 传统 重要。
        \item[D)] 喀麦隆的遗产 \textbf{给} 传统。
    \end{itemize}
    \item La particule marquant le complément d'agent dans une phrase passive (comme « l'héritage \textbf{est divisé par} la loi ») est :
    \begin{itemize}
        \item[A)] 被
        \item[B)] 把
        \item[C)] 给
        \item[D)] 让
    \end{itemize}
    \item « Les fils et les filles ont les \textbf{mêmes} droits » se traduit par :
    \begin{itemize}
        \item[A)] 儿子和女儿有 \textbf{不一样} 的权利。
        \item[B)] 儿子和女儿有 \textbf{一样} 的权利。
        \item[C)] 儿子 \textbf{比} 女儿 有 权利。
        \item[D)] 儿子 \textbf{跟} 女儿 \textbf{不} 一样 权利。
    \end{itemize}
\end{enumerate}
\end{exercice}

\begin{exercice}[Vrai ou Faux ? Justifie en français]
Lis les affirmations suivantes sur les pratiques successorales. Coche V (Vrai) ou F (Faux) et justifie brièvement en te basant sur les faits comparés vus en cours.
\begin{enumerate}
    \item [V / F] En Chine, depuis le Code civil de 2021, les filles ont légalement les mêmes droits d'héritage que les fils.
    \item [V / F] Au Cameroun, la coutume prévoit systématiquement que l'aîné des fils hérite de la totalité du patrimoine (terres, maisons, argent) à l'exclusion des filles.
    \item [V / F] La rédaction d'un testament (遗嘱) est une pratique courante et juridiquement contraignante en Chine moderne, mais reste rare dans les zones rurales camerounaises.
    \item [V / F] Le terme 继承人 (jìchéngrén) désigne exclusivement l'héritier désigné par la coutume, jamais par la loi.
\end{enumerate}
\end{exercice}

\begin{exercice}[Vocabulaire : associe les caractères, le pinyin et la traduction]
Relie chaque élément de la colonne A (caractères) à son pinyin (colonne B) et à sa traduction (colonne C). Écris les numéros correspondants.

\begin{center}
\begin{tabularx}{\linewidth}{|X|X|X|X|}
\hline
\rowcolor{popblueL} \textbf{A. Caractères} & \textbf{B. Pinyin} & \textbf{C. Traduction} & \textbf{Ton (1-4)} \\
\hline
1. 遗产 & a. fǎdìng & a. héritier & 1. \_\_\_ \\
2. 遗嘱 & b. jìchéngrén & b. héritage / patrimoine & 2. \_\_\_ \\
3. 法定 & c. yíchǎn & c. testament & 3. \_\_\_ \\
4. 继承人 & d. yízhǔ & d. légal / statutaire & 4. \_\_\_ \\
\hline
\end{tabularx}
\end{center}
\end{exercice}

\begin{exercice}[Pinyin et tons : écris le pinyin avec les marques de tons]
Donne le pinyin complet (avec accents toniques) pour les mots suivants. Attention aux changements de ton (砂纸 3-3 $\to$ 2-3) et au « ü ».
\begin{enumerate}
    \item 分配 \_\_\_\_\_\_\_\_\_\_
    \item 男女平等 \_\_\_\_\_\_\_\_\_\_
    \item 习俗 \_\_\_\_\_\_\_\_\_\_
    \item 选择 \_\_\_\_\_\_\_\_\_\_
    \item 公平 \_\_\_\_\_\_\_\_\_\_
    \item 长子 \_\_\_\_\_\_\_\_\_\_
\end{enumerate}
\end{exercice}

\courssub{Je m'entraîne}

\begin{exercice}[Traduction : du français vers le chinois]
Traduis les phrases suivantes en chinois (caractères + pinyin). Utilise le vocabulaire de la leçon (\emph{遗产, 遗嘱, 法定, 习俗, 比, 一样, 不如, 觉得}).
\begin{enumerate}
    \item Au Cameroun, l'héritage est souvent partagé selon la coutume.
    \item En Chine, la loi stipule que les fils et les filles ont les mêmes droits.
    \item Je pense que l'égalité homme-femme dans l'héritage est plus juste.
    \item Mon père a écrit un testament pour choisir l'héritier.
    \item La tradition camerounaise est différente de la loi chinoise.
\end{enumerate}
\end{exercice}

\begin{exercice}[Transformation de phrases : structures comparatives]
Réécris chaque phrase en utilisant le mot ou la structure indiquée entre parenthèses, sans changer le sens général.
\begin{enumerate}
    \item 中国的法律规定儿子和女儿权利相同。 (用 \textbf{一样})
    \item 喀麦隆的习俗跟中国的法律不一样。 (用 \textbf{比} ... \textbf{不同} ou \textbf{不如})
    \item 女儿也可以继承遗产。 (用 \textbf{也} ... \textbf{有权利} / \textbf{能})
    \item 很多人不写遗嘱。 (用 \textbf{没有} ... \textbf{写})
    \item 我觉得公平很重要。 (用 \textbf{认为} / \textbf{觉得} ... \textbf{重要})
\end{enumerate}
\end{exercice}

\begin{exercice}[Texte à trous : mots-outils et particules]
Complète le texte suivant avec les mots manquants : \textbf{的, 得, 地, 比, 和, 一样, 不如, 按照, 把, 被}.

\begin{textebox}[L'héritage au Cameroun et en Chine : comparaison]
中国 \_\_\_\_\_\_ 喀麦隆 \_\_\_\_\_\_ 继承 方式 \_\_\_\_\_\_ \_\_\_\_\_\_。\\
中国 \_\_\_\_\_\_ 法律 规定，儿子 \_\_\_\_\_\_ 女儿 权利 \_\_\_\_\_\_。\\
喀麦隆 很多 地方 \_\_\_\_\_\_ 习俗 办事，女儿 往往 \_\_\_\_\_\_ 儿子 受 重视。\\
现在，越来越 多 人 选择 写 遗嘱 \_\_\_\_\_\_ 分配 遗产，让 事情 办 \_\_\_\_\_\_ 更 公平。
\end{textebox}
\end{exercice}

\begin{exercice}[Remise en ordre et appariement]
\begin{enumerate}
    \item Remets les mots dans l'ordre pour former une phrase correcte.
    \begin{enumerate}
        \item 法律 / 规定 / 遗产 / 分配 / 公平 / 必须
        \item 遗嘱 / 父亲 / 留下 / 一份 / 给 / 家人
        \item 习俗 / 传统 / 虽然 / 重要 / 法律 / 但是 / 更 / 有 / 效力
    \end{enumerate}
    \item Apparie la première partie de la phrase (1-4) à sa fin logique (a-d).
    \begin{tabular}{ll}
    1. 如果 没有 遗嘱， & a. 女儿 也能 继承 遗产。 \\
    2. 中国 法律 规定， & b. 遗产 按 法定 顺序 分配。 \\
    3. 在 很多 传统 里， & c. 长子 往往 是 主要 继承人。 \\
    4. 写 遗嘱 可以 让， & d. 分配 更 符合 死者 意愿。
    \end{tabular}
\end{enumerate}
\end{exercice}

\courssub{Je me prépare au Bac}

\begin{exercice}[Compréhension de l'écrit : texte original et questions]
Lis le texte suivant attentivement, puis réponds aux questions \textbf{en français}.

\begin{textebox}[La succession dans la Chine contemporaine : entre loi et coutume]
中国 的 《民法典》 于 2021 年 实施，明确 规定 遗产 继承 坚持 男女 平等 原则。法定 继承 中，第一 顺序 继承人 包括 配偶、子女 和 父母。这里 的 “子女”，不 分 婚生、非婚生、养子女 和 继子女，权利 完全 一样。

然而，在 农村 地区 或 一些 传统 家庭，重男轻女 的 观念 依然 存在。有些 父母 认为 只有 儿子 能 “传宗接代”，所以 想 把 房子 和 土地 留给 儿子。如果 立 了 遗嘱，法院 一般 尊重 遗嘱 内容；但 如果 遗嘱 剥夺 了 无劳动 能力 且 无 生活 来源 的 女儿 继承 权，法院 可能 判决 遗嘱 部分 无效，保留 她 的 必留份。

现在，越来越 多 的 年轻 人 选择 通过 公证 立 遗嘱，或者 购买 保险、设立 信托 来 安排 遗产。这 说明 法治 观念 正在 深入 人心，但 改变 根深蒂固 的 习俗 还 需要 时间。
\end{textebox}

\begin{enumerate}
    \item Selon le texte, que stipule le Code civil chinois (民法典) de 2021 concernant l'ordre des héritiers et l'égalité des sexes ?
    \item Quelles sont les catégories de « enfants » (子女) qui ont des droits identiques selon la loi ?
    \item Pourquoi certains parents en zone rurale veulent-ils encore léguer leurs biens uniquement à leur fils ? Cite le terme chinois utilisé dans le texte pour cette mentalité.
    \item Dans quel cas précis un tribunal peut-il déclarer un testament « partiellement invalide » (部分 无效) ?
    \item Quelles sont les nouvelles pratiques mentionnées à la fin du texte pour organiser sa succession ? Que cela révèle-t-il sur l'évolution de la société chinoise ?
\end{enumerate}
\end{exercice}

\begin{exercice}[Thème et Version]
\begin{enumerate}
    \item \textbf{Thème (Français $\to$ Chinois) :} Traduis le paragraphe suivant en chinois (caractères + pinyin). \emph{Visée : 80-100 caractères.}
    \begin{quote}
    « Au Cameroun, la coutume choisit souvent l'héritier principal. En Chine, la loi protège les droits des filles. Je pense que l'égalité est meilleure que la tradition. Écrire un testament est une bonne méthode pour éviter les conflits. »
    \end{quote}
    \item \textbf{Version (Chinois $\to$ Français) :} Traduis le texte suivant en français.
    \begin{textebox}[Texte à traduire]
    现在 的 年轻人 既 尊重 传统，又 相信 法律。他们 觉得，无论 是 儿子 还是 女儿，都 应该 照顾 父母，也 都 有 权利 继承 遗产。如果 家里 人 能 坐 下来 好好 谈 一 谈，就 不会 吵架 了。
    \end{textebox}
\end{enumerate}
\end{exercice}

\begin{exercice}[Expression écrite : Mini-exposé comparatif]
\textbf{Sujet :} « \emph{Les filles doivent-elles hériter comme les fils ? Comparez la situation au Cameroun et en Chine et donnez votre avis personnel.} »

\textbf{Consignes :}
\begin{itemize}
    \item Rédige un court exposé structuré (introduction, développement comparé, conclusion/avis).
    \item Longueur visée : \textbf{120 à 150 caractères chinois} (ponctuation comprise).
    \item Utilise \textbf{au moins 5 mots} du vocabulaire de la leçon : \emph{遗产, 遗嘱, 法定, 习俗, 男女平等, 分配, 权利, 传统, 公平, 选择}.
    \item Utilise \textbf{au moins 3 structures comparatives} : \emph{比, 跟...一样/不一样, 不如, 觉得/认为...}.
    \item Soigne l'ordre des mots (Sujet + Temps + Lieu + Manière + Verbe + Complément) et les particules \emph{的/得/地}.
\end{itemize}

\begin{textebox}[Espace de rédaction (brouillon)]
\begin{center}
\begin{tabular}{|p{0.9\linewidth}|}
\hline
\vspace{6cm} \\
\hline
\end{tabular}
\end{center}
\end{textebox}
\end{exercice}

\newpage

\coursec{Corrigés des exercices}

\begin{corrige}[Exercice 1 — QCM : vocabulaire et grammaire du thème]
\textbf{Réponses :}
\begin{enumerate}
    \item \textbf{B)} \zh{遗嘱} (yízhǔ) = testament. \zh{遗产} (yíchǎn) = héritage/patrimoine ; \zh{继承} (jìchéng) = hériter ; \zh{继承人} (jìchéngrén) = héritier.
    \item \textbf{B)} \zh{喀麦隆的遗产按传统分配。} (Kāmàilóng de yíchǎn àn chuántǒng fēnpèi.) « Au Cameroun, l'héritage est partagé selon la tradition. » La préposition \zh{按} (àn) signifie « selon, conformément à » ; c'est la structure attendue. \zh{跟...一样} exprime la similitude, \zh{比} la comparaison de supériorité, \zh{给} le bénéficiaire : aucune ne convient au sens.
    \item \textbf{A)} \zh{被} (bèi). Structure passive : \cle{Objet + 被 + agent + verbe}. Exemple : \zh{遗产被法律分配。} (Yíchǎn bèi fǎlǜ fēnpèi.) « L'héritage est réparti par la loi. » \zh{把} est la particule de l'objet (construction active), \zh{让} et \zh{给} peuvent marquer le passif à l'oral mais la réponse canonique du programme est \zh{被}.
    \item \textbf{B)} \zh{儿子和女儿有一样的权利。} (Érzi hé nǚ'ér yǒu yíyàng de quánlì.) « Les fils et les filles ont les mêmes droits. » Structure : \cle{A 和 B + 一样}. L'option A dit le contraire (\zh{不一样}), C est une comparaison de supériorité, D est grammaticalement fautive (il manque \zh{的} et la structure est brisée).
\end{enumerate}
\textbf{Points de vigilance :} ne pas confondre \zh{遗产} (le bien transmis) et \zh{遗嘱} (le document) ; dans \zh{一样}, l'adjectif \zh{一} prend le 4\textsuperscript{e} ton devant un 4\textsuperscript{e} ton (\emph{yíyàng}).
\end{corrige}

\begin{corrige}[Exercice 2 — Vrai ou Faux ?]
\begin{enumerate}
    \item \textbf{Vrai.} Le Code civil chinois (民法典), entré en vigueur en 2021, consacre le principe d'égalité homme-femme (男女平等) en matière de succession : fils et filles ont légalement les mêmes droits d'héritage.
    \item \textbf{Faux.} « Systématiquement » est exagéré : les coutumes camerounaises sont très diverses selon les ethnies et les régions (systèmes patrilinéaires et matrilinéaires). Souvent l'aîné ou un héritier désigné administre les biens, mais il n'y a pas de règle unique excluant totalement les filles partout ; de plus, le droit moderne camerounais reconnaît les droits des enfants des deux sexes.
    \item \textbf{Vrai.} En Chine, le testament notarié (公证遗嘱) est juridiquement contraignant et de plus en plus courant en ville. Au Cameroun rural, la transmission se fait encore majoritairement par la coutume orale, et le testament écrit reste rare.
    \item \textbf{Faux.} \zh{继承人} (jìchéngrén) désigne tout héritier, qu'il soit \emph{légal} (法定继承人, désigné par la loi) ou \emph{testamentaire} (遗嘱继承人, désigné par testament), et non exclusivement l'héritier coutumier.
\end{enumerate}
\textbf{Point de méthode :} dans une justification, cite toujours un fait précis du cours (loi, date, pratique) plutôt qu'une impression générale.
\end{corrige}

\begin{corrige}[Exercice 3 — Association caractères / pinyin / traduction]
\textbf{Correspondances :}
\begin{itemize}
    \item 1. \zh{遗产} $\to$ \textbf{c. yíchǎn} $\to$ \textbf{b. héritage / patrimoine}. Tons : \textbf{2-3}.
    \item 2. \zh{遗嘱} $\to$ \textbf{d. yízhǔ} $\to$ \textbf{c. testament}. Tons : \textbf{2-3}.
    \item 3. \zh{法定} $\to$ \textbf{a. fǎdìng} $\to$ \textbf{d. légal / statutaire}. Tons : \textbf{3-4}.
    \item 4. \zh{继承人} $\to$ \textbf{b. jìchéngrén} $\to$ \textbf{a. héritier}. Tons : \textbf{4-2-2}.
\end{itemize}
\textbf{Points de vigilance :} attention au \zh{ü} de \zh{女} (nǚ) qui apparaît dans les mots du thème ; après \zh{j, q, x}, le « ü » s'écrit sans tréma mais se prononce toujours [ü] (ex. \zh{继承} jìchéng ne contient pas de ü, mais \zh{女儿} nǚ'ér oui). Distingue bien \zh{遗} (yí, « laisser ») présent dans \zh{遗产} et \zh{遗嘱}.
\end{corrige}

\begin{corrige}[Exercice 4 — Pinyin et tons]
\begin{enumerate}
    \item \zh{分配} : \textbf{fēnpèi} (1-4)
    \item \zh{男女平等} : \textbf{nánnǚ píngděng} (2-3 / 2-3). Attention : \zh{女} = nǚ avec tréma.
    \item \zh{习俗} : \textbf{xísú} (2-2)
    \item \zh{选择} : \textbf{xuǎnzé} (3-2). Règle 3-3 $\to$ 2-3 : ici le deuxième caractère est au 2\textsuperscript{e} ton, donc \emph{pas} de changement de ton ; on garde xuǎn au 3\textsuperscript{e} ton.
    \item \zh{公平} : \textbf{gōngpíng} (1-1)
    \item \zh{长子} : \textbf{zhǎngzǐ} (3-3). Deux 3\textsuperscript{e} tons se suivent : le premier passe au 2\textsuperscript{e} ton à l'oral (\emph{zhángzǐ}), mais on \emph{écrit} les tons d'origine : zhǎngzǐ.
\end{enumerate}
\textbf{Rappel de la règle :} quand deux syllabes au 3\textsuperscript{e} ton se suivent, la première se prononce au 2\textsuperscript{e} ton, mais la transcription pinyin conserve la marque du 3\textsuperscript{e} ton. Exemple classique : \zh{你好} s'écrit nǐ hǎo, se prononce ní hǎo.
\end{corrige}

\begin{corrige}[Exercice 5 — Traduction : du français vers le chinois]
\textbf{Traductions proposées :}
\begin{enumerate}
    \item \zh{在喀麦隆，遗产常常按照习俗分配。}\\
    \emph{Zài Kāmàilóng, yíchǎn chángcháng ànzhào xísú fēnpèi.}\\
    Structure : lieu (\zh{在} + lieu) en tête, puis sujet + fréquence + \zh{按照} + verbe.
    \item \zh{在中国，法律规定儿子和女儿有一样的权利。}\\
    \emph{Zài Zhōngguó, fǎlǜ guīdìng érzi hé nǚ'ér yǒu yíyàng de quánlì.}\\
    « Stipuler » = \zh{规定} ; « les mêmes droits » = \zh{一样的权利} (ne pas oublier le \zh{的}).
    \item \zh{我觉得遗产中的男女平等更公平。} ou \zh{我认为在遗产方面，男女平等比传统更公平。}\\
    \emph{Wǒ juéde yíchǎn zhōng de nánnǚ píngděng gèng gōngpíng.}\\
    « Je pense que » = \zh{我觉得/我认为} + proposition complète (sans « que » en chinois).
    \item \zh{我爸爸写了遗嘱来选择继承人。} ou \zh{我父亲为了选择继承人写了一份遗嘱。}\\
    \emph{Wǒ bàba xiě le yízhǔ lái xuǎnzé jìchéngrén.}\\
    \zh{来} introduit le but ; on peut ajouter le classificateur : \zh{一份遗嘱} (« un testament »).
    \item \zh{喀麦隆的传统跟中国的法律不一样。}\\
    \emph{Kāmàilóng de chuántǒng gēn Zhōngguó de fǎlǜ bù yíyàng.}\\
    Structure de la différence : \cle{A 跟 B 不一样}. Négation placée avant \zh{一样}.
\end{enumerate}
\textbf{Points de vigilance :} ordre des mots (le lieu avec \zh{在} se place avant le verbe) ; \zh{一样} exige \zh{的} devant le nom ; pas de « que » après \zh{觉得}.
\end{corrige}

\begin{corrige}[Exercice 6 — Transformation de phrases]
\begin{enumerate}
    \item \zh{中国的法律规定，儿子和女儿有一样的权利。}\\
    \emph{Zhōngguó de fǎlǜ guīdìng, érzi hé nǚ'ér yǒu yíyàng de quánlì.}\\
    \zh{相同} devient \zh{一样} avec la structure \cle{A 和 B 一样}.
    \item \zh{喀麦隆的习俗跟中国的法律不同。} ou \zh{喀麦隆的习俗不如中国的法律那么明确。}\\
    \emph{Kāmàilóng de xísú gēn Zhōngguó de fǎlǜ bùtóng.}\\
    \zh{不一样} $\to$ \zh{不同} (synonyme, plus soutenu) ; avec \zh{不如} : \cle{A 不如 B + adjectif} (« A n'est pas aussi... que B »).
    \item \zh{女儿也有权利继承遗产。} ou \zh{女儿也能继承遗产。}\\
    \emph{Nǚ'ér yě yǒu quánlì jìchéng yíchǎn.}\\
    \zh{也} se place juste après le sujet, avant le verbe.
    \item \zh{很多人没有写遗嘱。}\\
    \emph{Hěn duō rén méiyǒu xiě yízhǔ.}\\
    Négation du passé/de l'accompli : \zh{没有} + verbe (sans \zh{了}).
    \item \zh{我认为公平很重要。} ou \zh{我觉得公平很重要。}\\
    \emph{Wǒ rènwéi gōngpíng hěn zhòngyào.}\\
    \zh{认为} est plus formel que \zh{觉得} ; l'adjectif prédicat exige un adverbe de degré (\zh{很}).
\end{enumerate}
\textbf{Point de vigilance :} avec \zh{没有}, ne jamais garder \zh{了} après le verbe (\zh{没有写了} est fautif).
\end{corrige}

\begin{corrige}[Exercice 7 — Texte à trous]
\textbf{Texte complété :}
\begin{itemize}
    \item \zh{中国} \textbf{和} \zh{喀麦隆} \textbf{的} \zh{继承方式} \textbf{不一样}。\\
    \emph{Zhōngguó hé Kāmàilóng de jìchéng fāngshì bù yíyàng.}\\
    « La Chine et le Cameroun ont des modes de succession différents. »
    \item \zh{中国} \textbf{的} \zh{法律规定，儿子} \textbf{和} \zh{女儿权利} \textbf{一样}。\\
    \emph{Zhōngguó de fǎlǜ guīdìng, érzi hé nǚ'ér quánlì yíyàng.}
    \item \zh{喀麦隆很多地方} \textbf{按照} \zh{习俗办事，女儿往往} \textbf{不如} \zh{儿子受重视。}\\
    \emph{Kāmàilóng hěn duō dìfang ànzhào xísú bànshì, nǚ'ér wǎngwǎng bùrú érzi shòu zhòngshì.}\\
    « Dans beaucoup de régions du Cameroun, on agit selon la coutume ; les filles sont souvent moins valorisées que les fils. »
    \item \zh{现在，越来越多人选择写遗嘱} \textbf{来} \zh{分配遗产，让事情办} \textbf{得} \zh{更公平。}\\
    \emph{Xiànzài, yuè lái yuè duō rén xuǎnzé xiě yízhǔ lái fēnpèi yíchǎn, ràng shìqing bàn de gèng gōngpíng.}
\end{itemize}
\textbf{Justifications grammaticales :}
\begin{itemize}
    \item \zh{的} : marque de possession/détermination (\zh{中国的法律}, \zh{喀麦隆的继承方式}).
    \item \zh{得} : complément de manière/degré après le verbe (\zh{办得更公平} : « que les choses soient réglées plus équitablement »).
    \item \zh{不如} : comparaison d'infériorité \cle{A 不如 B + verbe/adjectif}.
    \item \zh{按照} : préposition « selon, conformément à » + nom, placée avant le verbe.
    \item \zh{一样} / \zh{不一样} : égalité / différence, avec \zh{和} pour joindre les deux termes comparés.
\end{itemize}
\textbf{Point de vigilance :} \zh{得} (degré, après le verbe) $\neq$ \zh{的} (détermination, avant le nom) $\neq$ \zh{地} (adverbe, avant le verbe). Les mots \zh{比} et \zh{被} n'étaient pas nécessaires ici : vérifie toujours le sens global avant de placer une particule.
\end{corrige}

\begin{corrige}[Exercice 8 — Remise en ordre et appariement]
\textbf{1. Phrases remises en ordre :}
\begin{enumerate}
    \item \zh{法律规定遗产分配必须公平。}\\
    \emph{Fǎlǜ guīdìng yíchǎn fēnpèi bìxū gōngpíng.}\\
    « La loi stipule que la répartition de l'héritage doit être équitable. » Ordre : sujet (\zh{法律}) + verbe (\zh{规定}) + subordonnée (sujet \zh{遗产分配} + modal \zh{必须} + adjectif \zh{公平}).
    \item \zh{父亲给家人留下一份遗嘱。}\\
    \emph{Fùqīn gěi jiārén liúxià yí fèn yízhǔ.}\\
    « Le père laisse un testament à la famille. » Structure : sujet + \zh{给} + bénéficiaire + verbe \zh{留下} + classificateur \zh{一份} + objet.
    \item \zh{虽然传统习俗重要，但是法律更有效力。}\\
    \emph{Suīrán chuántǒng xísú zhòngyào, dànshì fǎlǜ gèng yǒu xiàolì.}\\
    « Bien que les coutumes traditionnelles soient importantes, la loi a plus de force. » Paire corrélative : \cle{虽然...但是...} (les deux se combinent en chinois, contrairement au français).
\end{enumerate}
\textbf{2. Appariement :}
\begin{itemize}
    \item 1 $\to$ \textbf{b} : \zh{如果没有遗嘱，遗产按法定顺序分配。} « S'il n'y a pas de testament, l'héritage est réparti selon l'ordre légal. »
    \item 2 $\to$ \textbf{a} : \zh{中国法律规定，女儿也能继承遗产。} « La loi chinoise stipule que les filles aussi peuvent hériter. »
    \item 3 $\to$ \textbf{c} : \zh{在很多传统里，长子往往是主要继承人。} « Dans beaucoup de traditions, le fils aîné est souvent l'héritier principal. »
    \item 4 $\to$ \textbf{d} : \zh{写遗嘱可以让分配更符合死者意愿。} « Rédiger un testament permet une répartition plus conforme à la volonté du défunt. »
\end{itemize}
\textbf{Point de vigilance :} le classificateur de \zh{遗嘱} est \zh{份} (\zh{一份遗嘱}) ; \zh{给} + bénéficiaire se place avant le verbe principal.
\end{corrige}

\begin{corrige}[Exercice 9 — Compréhension de l'écrit]
\textbf{Réponses rédigées en français :}
\begin{enumerate}
    \item Le Code civil chinois (民法典), entré en vigueur en 2021, consacre le principe d'égalité homme-femme (男女平等原则) en matière de succession. Dans la succession légale (法定继承), les héritiers de premier ordre (第一顺序继承人) sont le conjoint (配偶), les enfants (子女) et les parents (父母).
    \item Ont des droits parfaitement identiques (权利完全一样) : les enfants légitimes (婚生), les enfants nés hors mariage (非婚生), les enfants adoptés (养子女) et les enfants du conjoint / beaux-enfants (继子女).
    \item Certains parents ruraux veulent léguer la maison et la terre uniquement au fils parce qu'ils pensent que seul le fils peut « perpétuer la lignée ». Le texte nomme cette mentalité \zh{重男轻女} (zhòng nán qīng nǚ) : « valoriser les garçons et mépriser les filles ».
    \item Un tribunal peut déclarer un testament partiellement invalide (部分无效) lorsque celui-ci prive de son droit d'héritage une fille qui est à la fois incapable de travailler (无劳动能力) et sans source de revenus (无生活来源) ; le tribunal lui conserve alors sa « part obligatoire » (必留份).
    \item Les nouvelles pratiques sont : le testament notarié (公证立遗嘱), l'achat d'assurances (购买保险) et la création de fiducies (设立信托). Cela révèle que l'esprit de légalité (法治观念) pénètre progressivement la société chinoise, même si changer des coutumes profondément enracinées (根深蒂固的习俗) demande encore du temps.
\end{enumerate}
\textbf{Barème indicatif (10 pts) :} 2 pts par question (1 pt pour l'information exacte, 1 pt pour la formulation claire et la citation du terme chinois quand il est demandé).
\textbf{Points de vigilance :} ne pas traduire \zh{子女} par « fils » uniquement : c'est « enfants » (des deux sexes) ; \zh{传宗接代} (chuán zōng jiē dài) signifie « perpétuer la lignée familiale » ; répondre en phrases complètes, jamais en recopiant le texte chinois.
\end{corrige}

\begin{corrige}[Exercice 10 — Thème et Version]
\textbf{1. Thème — traduction modèle (environ 90 caractères) :}

\zh{在喀麦隆，习俗常常选择主要继承人。在中国，法律保护女儿的权利。我觉得平等比传统更好。写遗嘱是避免冲突的好办法。}

\emph{Zài Kāmàilóng, xísú chángcháng xuǎnzé zhǔyào jìchéngrén. Zài Zhōngguó, fǎlǜ bǎohù nǚ'ér de quánlì. Wǒ juéde píngděng bǐ chuántǒng gèng hǎo. Xiě yízhǔ shì bìmiǎn chōngtū de hǎo bànfǎ.}

\textbf{Justifications :} « l'héritier principal » = \zh{主要继承人} ; « protéger » = \zh{保护} ; comparaison \cle{A 比 B 更 + adjectif} (\zh{平等比传统更好}) ; « méthode » = \zh{办法} ; « éviter les conflits » = \zh{避免冲突}.

\textbf{2. Version — traduction modèle :}

« Les jeunes d'aujourd'hui respectent à la fois la tradition et croient en la loi. Ils pensent que, qu'il s'agisse d'un fils ou d'une fille, tous doivent prendre soin de leurs parents et tous ont le droit d'hériter du patrimoine. Si les membres de la famille peuvent s'asseoir et discuter calmement, ils ne se disputeront plus. »

\textbf{Justifications :} \zh{既...又...} = « à la fois... et... » ; \zh{无论是...还是...都...} = « que ce soit... ou..., tous... » ; \zh{坐下来好好谈一谈} = « s'asseoir pour bien discuter » (\zh{好好} = « bien, calmement ») ; \zh{就...了} marque la conséquence.

\textbf{Barème indicatif (10 pts) :} exactitude du vocabulaire du thème (3 pts), correction grammaticale : ordre des mots, \zh{的/得}, comparatifs (4 pts), respect de la longueur et naturel (2 pts), caractères corrects (1 pt).
\textbf{Points de vigilance :} pas de « que » après \zh{觉得} ; \zh{比} exige souvent \zh{更} ou un complément de degré ; en version, restitue les corrélations (\zh{既...又}, \zh{无论...都}) par des tournures françaises naturelles.
\end{corrige}

\begin{corrige}[Exercice 11 — Expression écrite : mini-exposé comparatif]
\textbf{Proposition de rédaction modèle (environ 140 caractères) :}

\zh{女儿应该不应该跟儿子一样继承遗产？在喀麦隆，很多地方按照传统习俗办事，长子往往是主要继承人，女儿得到的遗产不如儿子多。在中国，法律规定男女平等，女儿跟儿子有一样的权利。现在，越来越多的人写遗嘱，让遗产分配更公平。我认为，女儿也应该继承遗产，因为平等比旧习俗更公平。}

\emph{Nǚ'ér yīnggāi bù yīnggāi gēn érzi yíyàng jìchéng yíchǎn? Zài Kāmàilóng, hěn duō dìfang ànzhào chuántǒng xísú bànshì, zhǎngzǐ wǎngwǎng shì zhǔyào jìchéngrén, nǚ'ér dédào de yíchǎn bùrú érzi duō. Zài Zhōngguó, fǎlǜ guīdìng nánnǚ píngděng, nǚ'ér gēn érzi yǒu yíyàng de quánlì. Xiànzài, yuè lái yuè duō de rén xiě yízhǔ, ràng yíchǎn fēnpèi gèng gōngpíng. Wǒ rènwéi, nǚ'ér yě yīnggāi jìchéng yíchǎn, yīnwèi píngděng bǐ jiù xísú gèng gōngpíng.}

« Les filles doivent-elles hériter comme les fils ? Au Cameroun, dans beaucoup de régions, on suit la coutume traditionnelle : le fils aîné est souvent l'héritier principal et les filles reçoivent moins que les fils. En Chine, la loi stipule l'égalité homme-femme : les filles ont les mêmes droits que les fils. Aujourd'hui, de plus en plus de gens rédigent un testament pour rendre la répartition plus équitable. Je pense que les filles doivent aussi hériter, car l'égalité est plus juste que les vieilles coutumes. »

\textbf{Vérification des consignes :}
\begin{itemize}
    \item Vocabulaire de la leçon (au moins 5) : \zh{遗产, 遗嘱, 习俗, 男女平等, 分配, 权利, 传统, 公平} — 8 mots employés.
    \item Structures comparatives (au moins 3) : \zh{跟...一样} (\zh{跟儿子一样继承}), \zh{不如} (\zh{不如儿子多}), \zh{比} (\zh{平等比旧习俗更公平}), \zh{认为} (\zh{我认为...}) — 4 structures.
    \item Plan respecté : question d'introduction, comparaison Cameroun/Chine, avis personnel en conclusion.
\end{itemize}
\textbf{Barème indicatif (20 pts) :} contenu et comparaison factuelle Cameroun/Chine (5 pts) ; vocabulaire du thème, au moins 5 mots (4 pts) ; structures comparatives, au moins 3 (4 pts) ; grammaire générale : ordre des mots, \zh{的/得/地}, particules (4 pts) ; caractères et longueur 120-150 (3 pts).
\textbf{Points de vigilance :}
\begin{itemize}
    \item Ne pas écrire \zh{一样继承} sans \zh{跟} : il faut \zh{跟儿子一样继承遗产}.
    \item \zh{不如} s'utilise sans \zh{比} : \zh{女儿不如儿子多}, jamais \zh{女儿比不如儿子}.
    \item Après \zh{认为/觉得}, la proposition suit directement, sans mot de liaison.
    \item Le complément de lieu avec \zh{在} se place avant le verbe : \zh{在中国，法律规定...}.
    \item Soigner les tons dans la copie pinyin éventuelle : nánnǚ píngděng, yízhǔ, gōngpíng.
\end{itemize}
\end{corrige}

\newpage

\coursec{Fiche bilan}

\coursec{Leçon 8 — Éléments socioculturels : processus d'hérédité au Cameroun et en Chine (suite : comparaison et exposé)}

\courssub{Carte mentale de la leçon}
\begin{popfigure}
\begin{center}\tcbox[colback=popgreen,colframe=popgreen,coltext=white,arc=9pt,boxsep=4pt,left=6pt,right=6pt]{\bfseries\small Hérédité Cameroun vs Chine}\end{center}
\vspace{-2pt}
\begin{tcbraster}[raster columns=3,raster equal height=rows,raster column skip=6pt,raster row skip=6pt,enhanced,blankest]
\begin{tcolorbox}[enhanced,colback=white,colframe=popgreen,boxrule=0.9pt,arc=3pt,title={Cameroun Coutumes \& Lignée},fonttitle=\bfseries\scriptsize,coltitle=white,colbacktitle=popgreen,left=4pt,right=4pt,top=2pt,bottom=2pt,boxsep=1pt,toptitle=1.5pt,bottomtitle=1.5pt,halign=left]
\begin{itemize}[nosep,leftmargin=*,itemsep=1pt,font=\scriptsize]
\item Patrilinéaire / Matrilinéaire
\item Terre, titres, chefferie
\item Rôle du chef de famille
\end{itemize}
\end{tcolorbox}
\begin{tcolorbox}[enhanced,colback=white,colframe=poporange,boxrule=0.9pt,arc=3pt,title={Chine Droit \& Filialité},fonttitle=\bfseries\scriptsize,coltitle=white,colbacktitle=poporange,left=4pt,right=4pt,top=2pt,bottom=2pt,boxsep=1pt,toptitle=1.5pt,bottomtitle=1.5pt,halign=left]
\begin{itemize}[nosep,leftmargin=*,itemsep=1pt,font=\scriptsize]
\item Code civil : égalité H/F
\item Testament (\zh{遗嘱}) \& Notaire
\item Registre familial (\zh{户口})
\end{itemize}
\end{tcolorbox}
\begin{tcolorbox}[enhanced,colback=white,colframe=poppurple,boxrule=0.9pt,arc=3pt,title={Vocabulaire clé},fonttitle=\bfseries\scriptsize,coltitle=white,colbacktitle=poppurple,left=4pt,right=4pt,top=2pt,bottom=2pt,boxsep=1pt,toptitle=1.5pt,bottomtitle=1.5pt,halign=left]
\begin{itemize}[nosep,leftmargin=*,itemsep=1pt,font=\scriptsize]
\item \zh{遗产} 、 \zh{继承} 、 \zh{遗嘱}
\item \zh{习俗} 、 \zh{法定} 、 \zh{分配}
\item \zh{宗族} 、 \zh{族长} 、 \zh{公证}
\end{itemize}
\end{tcolorbox}
\begin{tcolorbox}[enhanced,colback=white,colframe=poppink,boxrule=0.9pt,arc=3pt,title={Structures comparatives},fonttitle=\bfseries\scriptsize,coltitle=white,colbacktitle=poppink,left=4pt,right=4pt,top=2pt,bottom=2pt,boxsep=1pt,toptitle=1.5pt,bottomtitle=1.5pt,halign=left]
\begin{itemize}[nosep,leftmargin=*,itemsep=1pt,font=\scriptsize]
\item \zh{A 比 B + Adj}
\item \zh{跟 / 和 ... 一样 / 不同}
\item \zh{与其 ... 不如 ...}
\end{itemize}
\end{tcolorbox}
\begin{tcolorbox}[enhanced,colback=white,colframe=popgold,boxrule=0.9pt,arc=3pt,title={Expression orale},fonttitle=\bfseries\scriptsize,coltitle=white,colbacktitle=popgold,left=4pt,right=4pt,top=2pt,bottom=2pt,boxsep=1pt,toptitle=1.5pt,bottomtitle=1.5pt,halign=left]
\begin{itemize}[nosep,leftmargin=*,itemsep=1pt,font=\scriptsize]
\item Mini-exposé 2-3 min
\item Débat : coutume vs loi
\item Connecteurs : \zh{然而} 、 \zh{此外}
\end{itemize}
\end{tcolorbox}
\begin{tcolorbox}[enhanced,colback=white,colframe=popdark,boxrule=0.9pt,arc=3pt,title={Expression écrite},fonttitle=\bfseries\scriptsize,coltitle=white,colbacktitle=popdark,left=4pt,right=4pt,top=2pt,bottom=2pt,boxsep=1pt,toptitle=1.5pt,bottomtitle=1.5pt,halign=left]
\begin{itemize}[nosep,leftmargin=*,itemsep=1pt,font=\scriptsize]
\item Paragraphe comparatif
\item Traduction thème/version
\item Caractères : \zh{继} 、 \zh{承} 、 \zh{遗}
\end{itemize}
\end{tcolorbox}
\begin{tcolorbox}[enhanced,colback=white,colframe=popmuted,boxrule=0.9pt,arc=3pt,title={Méthode},fonttitle=\bfseries\scriptsize,coltitle=white,colbacktitle=popmuted,left=4pt,right=4pt,top=2pt,bottom=2pt,boxsep=1pt,toptitle=1.5pt,bottomtitle=1.5pt,halign=left]
\begin{itemize}[nosep,leftmargin=*,itemsep=1pt,font=\scriptsize]
\item 1. Collecter faits
\item 2. Classer : similitudes / différences
\item 3. Structurer : intro / corps / conclusion
\end{itemize}
\end{tcolorbox}
\end{tcbraster}

\legende{Synthèse visuelle : comparaison des systèmes d'hérédité, lexique, grammaire et compétences visées.}
\end{popfigure}

\courssub{Bilan récapitulatif}

\begin{aretenir}[Lexique clé (HSK 3--4) — Thème : Hérédité et succession]
\begin{center}
\begin{tabularx}{\linewidth}{|X|X|X|X|}
\hline
\rowcolor{popblueL} \textbf{Caractères} & \textbf{Pinyin} & \textbf{Nature} & \textbf{Traduction} \\ \hline
遗产 & yíchǎn & n. & héritage, succession (biens) \\ \hline
继承 & jìchéng & v. & hériter, succéder à \\ \hline
遗嘱 & yízhǔ & n. & testament \\ \hline
法定继承 & fǎdìng jìchéng & n. & succession légale (ab intestat) \\ \hline
遗嘱继承 & yízhǔ jìchéng & n. & succession testamentaire \\ \hline
习俗 & xísú & n. & coutume, tradition \\ \hline
宗族 & zōngzú & n. & clan, lignage patrilinéaire \\ \hline
族长 & zúzhǎng & n. & chef de clan / de famille élargie \\ \hline
分配 & fēnpèi & v. & répartir, partager (héritage) \\ \hline
公证 & gōngzhèng & v./n. & notarier, acte notarié \\ \hline
公证处 & gōngzhèngchù & n. & office notarial \\ \hline
户口 & hùkǒu & n. & registre de population (hukou) \\ \hline
孝道 & xiàodào & n. & piété filiale \\ \hline
土地使用权 & tǔdì shǐyòngquán & n. & droit d'usage du sol (Chine) \\ \hline
\end{tabularx}
\end{center}
\vspace{0.5em}
\textbf{Point de vigilance :} \zh{遗产} (héritage matériel) $\neq$ \zh{遗传} (yíchuán, hérédité biologique / génétique). \zh{分配} se prononce \emph{fēnpèi} (tons 1 et 4), pas ton neutre.
\end{aretenir}

\begin{propriete}[Règles de grammaire — Comparer et nuancer (niveau HSK 3-4)]
\begin{center}
\begin{tabularx}{\linewidth}{|X|X|X|}
\hline
\rowcolor{popblueL} \textbf{Structure} & \textbf{Exemple (Caractères + Pinyin)} & \textbf{Traduction} \\ \hline
\zh{A 比 B + Adj} (supériorité) & \zh{中国的法律 比 习俗 完善。}\newline Zhōngguó de fǎlǜ \textbf{bǐ} xísú \textbf{wánshàn}. & La loi chinoise est plus complète que la coutume. \\ \hline
\zh{A 没有 B + Adj} (infériorité / non-supériorité) & \zh{喀麦隆的习俗 没有 法律 完善。}\newline Kāmàilóng de xísú \textbf{méiyǒu} fǎlǜ \textbf{wánshàn}. & La coutume camerounaise est moins complète que la loi. \\ \hline
\zh{A 跟 / 和 B 一样 / 不一样} & \zh{喀麦隆 跟 中国 不一样。}\newline Kāmàilóng \textbf{gēn} Zhōngguó \textbf{bù yíyàng}. & Le Cameroun n'est pas pareil que la Chine. \\ \hline
\zh{A 与 B 不同} (écrit / formel) & \zh{法定继承 与 习俗继承 不同。}\newline Fǎdìng jìchéng \textbf{yǔ} xísú jìchéng \textbf{bùtóng}. & La succession légale diffère de la coutumière. \\ \hline
\zh{与其 A 不如 B} (préférence / conseil) & \zh{与其 争吵 不如 协商。}\newline \textbf{Yǔqí} zhēngchǎo \textbf{bùrú} xiéshāng. & Plutôt que de se disputer, mieux vaut négocier. \\ \hline
\zh{相比之下} (connecteur de comparaison) & \zh{中国有公证处。相比之下，喀麦隆靠族长。}\newline Zhōngguó yǒu gōngzhèngchù. \textbf{Xiāngbǐ zhīxià}, Kāmàilóng kào zúzhǎng. & La Chine a des notaires. En comparaison, le Cameroun repose sur le chef de clan. \\ \hline
\zh{不仅...而且...} (progression) & \zh{中国不仅有法律，而且有公证。}\newline Zhōngguó \textbf{bùjǐn} yǒu fǎlǜ, \textbf{érqiě} yǒu gōngzhèng. & La Chine a non seulement la loi, mais encore le notariat. \\ \hline
\end{tabularx}
\end{center}
\vspace{0.5em}
\textbf{Rappel :} La négation de \zh{比} est \zh{没有} (pas \zh{不比} sauf style très formel/littéraire). L'adjectif ne prend pas de \zh{很} après \zh{比}.
\end{propriete}

\begin{methode}[Méthodes express pour l'évaluation]
\textbf{1. Mini-exposé oral (2--3 minutes) — Structure type :}
\begin{enumerate}
\item \textbf{Accroche / Annonce :} \zh{今天我谈谈喀麦隆和中国的继承方式。} (Aujourd'hui je parle des modes de succession au Cameroun et en Chine.)
\item \textbf{Présentation Cameroun (faits + vocab) :} \zh{喀麦隆有多种习俗。比如，巴米莱克人是父系继承，土地归族长管理。} (Le Cameroun a diverses coutumes. Ex: les Bamiléké sont patrilinéaires, la terre gérée par le chef de clan.)
\item \textbf{Présentation Chine (faits + vocab) :} \zh{中国实行法定继承。男女平等，有遗嘱按遗嘱办，没遗嘱按法律分。公证处负责公证。} (Chine : succession légale. Égalité H/F. Si testament : selon testament. Sinon : selon loi. Office notarial pour notarié.)
\item \textbf{Comparaison (structures) :} \zh{相比之下，中国法律比较完善，喀麦隆更依赖习俗。中国与喀麦隆不同，中国强调公证。} (En comparaison, la loi chinoise plus complète, Cameroun plus coutume. Chine diffère du Cameroun, Chine insiste sur notariat.)
\item \textbf{Opinion personnelle :} \zh{我认为法定继承更公平，因为男女平等。但是习俗也有文化价值。} (Je pense succession légale plus juste car égalité H/F. Mais coutume a valeur culturelle.)
\item \textbf{Conclusion :} \zh{总之，了解两国差异有助于跨文化交流。} (En somme, comprendre différences aide communication interculturelle.)
\end{enumerate}

\textbf{2. Rédaction comparative (80--100 caractères chinois) — Plan type :}
\begin{enumerate}
\item \textbf{Phrase d'intro :} \zh{中喀两国在继承方面有异同。} (Chine-Cameroun deux pays dans succession ont similitudes/différences.)
\item \textbf{Paragraphe 1 (Cameroun) :} \zh{喀麦隆多按习俗办事，族长分配遗产，土地往往不可买卖。} (Cameroun souvent selon coutume, chef clan partage héritage, terre souvent inaliénable.)
\item \textbf{Paragraphe 2 (Chine) :} \zh{中国依法办事，男女平等继承遗产，遗嘱和公证很常见。} (Chine selon loi, H/F égaux hériter héritage, testament et notariat courants.)
\item \textbf{Paragraphe 3 (Comparatif + avis) :} \zh{相比之下，法律更规范，习俗更传统。我建议结合两者优点。} (En comparaison, loi plus normative, coutume plus traditionnelle. Je suggère combiner avantages des deux.)
\end{enumerate}

\textbf{3. Traduction Thème / Version — Points de vigilance :}
\begin{itemize}
\item \textbf{Termes juridiques :} \zh{法定继承人} (héritier légal) ; \zh{遗嘱继承} (succession testamentaire) ; \zh{放弃继承} (renoncer à la succession) ; \zh{遗产管理人} (administrateur de succession).
\item \textbf{Ordre des mots :} Compléments de lieu (\zh{在喀麦隆}, \zh{在中国}) et de manière (\zh{按习俗}, \zh{依法}) se placent \textbf{AVANT} le verbe.
\item \textbf{Classificateurs :} \zh{一份遗嘱} (un testament), \zh{一位族长} (un chef de clan), \zh{一套房子} (un appartement/maison).
\end{itemize}
\end{methode}

\courssub{Checklist avant le Bac}
\begin{aretenir}[Checklist avant le Bac -- Leçon ZH08]
\begin{itemize}
\item $\square$ Je connais par cœur le vocabulaire de l'héritage (\zh{遗产、继承、遗嘱、法定、习俗、族长、公证、分配、公证处、户口、孝道、土地使用权}).
\item $\square$ J'écris correctement les caractères clés : \zh{继} (ordre : \zh{纟} + \zh{米} + \zh{页}), \zh{承} (\zh{受} + \zh{手}), \zh{遗} (\zh{辶} + \zh{贵}), \zh{嘱} (\zh{口} + \zh{烛}).
\item $\square$ Je maîtrise la structure comparative \zh{A 比 B Adj} et sa négation \zh{A 没有 B Adj} (pas \zh{不比} à l'oral courant).
\item $\square$ J'utilise les connecteurs de comparaison : \zh{相比之下、与...不同、跟...一样、不仅...而且...}.
\item $\square$ Je sais expliquer \textbf{3 faits camerounais} : lignée (patri/matri selon ethnies), rôle du chef/chef de clan (\zh{族长}) dans le partage, terre coutumière inaliénable (\zh{土地不可买卖}).
\item $\square$ Je sais expliquer \textbf{3 faits chinois} : Code civil (\zh{民法典}) égalité H/F (\zh{男女平等}), testament notarié (\zh{公证遗嘱}), \zh{户口} (hukou) et droit d'usage du sol (\zh{土地使用权}, 70 ans pour résidentiel).
\item $\square$ Je produis un mini-exposé oral structuré (Intro / Cameroun / Chine / Comparatif / Avis / Conclusion) en 2--3 minutes, avec connecteurs logiques.
\item $\square$ Je rédige un paragraphe comparatif cohérent (80--100 caractères) avec pinyin et tons corrects.
\item $\square$ Je distingue \zh{遗产} (héritage matériel) de \zh{遗传} (héritage génétique / hérédité biologique).
\item $\square$ Je gère les tons du 3e ton (sandhi : \zh{很好} $\rightarrow$ prononcé \zh{hén hǎo}, \zh{继承} \zh{jìchéng} $\rightarrow$ \zh{jíchéng} devant 4e ton) et vérifie les tons du vocabulaire thème (\zh{分配} \emph{fēnpèi}).
\item $\square$ J'emploie le registre formel (\zh{您、与其...不如、鉴于、依据}) pour le débat ou la rédaction officielle.
\item $\square$ Je cite des exemples concrets (terre à Buea / région de l'Ouest, appartement à Pékin / Shanghai, chefferie Bamiléké / Sawa, notaire à Shanghai) pour illustrer mes propos.
\end{itemize}
\end{aretenir}

\findecours

\end{document}
